% Thalès dans le triangle — Mathématiques 3ème — Cours signé OnBuch+
% Programme officiel MINESEC. Compiler avec : tectonic N02-thales-triangle.tex
\newcommand{\DOCMATIERE}{Mathématiques}
\newcommand{\DOCNIVEAU}{3ème}
\newcommand{\DOCMODULE}{Mathématiques – classe de 3ème}
\newcommand{\DOCLECON}{N02}
\newcommand{\DOCTITRE}{Thalès dans le triangle}
% OnBuch+ — préambule « cours signés » ENRICHI (compilé avec tectonic / XeLaTeX)
% Système visuel « Pop » d'OnBuch+ : fond crème (jamais blanc), bordures encre
% épaisses, ombres « dures » décalées, titres Archivo Black en capitales.
% Version MATHÉMATIQUES : graphes (pgfplots/tikz), boîtes pédagogiques
% (définition, propriété, méthode, attention, le savais-tu, exercice résolu,
% exercice, TP, bilan), page de garde et signature OnBuch+ ancrée en bas.
%
% Macros attendues dans main.tex AVANT \input{preamble} :
%   \DOCMATIERE  \DOCNIVEAU  \DOCMODULE  \DOCLECON (numéro)  \DOCTITRE
\documentclass[11pt]{article}

\usepackage{fontspec}
\usepackage[french]{babel}
\usepackage[a4paper,margin=2cm,top=3.4cm,bottom=2.8cm,headheight=1.4cm,footskip=1.3cm]{geometry}
\usepackage{amsmath,amssymb,amsfonts}
\usepackage{mathtools} % \xmapsto, \coloneqq... souvent utilisés par les modèles
\usepackage{cancel} % simplifications barrées
% \abs{z} (module, valeur absolue) et \norm{u} (norme), utilisés par les modèles
\providecommand{\abs}{}\let\abs\relax\DeclarePairedDelimiter{\abs}{\lvert}{\rvert}
\providecommand{\norm}{}\let\norm\relax\DeclarePairedDelimiter{\norm}{\lVert}{\rVert}
\usepackage[table]{xcolor} % [table] : fournit aussi \rowcolors (alternance de lignes)
\usepackage{pagecolor}
\usepackage{tikz}
\usetikzlibrary{arrows.meta,positioning,calc,shapes.geometric,shapes.misc,shapes.arrows,decorations.pathmorphing,decorations.pathreplacing,decorations.markings,patterns,shadows,mindmap,trees,fit,backgrounds,matrix,intersections,angles,quotes}
\usepackage{pgfplots}
\usepackage{pgf-pie} % diagrammes circulaires \pie (statistiques)
\pgfplotsset{compat=1.18}
\usepgfplotslibrary{fillbetween,groupplots} % aires entre courbes (calcul intégral)
\usepackage{enumitem}
\usepackage{fancyhdr}
% [most] charge toutes les bibliothèques tcolorbox (listings, minted, raster,
% documentation...) et gonfle inutilement la mémoire TeX sur un document long
% (~300 boîtes) : on ne charge que ce qui est réellement utilisé.
\usepackage[skins,breakable]{tcolorbox}
\usepackage{graphicx}
\usepackage{colortbl}
\usepackage{booktabs}
\usepackage{tabularx}
\usepackage{diagbox} % case d'en-tête diagonale des tableaux à double entrée
% Colonnes extensibles centrée / gauche / droite (C, L, R), souvent utilisées par les modèles
\newcolumntype{C}{>{\centering\arraybackslash}X}
\newcolumntype{L}{>{\raggedright\arraybackslash}X}
\newcolumntype{R}{>{\raggedleft\arraybackslash}X}
\usepackage{array}
\usepackage{multicol}
\usepackage{siunitx}
% parse-numbers=false : accepte aussi \SI{2,5 \times 10^{-3}}{...}
\sisetup{locale=FR,output-decimal-marker={,},group-minimum-digits=4,parse-numbers=false}
\DeclareSIUnit\ohm{\ensuremath{\symup{\Omega}}} % U+2126 absent de la police
\usepackage{qrcode}
% \degree : commande fréquemment utilisée par les modèles pour un angle ou une
% température en dehors de \SI{}{} (non fournie par les paquets chargés).
\providecommand{\degree}{^{\circ}}
% \qed (fin de démonstration), utilisé par les modèles sans amsthm
\providecommand{\qed}{\hfill$\square$}
% \barre : double barre des valeurs interdites dans un tableau de variations (tkz-tab)
\providecommand{\barre}{\ensuremath{\|}}
% environnement proof (amsthm non chargé) utilisé par les modèles
\ifdefined\proof\else\newenvironment{proof}[1][Démonstration]{\par\noindent\textit{#1.}\ }{\qed\par}\fi
\usepackage{lastpage}
\usepackage{hyperref}
\hypersetup{hidelinks}

% ============================================================
% Palette « Pop » OnBuch+ — mêmes hex que la classe P (plus_ui.dart)
% ============================================================
\definecolor{poporange}{HTML}{F59321}
\definecolor{popdark}{HTML}{E07A0C}
\definecolor{popcream}{HTML}{FFF6E8}
\definecolor{popink}{HTML}{1C1714}
\definecolor{poppurple}{HTML}{7A5AE0}
\definecolor{popgreen}{HTML}{2E9E6A}
\definecolor{poppink}{HTML}{E0457B}
\definecolor{popblue}{HTML}{2F7DE1}
\definecolor{popgold}{HTML}{C98A12}
\definecolor{popmuted}{HTML}{8A7F76}
\definecolor{popline}{HTML}{EADFD2}
\definecolor{popgray}{HTML}{8A7F76}
\colorlet{popgrey}{popgray}
% Alias tolérants (le modèle nomme parfois les couleurs différemment)
\colorlet{popwhite}{white}
\colorlet{popblack}{popink}
\colorlet{popred}{poppink}
\colorlet{popyellow}{popgold}
% Teintes claires pour fonds de tableaux / zones
\colorlet{popblueL}{popblue!10!white}
\colorlet{poporangeL}{poporange!14!white}
\colorlet{popgreenL}{popgreen!12!white}
\colorlet{poppurpleL}{poppurple!12!white}
\colorlet{poppinkL}{poppink!12!white}
\colorlet{popgoldL}{popgold!14!white}

\pagecolor{popcream}

% ============================================================
% Typographie
% ============================================================
\defaultfontfeatures{Ligatures=TeX}
\setmainfont{PlusJakartaSans-Medium.ttf}[Path=../../../fonts/,BoldFont=PlusJakartaSans-Bold.ttf]
% Maths en sans-serif (Fira Math) avec chiffres et lettres droites de Plus
% Jakarta, pour que les formules se fondent dans le texte.
\usepackage[math-style=ISO,bold-style=ISO]{unicode-math}
\setmathfont{FiraMath-Regular.otf}[Path=../../../fonts/]
\setmathfont{PlusJakartaSans-Medium.ttf}[Path=../../../fonts/,range={up/num,up/{Latin,latin},\mathup}]
\usepackage{icomma} % virgule décimale sans espace en mode math
% Caractères mathématiques Unicode absents des polices de texte (Plus Jakarta,
% Archivo Black) : rendus par leur équivalent mathématique, en texte comme en
% formule — sinon ils s'affichent en carré vide (ex. « Arithmétique dans ℤ »).
\usepackage{newunicodechar}
\newunicodechar{ℝ}{\ensuremath{\mathbb{R}}}
\newunicodechar{ℤ}{\ensuremath{\mathbb{Z}}}
\newunicodechar{ℂ}{\ensuremath{\mathbb{C}}}
\newunicodechar{ℕ}{\ensuremath{\mathbb{N}}}
\newunicodechar{ℚ}{\ensuremath{\mathbb{Q}}}
\newunicodechar{𝔻}{\ensuremath{\mathbb{D}}}
\newunicodechar{α}{\ensuremath{\alpha}}
\newunicodechar{β}{\ensuremath{\beta}}
\newunicodechar{γ}{\ensuremath{\gamma}}
\newunicodechar{δ}{\ensuremath{\delta}}
\newunicodechar{ε}{\ensuremath{\varepsilon}}
\newunicodechar{θ}{\ensuremath{\theta}}
\newunicodechar{λ}{\ensuremath{\lambda}}
\newunicodechar{μ}{\ensuremath{\mu}}
\newunicodechar{σ}{\ensuremath{\sigma}}
\newunicodechar{τ}{\ensuremath{\tau}}
\newunicodechar{φ}{\ensuremath{\varphi}}
\newunicodechar{ψ}{\ensuremath{\psi}}
\newunicodechar{ω}{\ensuremath{\omega}}
\newunicodechar{Σ}{\ensuremath{\Sigma}}
% majuscules produites par \MakeUppercase dans les titres (α-aminés -> Α)
\newunicodechar{Α}{\ensuremath{\alpha}}
\newunicodechar{Β}{\ensuremath{\beta}}
\newunicodechar{Γ}{\ensuremath{\Gamma}}
\newunicodechar{Δ}{\ensuremath{\Delta}}
\newunicodechar{Ε}{\ensuremath{\varepsilon}}
\newunicodechar{Θ}{\ensuremath{\Theta}}
\newunicodechar{Λ}{\ensuremath{\Lambda}}
\newunicodechar{Μ}{\ensuremath{\mu}}
\newunicodechar{Τ}{\ensuremath{\tau}}
\newunicodechar{Φ}{\ensuremath{\Phi}}
\newunicodechar{Ψ}{\ensuremath{\Psi}}
\newunicodechar{Ω}{\ensuremath{\Omega}}
\newunicodechar{↦}{\ensuremath{\mapsto}}
\newunicodechar{∈}{\ensuremath{\in}}
\newunicodechar{∉}{\ensuremath{\notin}}
\newunicodechar{∧}{\ensuremath{\wedge}}
\newunicodechar{⇔}{\ensuremath{\Leftrightarrow}}
\newunicodechar{⟺}{\ensuremath{\Longleftrightarrow}}
\newunicodechar{⇒}{\ensuremath{\Rightarrow}}
\newunicodechar{⟹}{\ensuremath{\Longrightarrow}}
\newunicodechar{∘}{\ensuremath{\circ}}
\newunicodechar{∪}{\ensuremath{\cup}}
\newunicodechar{∩}{\ensuremath{\cap}}
\newunicodechar{≡}{\ensuremath{\equiv}}
\newunicodechar{⊂}{\ensuremath{\subset}}
\newunicodechar{⊥}{\ensuremath{\perp}}
\newunicodechar{∃}{\ensuremath{\exists}}
\newunicodechar{∀}{\ensuremath{\forall}}
\newunicodechar{∛}{\ensuremath{\sqrt[3]{\,}}}
\newunicodechar{∜}{\ensuremath{\sqrt[4]{\,}}}
\newunicodechar{⃗}{\ensuremath{{}^{\to}}}
\newunicodechar{ⁿ}{\ifmmode^{n}\else\textsuperscript{\ensuremath{n}}\fi}
\newunicodechar{ˣ}{\ifmmode^{x}\else\textsuperscript{\ensuremath{x}}\fi}
\newunicodechar{ᵇ}{\ifmmode^{b}\else\textsuperscript{\ensuremath{b}}\fi}
\newunicodechar{ʳ}{\ifmmode^{r}\else\textsuperscript{\ensuremath{r}}\fi}
\newunicodechar{ᵘ}{\ifmmode^{u}\else\textsuperscript{\ensuremath{u}}\fi}
\newunicodechar{ʸ}{\ifmmode^{y}\else\textsuperscript{\ensuremath{y}}\fi}
\newunicodechar{ᵃ}{\ifmmode^{a}\else\textsuperscript{\ensuremath{a}}\fi}
\newunicodechar{ᵉ}{\ifmmode^{e}\else\textsuperscript{\ensuremath{e}}\fi}
\newunicodechar{ᵏ}{\ifmmode^{k}\else\textsuperscript{\ensuremath{k}}\fi}
\newunicodechar{ᵖ}{\ifmmode^{p}\else\textsuperscript{\ensuremath{p}}\fi}
\newunicodechar{ᴺ}{\ifmmode^{N}\else\textsuperscript{\ensuremath{N}}\fi}
\newunicodechar{ᶜ}{\ifmmode^{c}\else\textsuperscript{\ensuremath{c}}\fi}
\newunicodechar{ʰ}{\ifmmode^{h}\else\textsuperscript{\ensuremath{h}}\fi}
\newunicodechar{ᵠ}{\ifmmode^{q}\else\textsuperscript{\ensuremath{q}}\fi}
\newunicodechar{ₙ}{\ifmmode_{n}\else\textsubscript{\ensuremath{n}}\fi}
\newunicodechar{ₐ}{\ifmmode_{a}\else\textsubscript{\ensuremath{a}}\fi}
\newunicodechar{ᵢ}{\ifmmode_{i}\else\textsubscript{\ensuremath{i}}\fi}
\newunicodechar{ᵣ}{\ifmmode_{r}\else\textsubscript{\ensuremath{r}}\fi}
\newunicodechar{ₖ}{\ifmmode_{k}\else\textsubscript{\ensuremath{k}}\fi}
\newunicodechar{ᵦ}{\ifmmode_{\beta}\else\textsubscript{\ensuremath{\beta}}\fi}
\newunicodechar{ₓ}{\ifmmode_{x}\else\textsubscript{\ensuremath{x}}\fi}
\newfontfamily\popdisplay{ArchivoBlack-Regular.ttf}[Path=../../../fonts/]
\newfontfamily\poplogo{SpaceGrotesk-Bold.ttf}[Path=../../../fonts/]
\newfontfamily\popsemi{PlusJakartaSans-SemiBold.ttf}[Path=../../../fonts/]
\newfontfamily\popxbold{PlusJakartaSans-ExtraBold.ttf}[Path=../../../fonts/]
\newfontfamily\popxboldls{PlusJakartaSans-ExtraBold.ttf}[Path=../../../fonts/,LetterSpace=3.0]

\setlist[enumerate]{leftmargin=1.7em,itemsep=5pt,topsep=4pt}
\setlist[itemize]{leftmargin=1.7em,itemsep=4pt,topsep=4pt,label={\color{poporange}\textbullet}}
\setlength{\parindent}{0pt}
\setlength{\parskip}{5pt}
\renewcommand{\arraystretch}{1.25}

% pgfplots : style Pop par défaut
\pgfplotsset{
  every axis/.append style={axis line style={popink,line width=1pt},tick style={popink},
    grid style={popline,line width=0.5pt},label style={font=\small},tick label style={font=\footnotesize}},
  popcurve/.style={poporange,line width=1.8pt,no markers,smooth},
}
% Même style disponible dans un \draw TikZ simple (hors pgfplots)
\tikzset{popcurve/.style={poporange,line width=1.8pt}}

% ============================================================
% Pill / squiggle
% ============================================================
\newcommand{\poppill}[3]{%
  \tikz[baseline=-0.55ex]\node[draw=popink,line width=1.2pt,rounded corners=2.6pt,%
    fill=#2,inner xsep=6.5pt,inner ysep=3pt]{%
    \popxboldls\fontsize{7}{7}\selectfont\color{#3}\MakeUppercase{#1}};%
}
\newcommand{\popsquiggle}[1]{%
  \tikz[baseline=-0.4ex]{
    \draw[#1,line width=2.6pt,line cap=round]
      plot [smooth] coordinates {(0,0.09) (0.34,0.01) (0.68,0.21) (1.02,0.01) (1.36,0.10)};
  }%
}
% Mot-clé surligné (marqueur orange sous le texte)
\newcommand{\cle}[1]{{\popxbold\color{popdark}#1}}

% ============================================================
% En-tête / pied
% ============================================================
\pagestyle{fancy}
\fancyhf{}
\renewcommand{\headrulewidth}{0pt}
\renewcommand{\footrulewidth}{0pt}
\lhead{\raisebox{-0.15ex}{\poplogo\fontsize{13}{13}\selectfont\color{popink}OnBuch\color{poporange}+}}
\rhead{\poppill{\DOCMATIERE\ \textbullet\ \DOCNIVEAU\ \textbullet\ Leçon \DOCLECON}{white}{popink}}
\lfoot{\popxbold\fontsize{7.6}{7.6}\selectfont\color{popmuted}\MakeUppercase{Cours signé OnBuch+}\ \textbullet\ {\popsemi\DOCTITRE}}
\rfoot{\tikz[baseline=-0.6ex]\node[rounded corners=8.5pt,fill=popink,minimum height=17pt,minimum width=17pt,inner xsep=5pt,inner sep=0]{\popxbold\fontsize{8}{8}\selectfont\color{popcream}\thepage\,/\,\pageref*{LastPage}};}

% ============================================================
% Titres
% ============================================================
\newcommand{\chaptitre}[2]{%
  \begin{tcolorbox}[enhanced,colback=poporange,colframe=popink,boxrule=2.6pt,arc=11pt,
      drop shadow=popink,left=15pt,right=15pt,top=12pt,bottom=11pt,before skip=3pt,after skip=18pt]
    \poppill{#2}{popink}{popcream}\par\vspace{9pt}
    {\raggedright\hyphenpenalty=10000\exhyphenpenalty=10000\popdisplay\fontsize{22}{24}\selectfont\color{popcream}\MakeUppercase{#1}\par}
  \end{tcolorbox}
}
% Section numérotée : pastille orange avec le numéro + titre Archivo
\newcounter{popsec}
\newcommand{\coursec}[1]{%
  \stepcounter{popsec}\setcounter{popsub}{0}%
  \par\vspace{16pt}\noindent
  \begin{minipage}{\linewidth}
  \tikz[baseline=-0.7ex]\node[draw=popink,line width=1.6pt,fill=poporange,rounded corners=4pt,minimum size=22pt,inner sep=0]{\popdisplay\fontsize{12}{12}\selectfont\color{popcream}\thepopsec};%
  \hspace{8pt}{\popdisplay\fontsize{15}{17}\selectfont\color{popink}\MakeUppercase{#1}}\par
  \vspace{4pt}\noindent{\color{poporange}\rule{36pt}{3.6pt}}
  \end{minipage}\par\nopagebreak\vspace{8pt}
}
\newcounter{popsub}
\newcommand{\courssub}[1]{%
  \stepcounter{popsub}%
  \par\vspace{9pt}\noindent{\popxbold\fontsize{11.5}{13}\selectfont\color{popink}{\color{poporange}\thepopsec.\thepopsub}\hspace{6pt}#1}\par\nopagebreak\vspace{4pt}
}

% ============================================================
% Boîtes pédagogiques — même recette : fond blanc, bordure épaisse colorée,
% ombre dure encre, badge plein. Titre optionnel [#1].
% ============================================================
\tcbset{popbase/.style={breakable,enhanced,colback=white,boxrule=2.3pt,arc=9pt,
  shadow={3pt}{-3pt}{0pt}{popink},left=14pt,right=14pt,top=11pt,bottom=11pt,before skip=11pt,after skip=15pt,
  fontupper=\popsemi\fontsize{10.6}{14.5}\selectfont\color{popink},
  fontlower=\popsemi\fontsize{10.6}{14.5}\selectfont\color{popink}}}
% Coche dessinée (le glyphe ✓ n'existe pas dans les polices du document)
\newcommand{\popcheck}[1]{\tikz[baseline=-0.5ex]\draw[#1,line width=1.6pt,line cap=round,line join=round](-0.13,0.0)--(-0.04,-0.09)--(0.13,0.1);}
\newcommand{\popboxhead}[3]{\poppill{#1}{#2}{white}\ifx\relax#3\relax\else\hspace{6pt}{\popxbold\fontsize{10.6}{12}\selectfont\color{popink}#3}\fi\par\vspace{7pt}}

\newtcolorbox{aretenir}[1][]{popbase,colframe=popgreen,before upper={\popboxhead{À retenir}{popgreen}{#1}}}
\newtcolorbox{exemplebox}[1][]{popbase,colframe=poporange,before upper={\popboxhead{Exemple}{poporange}{#1}}}
\newtcolorbox{definition}[1][]{popbase,colframe=popblue,before upper={\popboxhead{Définition}{popblue}{#1}}}
\newtcolorbox{propriete}[1][]{popbase,colframe=poppurple,before upper={\popboxhead{Propriété}{poppurple}{#1}}}
\newtcolorbox{methode}[1][]{popbase,colframe=popgold,before upper={\popboxhead{Méthode}{popgold}{#1}}}
\newtcolorbox{attention}[1][]{popbase,colframe=poppink,before upper={\popboxhead{Attention}{poppink}{#1}}}
\newtcolorbox{experience}[1][]{popbase,colframe=popblue,colback=popblueL,before upper={\popboxhead{Expérience}{popblue}{#1}}}
% « Le savais-tu ? » : carte violette pleine (contexte, culture, Cameroun)
\newtcolorbox{savaistu}[1][]{popbase,colback=poppurple,colframe=popink,
  fontupper=\popsemi\fontsize{10.4}{14.2}\selectfont\color{white},
  before upper={\poppill{Le savais-tu\,?}{white}{poppurple}\ifx\relax#1\relax\else\hspace{6pt}{\popxbold\color{white}#1}\fi\par\vspace{7pt}}}
% Exercice résolu : énoncé en haut, \tcblower puis la solution sur fond crème
\newcounter{popexo}\newcounter{popexr}
\newtcolorbox{exoresolu}[1][]{popbase,colframe=poporange,colbacklower=popcream,
  segmentation style={popink,line width=1.2pt,dashed},
  before upper={\stepcounter{popexr}\popboxhead{Exercice résolu \thepopexr}{poporange}{#1}},
  before lower={\poppill{Solution}{popink}{popcream}\par\vspace{6pt}}}
% Exercice (non résolu en place) — la correction est dans la partie Corrigés
\newtcolorbox{exercice}[1][]{popbase,colframe=popink,
  before upper={\stepcounter{popexo}\popboxhead{Exercice \thepopexo}{popink}{#1}}}
% Corrigé
\newtcolorbox{corrige}[1][]{popbase,colframe=popgreen,colback=popgreenL,
  before upper={\popboxhead{Corrigé}{popgreen}{#1}}}
% Objectifs / prérequis / bilan
\newtcolorbox{objectifs}{popbase,colframe=popink,colback=poporangeL,
  before upper={\popboxhead{Objectifs de la leçon}{popink}{}}}
\newtcolorbox{prerequis}{popbase,colframe=popink,colback=popblueL,
  before upper={\popboxhead{Prérequis}{popblue}{}}}
\newtcolorbox{bilan}{popbase,colframe=popink,colback=white,boxrule=2.8pt,
  before upper={\popboxhead{Fiche bilan}{poporange}{L'essentiel à connaître pour l'examen}}}
% Figure encadrée avec légende
\newtcolorbox{popfigure}[1][]{enhanced,breakable=false,colback=white,colframe=popline,boxrule=1.4pt,arc=8pt,
  left=10pt,right=10pt,top=10pt,bottom=8pt,before skip=10pt,after skip=12pt,halign=center,
  lower separated=false,halign lower=center,
  fontlower=\popsemi\fontsize{9}{11.5}\selectfont\color{popmuted},
  #1}
\newcommand{\legende}[1]{\par\vspace{6pt}{\popsemi\fontsize{9}{11.5}\selectfont\color{popmuted}#1}}

% ============================================================
% Signature OnBuch+
% ============================================================
\newcommand{\poplogoinline}[1]{{\poplogo\fontsize{#1}{#1}\selectfont\color{popink}OnBuch\color{poporange}+}}
\newcommand{\signatureonbuch}{%
  \begin{tcolorbox}[enhanced,colback=popink,colframe=popink,boxrule=2.2pt,arc=10pt,
      drop shadow=poporange,left=14pt,right=14pt,top=10pt,bottom=10pt,before skip=0pt,after skip=0pt]
    \tikz[baseline=-0.7ex]\node[circle,fill=popgreen,draw=popcream,line width=1.3pt,minimum size=22pt,inner sep=0]{\popcheck{white}};%
    \hspace{9pt}\begin{minipage}[c]{0.62\linewidth}
      {\popxbold\fontsize{12}{13}\selectfont\color{popcream}Signé\ \textbullet\ {\poplogo OnBuch\color{poporange}+}}\par\vspace{2pt}
      {\raggedright\popsemi\fontsize{8}{10.5}\selectfont\color{popline}\MakeUppercase{Équipe pédagogique OnBuch+}\\\MakeUppercase{Conforme au programme officiel MINESEC}\par}
    \end{minipage}\hfill
    {\poplogo\fontsize{22}{22}\selectfont\color{popcream}OnBuch\color{poporange}+}
  \end{tcolorbox}%
}

% ============================================================
% Page de garde
% #1 titre · #2 matière · #3 niveau · #4 module · #5 n° leçon · #6 durée ·
% #7 description / objectifs (texte ou itemize)
% ============================================================
\newcommand{\pagegarde}[7]{%
  \thispagestyle{empty}
  \newgeometry{a4paper,margin=1.7cm,top=1.6cm,bottom=1.4cm}%
  \begin{tikzpicture}[remember picture,overlay]
    \fill[poporange!18] ([xshift=-3.2cm,yshift=-3.6cm]current page.north east) circle (4.6cm);
    \fill[poppurple!14] ([xshift=2.4cm,yshift=7.6cm]current page.south west) circle (3.8cm);
  \end{tikzpicture}%
  {\poplogo\fontsize{30}{30}\selectfont\color{popink}OnBuch\color{poporange}+}\hfill
  \raisebox{0.6ex}{\poppill{Cours signé \textbullet\ Premium}{popink}{popcream}}\par
  \vspace{1pt}\popsquiggle{poppurple}\par
  \vspace{8pt}
  \begin{tcolorbox}[enhanced,colback=poporange,colframe=popink,boxrule=2.8pt,arc=13pt,
      drop shadow=popink,left=18pt,right=18pt,top=12pt,bottom=13pt,after skip=10pt]
    \poppill{#2 \textbullet\ #3}{popink}{popcream}\hspace{5pt}\poppill{#4}{white}{popink}\par\vspace{8pt}
    {\popdisplay\fontsize{14}{14}\selectfont\color{popink}LEÇON #5}\par\vspace{4pt}
    {\raggedright\hyphenpenalty=10000\exhyphenpenalty=10000\popdisplay\fontsize{25}{27}\selectfont\color{popcream}\MakeUppercase{#1}\par}
    \vspace{8pt}
    {\popxbold\fontsize{9.5}{9.5}\selectfont\color{popink}\MakeUppercase{Durée conseillée : #6}}
  \end{tcolorbox}
  \begin{tcolorbox}[enhanced,colback=white,colframe=popink,boxrule=2.2pt,arc=10pt,
      drop shadow=popink,left=16pt,right=16pt,top=10pt,bottom=10pt,after skip=8pt]
    \poppill{Dans cette leçon}{poppurple}{white}\par\vspace{5pt}
    \popsemi\fontsize{9.3}{12.6}\selectfont\color{popink}#7
  \end{tcolorbox}
  \noindent\begin{minipage}[t]{0.487\linewidth}
    \begin{tcolorbox}[enhanced,colback=popgreen,colframe=popink,boxrule=2.2pt,arc=10pt,
        drop shadow=popink,left=11pt,right=11pt,top=10pt,bottom=10pt,height=3.1cm,valign=center]
      \tcbox[colback=white,colframe=popink,boxrule=1.3pt,arc=5pt,boxsep=0pt,nobeforeafter,
        left=3pt,right=3pt,top=3pt,bottom=3pt,valign=center]{\qrcode[height=1.9cm]{https://play.google.com/store/apps/details?id=com.luvvix.onbuch&hl=fr}}%
      \hspace{8pt}\begin{minipage}[c]{0.5\linewidth}
        {\popdisplay\fontsize{11}{12.5}\selectfont\color{white}\MakeUppercase{Télécharge OnBuch}}\par\vspace{4pt}
        {\raggedright\popsemi\fontsize{8.4}{11}\selectfont\color{white}Gratuit \textbullet\ Google Play\\Tuteur IA, annales, concours, résultats\par}
      \end{minipage}
    \end{tcolorbox}
  \end{minipage}\hfill
  \begin{minipage}[t]{0.487\linewidth}
    \begin{tcolorbox}[enhanced,colback=poppurple,colframe=popink,boxrule=2.2pt,arc=10pt,
        drop shadow=popink,left=11pt,right=11pt,top=10pt,bottom=10pt,height=3.1cm,valign=center]
      \tikz[baseline=-0.7ex]\node[circle,fill=white,draw=popink,line width=1.3pt,minimum size=1.6cm,inner sep=0]{\popdisplay\fontsize{24}{24}\selectfont\color{poppurple}+};%
      \hspace{8pt}\begin{minipage}[c]{0.6\linewidth}
        {\popdisplay\fontsize{11}{12.5}\selectfont\color{white}\MakeUppercase{Passe à OnBuch+}}\par\vspace{4pt}
        {\raggedright\popsemi\fontsize{8.4}{11}\selectfont\color{white}Tous les cours signés, Tuteur IA illimité, fiches PDF — onglet OnBuch+ de l'app\par}
      \end{minipage}
    \end{tcolorbox}
  \end{minipage}
  \vspace{8pt}
  \popcontenu
  \vfill
  \signatureonbuch
  \restoregeometry
  \newpage
}

% Rangée « Ce cours contient » (page de garde)
\newcommand{\poptile}[3]{% #1 couleur #2 gros texte #3 légende
  \begin{tcolorbox}[enhanced,colback=#1,colframe=popink,boxrule=2pt,arc=9pt,shadow={2.5pt}{-2.5pt}{0pt}{popink},
      left=6pt,right=6pt,top=9pt,bottom=9pt,halign=center,valign=center,height=1.75cm,width=\linewidth,nobeforeafter]
    {\popdisplay\fontsize{12.5}{13}\selectfont\color{white}\MakeUppercase{#2}}\par\vspace{3pt}
    {\centering\popsemi\fontsize{7.8}{9.5}\selectfont\color{white}#3\par}
  \end{tcolorbox}}
\newcommand{\popcontenu}{%
  \par\noindent{\popxboldls\fontsize{7.5}{7.5}\selectfont\color{popmuted}CE COURS SIGNÉ CONTIENT}\par\vspace{7pt}
  \noindent\begin{minipage}[t]{0.235\linewidth}\poptile{popblue}{Cours}{complet, illustré, pas à pas}\end{minipage}\hfill
  \begin{minipage}[t]{0.235\linewidth}\poptile{poporange}{Méthodes}{et exercices résolus}\end{minipage}\hfill
  \begin{minipage}[t]{0.235\linewidth}\poptile{poppink}{Exercices}{type examen + corrigés}\end{minipage}\hfill
  \begin{minipage}[t]{0.235\linewidth}\poptile{popgreen}{Fiche bilan}{l'essentiel à retenir}\end{minipage}\par}

% Sommaire
\newtcolorbox{sommaire}{enhanced,colback=white,colframe=popink,boxrule=2.2pt,arc=10pt,
  drop shadow=popink,left=18pt,right=18pt,top=13pt,bottom=13pt,before skip=6pt,after skip=20pt,
  before upper={\poppill{Sommaire}{poppurple}{white}\par\vspace{9pt}\popsemi\fontsize{10.6}{14.5}\selectfont\color{popink}}}

% Signature de fin de cours — poussée tout en bas de la dernière page
\newcommand{\findecours}{%
  \par\vfill\nopagebreak
  \begin{center}\popsquiggle{poporange}\end{center}\vspace{4pt}
  \signatureonbuch
}
\AtBeginDocument{\def\llbracket{\mathopen{[\mkern-3.5mu[}}\def\rrbracket{\mathclose{]\mkern-3.5mu]}}} % intervalles d'entiers (glyphe ⟦ absent de la police)
% Glyphes absents de Fira Math : redéfinis à partir de glyphes disponibles
\AtBeginDocument{%
  \def\setminus{\mathbin{\backslash}}\let\smallsetminus\setminus
  \def\vdots{\mathord{\vbox{\baselineskip3.2pt\lineskiplimit0pt\kern4pt\hbox{.}\hbox{.}\hbox{.}}}}%
  \def\ddots{\mathinner{\mkern1mu\raise6.5pt\hbox{.}\mkern2mu\raise3.6pt\hbox{.}\mkern2mu\raise0.7pt\hbox{.}\mkern1mu}}%
  \def\triangle{\mathord{\scalebox{0.9}{$\symup{\Delta}$}}}%
  \def\nabla{\mathord{\raisebox{\depth}{\scalebox{1}[-1]{$\symup{\Delta}$}}}}%
  \def\checkmark{\popcheck{popdark}}%
  \def\star{\mathbin{\ast}}\def\bigstar{\mathord{\ast}}%
  \def\diamond{\mathbin{\scalebox{0.6}{$\Diamond$}}}%
  \def\bigcup{\mathop{\mathchoice{\raisebox{-0.3em}{\scalebox{1.7}{$\cup$}}}{\scalebox{1.25}{$\cup$}}{\cup}{\cup}}\displaylimits}%
  \def\bigcap{\mathop{\mathchoice{\raisebox{-0.3em}{\scalebox{1.7}{$\cap$}}}{\scalebox{1.25}{$\cap$}}{\cap}{\cap}}\displaylimits}%
  \def\lesseqqgtr{\mathrel{\lessgtr}}\def\bigoplus{\mathop{\oplus}\displaylimits}%
}
\providecommand{\ssi}{si et seulement si} % abréviation fréquente
\AtBeginDocument{\providecommand{\wideparen}{\overparen}} % arc de cercle

\begin{document}

\pagegarde{Thalès dans le triangle}{Mathématiques}{3ème}{Mathématiques – classe de 3ème}{N02}{2 séances (fiche de progression harmonisée)}{Cette leçon présente la propriété directe de Thalès et sa réciproque dans le triangle, outils fondamentaux pour calculer des longueurs inaccessibles et justifier que des droites sont parallèles. À travers des situations concrètes de la vie courante camerounaise, l'élève apprend à modéliser, rédiger et justifier rigoureusement ses démarches.
\vspace{4pt}
\begin{multicols}{2}\small
\begin{itemize}
\item À la fin de cette leçon, je sais reconnaître une configuration de Thalès dans une figure (droites parallèles coupant deux sécantes)
\item À la fin de cette leçon, je sais énoncer et appliquer la propriété directe de Thalès pour écrire des rapports de longueurs égaux
\item À la fin de cette leçon, je sais calculer une longueur inconnue en utilisant l'égalité des rapports et la quatrième proportionnelle
\item À la fin de cette leçon, je sais appliquer la propriété de Thalès à des situations concrètes de la vie courante (hauteurs, distances, plans, agrandissements)
\item À la fin de cette leçon, je sais énoncer et appliquer la propriété réciproque de Thalès pour démontrer que deux droites sont parallèles
\item À la fin de cette leçon, je sais utiliser la contraposée de la propriété de Thalès pour démontrer que deux droites ne sont pas parallèles
\end{itemize}\end{multicols}}

\begin{sommaire}
\begin{multicols}{2}
\begin{enumerate}
\item Découverte de la configuration de Thalès
\item Propriété directe de Thalès
\item Applications à la vie courante
\item Propriété réciproque de Thalès
\item Contraposée : prouver que des droites ne sont pas parallèles
\item Synthèse et prolongements
\item Activité d'intégration
\item Méthodes et exercices résolus
\item Exercices
\item Corrigés des exercices
\item Fiche bilan
\end{enumerate}
\end{multicols}
\end{sommaire}

\begin{objectifs}
\begin{itemize}
\item À la fin de cette leçon, je sais reconnaître une configuration de Thalès dans une figure (droites parallèles coupant deux sécantes)
\item À la fin de cette leçon, je sais énoncer et appliquer la propriété directe de Thalès pour écrire des rapports de longueurs égaux
\item À la fin de cette leçon, je sais calculer une longueur inconnue en utilisant l'égalité des rapports et la quatrième proportionnelle
\item À la fin de cette leçon, je sais appliquer la propriété de Thalès à des situations concrètes de la vie courante (hauteurs, distances, plans, agrandissements)
\item À la fin de cette leçon, je sais énoncer et appliquer la propriété réciproque de Thalès pour démontrer que deux droites sont parallèles
\item À la fin de cette leçon, je sais utiliser la contraposée de la propriété de Thalès pour démontrer que deux droites ne sont pas parallèles
\item À la fin de cette leçon, je sais rédiger et justifier une démarche complète : hypothèses, propriété utilisée, calculs, conclusion
\item À la fin de cette leçon, je sais distinguer les situations où utiliser la propriété directe de celles où utiliser la réciproque
\end{itemize}
\end{objectifs}

\begin{prerequis}
\begin{itemize}
\item Égalité de fractions et produits en croix (quatrième proportionnelle)
\item Proportionnalité et tableaux de proportionnalité (6ème-5ème)
\item Droites parallèles et points alignés (notions de 6ème-4ème)
\item Longueurs, mesures et conversions d'unités (m, cm)
\item Notion de rapport de deux longueurs et simplification de fractions
\end{itemize}
\end{prerequis}

\newpage

\coursec{Découverte de la configuration de Thalès}

\courssub{Une question posée dans la cour du collège}

Il est onze heures à Yaoundé, le soleil est haut. Moussa, élève en classe de 3ème, regarde le mât du drapeau de son collège et se demande : \emph{quelle est sa hauteur ?} Impossible d'y grimper, et aucune échelle n'est assez longue. Pourtant, Moussa a une idée lumineuse. Il plante verticalement un bâton de $1,5$~m à quelques mètres du mât, au soleil. Les rayons du soleil, qui arrivent tous parallèles, projettent au sol l'ombre du bâton et l'ombre du mât. En reculant légèrement, Moussa observe que le sommet du bâton, le sommet du mât et l'extrémité des deux ombres sont \textbf{alignés}. Il mesure alors les ombres au sol : l'ombre du bâton mesure $1$~m, celle du mât mesure $4$~m. Grâce à ces seules mesures, Moussa va pouvoir \emph{calculer} la hauteur du mât, sans jamais le toucher !

Comment est-ce possible ? Quel lien mystérieux existe-t-il entre les longueurs des ombres et les hauteurs des objets ? C'est le \cle{théorème de Thalès}, découvert il y a plus de 2500 ans, qui va nous donner la réponse. Toute cette leçon est consacrée à la découverte de la \emph{configuration géométrique} qui se cache derrière cette astuce.

\begin{popfigure}
\begin{tikzpicture}[scale=1.1]
% Sol
\draw[popline, thick] (-0.5,0) -- (8.5,0);
% Ombre du mât (hachurée)
\fill[poppinkL] (0,0) -- (4,0) -- (4,0.02) -- cycle;
\fill[poporangeL] (0,0) -- (1,0) -- (1,0.02) -- cycle;
% Mât
\draw[popblue, very thick] (4,0) -- (4,4);
\fill[popblue] (4,4) circle (0.06);
\node[popblue, right] at (4.05,3.4) {Mât};
\node[popblue, left] at (3.95,2) {$h = ?$};
% Bâton
\draw[poporange, very thick] (1,0) -- (1,1.5);
\fill[poporange] (1,1.5) circle (0.06);
\node[poporange, above left] at (1,1.5) {Bâton};
\node[poporange, right] at (1.05,0.75) {$1,5$ m};
% Rayon du soleil (aligné)
\draw[popink, thick, dashed] (0,0) -- (4,4);
\draw[popink, thick, dashed, ->] (4,4) -- (5.6,5.6) node[above right] {rayon du soleil};
% Extrémité des ombres
\fill[popdark] (0,0) circle (0.06);
\node[popdark, below left] at (0,0) {extrémité des ombres};
% Cotes des ombres
\draw[<->, popmuted] (0,-0.35) -- (1,-0.35) node[midway, below] {$1$ m};
\draw[<->, popmuted] (0,-0.75) -- (4,-0.75) node[midway, below] {$4$ m};
\end{tikzpicture}
\legende{La situation de Moussa : le bâton et le mât sont verticaux (donc parallèles), et le rayon du soleil aligne leurs sommets avec l'extrémité des ombres.}
\end{popfigure}

\courssub{De la réalité au dessin géométrique : la modélisation}

Pour résoudre le problème de Moussa, la première étape du mathématicien consiste à \emph{modéliser} la situation : on remplace les objets réels par des figures géométriques.

\begin{itemize}
\item Le mât et le bâton sont \textbf{verticaux} : ils sont donc représentés par deux segments \textbf{parallèles}.
\item Le rayon du soleil qui passe par les deux sommets est une \textbf{droite}.
\item Le sol, où l'on mesure les ombres, est une autre \textbf{droite}.
\item Ces deux droites (rayon du soleil et sol) se coupent en un point : l'extrémité des ombres.
\end{itemize}

On obtient ainsi un \textbf{triangle} : le grand triangle formé par le sol, le rayon du soleil et le mât. À l'intérieur de ce triangle, le bâton dessine un \textbf{segment parallèle au mât}, qui découpe un \emph{petit triangle} semblable au grand. C'est exactement la figure que nous allons étudier.

\begin{definition}[Configuration de Thalès]
On dit que des points forment une \cle{configuration de Thalès} lorsqu'on a un triangle $AMN$ tel que :
\begin{itemize}
\item le point $M$ appartient à la droite $(AB)$ ;
\item le point $N$ appartient à la droite $(AC)$ ;
\item les droites $(MN)$ et $(BC)$ sont \textbf{parallèles}.
\end{itemize}
Autrement dit : deux droites parallèles $(MN)$ et $(BC)$ sont coupées par deux droites \textbf{sécantes} $(AB)$ et $(AC)$ qui se rencontrent au point $A$.
\end{definition}

\begin{attention}[Bien lire les hypothèses]
La configuration de Thalès exige \textbf{deux} conditions indispensables :
\begin{enumerate}
\item les points $A$, $M$, $B$ sont alignés \emph{et} les points $A$, $N$, $C$ sont alignés ;
\item les droites $(MN)$ et $(BC)$ sont parallèles.
\end{enumerate}
Si l'une de ces conditions manque (par exemple si $(MN)$ n'est pas parallèle à $(BC)$), on ne peut \textbf{pas} appliquer le théorème de Thalès. Dans un exercice, il faut toujours vérifier et \emph{rédiger} ces deux conditions avant tout calcul.
\end{attention}

\courssub{Les deux formes de la configuration}

\courssub{Première forme : le triangle « emboîté »}

Dans la situation de Moussa, les points $M$ et $N$ se trouvent \emph{entre} $A$ et les sommets $B$, $C$ : le petit triangle $AMN$ est \textbf{à l'intérieur} du grand triangle $ABC$, comme deux poupées emboîtées. On parle de configuration \emph{emboîtée} (ou « triangle dans le triangle »).

\begin{popfigure}
\begin{tikzpicture}[scale=1.6]
% Sommets
\coordinate (A) at (0,3);
\coordinate (B) at (-1.8,0);
\coordinate (C) at (2.2,0);
% Points M et N (rapport 0.55)
\coordinate (M) at ($(A)!0.55!(B)$);
\coordinate (N) at ($(A)!0.55!(C)$);
% Triangles
\fill[popblueL, opacity=0.6] (A) -- (M) -- (N) -- cycle;
\draw[popdark, thick] (A) -- (B) -- (C) -- cycle;
\draw[popgreen, very thick] (M) -- (N);
% Points et labels
\fill[popink] (A) circle (0.045); \node[popink, above] at (A) {$A$};
\fill[popdark] (B) circle (0.045); \node[popdark, below left] at (B) {$B$};
\fill[popdark] (C) circle (0.045); \node[popdark, below right] at (C) {$C$};
\fill[popgreen] (M) circle (0.045); \node[popgreen, left] at (M) {$M$};
\fill[popgreen] (N) circle (0.045); \node[popgreen, right] at (N) {$N$};
% Cotes
\node[popblue, left, font=\small] at ($(A)!0.28!(M)$) {$AM$};
\node[popmuted, left, font=\small] at ($(M)!0.5!(B)$) {$MB$};
\node[popblue, right, font=\small] at ($(A)!0.28!(N)$) {$AN$};
\node[popmuted, right, font=\small] at ($(N)!0.5!(C)$) {$NC$};
\node[popgreen, above, font=\small] at ($(M)!0.5!(N)$) {$MN$};
\node[popdark, below, font=\small] at ($(B)!0.5!(C)$) {$BC$};
% Codage parallélisme
\draw[popgreen, ->, thick] ($(M)!0.4!(N)$) -- ($(M)!0.6!(N)$);
\draw[popgreen, ->, thick] ($(B)!0.4!(C)$) -- ($(B)!0.6!(C)$);
\end{tikzpicture}
\legende{Configuration « emboîtée » : $M \in [AB]$, $N \in [AC]$ et $(MN) \parallel (BC)$. Le petit triangle $AMN$ est à l'intérieur du triangle $ABC$.}
\end{popfigure}

\courssub{Deuxième forme : le triangle « papillon »}

Il existe une seconde façon de rencontrer cette figure. Imaginons que les deux droites sécantes $(AB)$ et $(AC)$ se croisent en $A$, et que les points $M$ et $N$ se trouvent \emph{de l'autre côté de $A$}, sur les prolongements des demi-droites. Les deux parallèles $(MN)$ et $(BC)$ sont alors de part et d'autre du point $A$ : la figure ressemble aux ailes d'un papillon. On parle de configuration \emph{papillon} (ou « en croix »).

\begin{popfigure}
\begin{tikzpicture}[scale=1.4]
% Sommets : A au centre, M et B de part et d'autre
\coordinate (A) at (0,0);
\coordinate (B) at (2.2,1.4);
\coordinate (C) at (-2.2,1.4);
\coordinate (M) at ($(A)!0.6!(B)$);
\coordinate (N) at ($(A)!0.6!(C)$);
% Prolongements au-delà de A
\coordinate (M') at ($(A)!-0.6!(B)$);
\coordinate (N') at ($(A)!-0.6!(C)$);
% On place le petit triangle de l'autre côté : M' et N'
\fill[poporangeL, opacity=0.7] (A) -- (M') -- (N') -- cycle;
\fill[popblueL, opacity=0.5] (A) -- (B) -- (C) -- cycle;
% Sécantes
\draw[popdark, thick] (N') -- (B);
\draw[popdark, thick] (M') -- (C);
% Parallèles
\draw[popblue, very thick] (B) -- (C);
\draw[poporange, very thick] (M') -- (N');
% Codage parallélisme
\draw[popgreen, ->, thick] ($(B)!0.45!(C)$) -- ($(B)!0.55!(C)$);
\draw[popgreen, ->, thick] ($(M')!0.45!(N')$) -- ($(M')!0.55!(N')$);
% Points et labels
\fill[popink] (A) circle (0.05); \node[popink, below right] at (A) {$A$};
\fill[popblue] (B) circle (0.045); \node[popblue, above right] at (B) {$B$};
\fill[popblue] (C) circle (0.045); \node[popblue, above left] at (C) {$C$};
\fill[poporange] (M') circle (0.045); \node[poporange, below left] at (M') {$M$};
\fill[poporange] (N') circle (0.045); \node[poporange, below right] at (N') {$N$};
\end{tikzpicture}
\legende{Configuration « papillon » : le point $A$ est \emph{entre} $M$ et $B$, et \emph{entre} $N$ et $C$. On a toujours $(MN) \parallel (BC)$.}
\end{popfigure}

\begin{aretenir}[Les deux formes à connaître]
Dans les deux cas, la structure est la même :
\begin{itemize}
\item deux droites \textbf{sécantes en $A$} : $(AB)$ et $(AC)$ ;
\item deux droites \textbf{parallèles} : $(MN)$ et $(BC)$, avec $M \in (AB)$ et $N \in (AC)$.
\end{itemize}
Seule la \emph{position} des points change : dans la forme emboîtée, $M$ et $N$ sont sur les segments $[AB]$ et $[AC]$ ; dans la forme papillon, $A$ est entre $M$ et $B$ et entre $N$ et $C$. Le théorème de Thalès s'appliquera dans les deux cas.
\end{aretenir}

\begin{methode}[Reconnaître une configuration de Thalès dans une figure]
Face à une figure géométrique, voici la démarche à suivre :
\begin{enumerate}
\item \textbf{Chercher deux droites parallèles} dans la figure (elles sont souvent codées par des flèches identiques, ou indiquées dans l'énoncé).
\item \textbf{Vérifier que ces deux parallèles sont coupées par deux droites sécantes} qui se rencontrent en un point : ce point sera le sommet commun, qu'on note en général $A$.
\item \textbf{Nommer les points d'intersection} : sur la première sécante, on lit $A$, $M$, $B$ ; sur la deuxième, $A$, $N$, $C$ — en veillant à ce que $M$ et $N$ soient sur la \emph{même} parallèle.
\item \textbf{Conclure} : « les points $A$, $M$, $B$ sont alignés, les points $A$, $N$, $C$ sont alignés, et $(MN) \parallel (BC)$ : on est dans une configuration de Thalès ».
\end{enumerate}
\end{methode}

\courssub{Activité de mesure : vers une conjecture}

\begin{experience}[Mesurer pour conjecturer]
Réalisons ensemble l'expérience suivante (tu peux la refaire sur ton cahier avec une règle graduée).
\begin{enumerate}
\item Trace un triangle $ABC$ avec $AB = 6$~cm, $AC = 8$~cm et $BC = 7$~cm.
\item Place le point $M$ sur le segment $[AB]$ tel que $AM = 3$~cm.
\item Trace la parallèle à $(BC)$ passant par $M$ ; elle coupe $[AC]$ en $N$.
\item Mesure soigneusement : $AN$, $MN$.
\item Calcule les trois rapports $\dfrac{AM}{AB}$, $\dfrac{AN}{AC}$ et $\dfrac{MN}{BC}$.
\end{enumerate}
En mesurant, on trouve $AN \approx 4$~cm et $MN \approx 3,5$~cm. Calculons :
\[ \frac{AM}{AB} = \frac{3}{6} = 0,5 \qquad ; \qquad \frac{AN}{AC} = \frac{4}{8} = 0,5 \qquad ; \qquad \frac{MN}{BC} = \frac{3,5}{7} = 0,5. \]
Les trois rapports sont \textbf{égaux} ! Recommence avec $AM = 2$~cm : tu trouveras $AN \approx 2,67$~cm, $MN \approx 2,33$~cm, et les trois rapports valent tous $\dfrac{1}{3}$.
\end{experience}

Ce n'est pas un hasard : quelle que soit la position de $M$ sur $[AB]$, dès que $(MN)$ est parallèle à $(BC)$, les trois rapports semblent toujours égaux. Nous venons de formuler une \cle{conjecture}, c'est-à-dire un énoncé que l'expérience rend très vraisemblable et qu'il reste à démontrer. Ce sera l'objet de la section suivante : la \emph{propriété directe de Thalès}.

\begin{exemplebox}[Le vocabulaire des segments proportionnels]
Reprenons les mesures de l'activité : $\dfrac{AM}{AB} = \dfrac{AN}{AC} = \dfrac{MN}{BC} = \dfrac{1}{2}$.
On dit que les droites parallèles $(MN)$ et $(BC)$ \textbf{déterminent sur les sécantes $(AB)$ et $(AC)$ des segments proportionnels}. Cela signifie que les longueurs découpées sur la première sécante ($AM$ et $AB$) sont dans le même rapport que les longueurs correspondantes découpées sur la deuxième sécante ($AN$ et $AC$). Ici, le rapport commun est $\dfrac{1}{2}$ : le petit triangle est une « réduction de moitié » du grand.
\end{exemplebox}

\begin{attention}[Erreur fréquente : confondre les sommets]
Dans les rapports de Thalès, toutes les longueurs doivent \textbf{partir du même sommet $A$}, le point de concours des deux sécantes. On écrit :
\[ \frac{AM}{AB} = \frac{AN}{AC} = \frac{MN}{BC}. \]
Il est \textbf{faux} d'écrire sans précaution $\dfrac{AM}{MB} = \dfrac{AN}{NC}$ dans un calcul direct, ou pire, de mélanger les sécantes en écrivant $\dfrac{AM}{AC} = \dfrac{AN}{AB}$. Un bon réflexe : vérifie que chaque rapport compare deux longueurs \emph{portées par la même droite} ($AM$ et $AB$ sont sur $(AB)$ ; $AN$ et $AC$ sont sur $(AC)$ ; $MN$ et $BC$ sont sur les deux parallèles).
\end{attention}

\begin{exoresolu}[Reconnaître une configuration et prévoir un rapport]
Dans un triangle $KLM$, on place le point $P$ sur $[KL]$ avec $KP = 4$~cm et $KL = 10$~cm, puis le point $Q$ sur $[KM]$ tel que $(PQ) \parallel (LM)$.
\begin{enumerate}
\item Justifier que l'on est dans une configuration de Thalès.
\item Sans mesurer, quelle valeur semble prendre le rapport $\dfrac{KQ}{KM}$ ?
\end{enumerate}
\tcblower
\textbf{Solution.}
\begin{enumerate}
\item Les points $K$, $P$, $L$ sont alignés (car $P \in [KL]$), les points $K$, $Q$, $M$ sont alignés (car $Q \in [KM]$), et les droites $(PQ)$ et $(LM)$ sont parallèles. On est donc bien dans une \textbf{configuration de Thalès}, de sommet commun $K$ (forme emboîtée).
\item D'après notre conjecture, les trois rapports sont égaux :
\[ \frac{KQ}{KM} = \frac{KP}{KL} = \frac{4}{10} = \frac{2}{5} = 0,4. \]
Le rapport $\dfrac{KQ}{KM}$ semble donc valoir $0,4$ : le point $Q$ est situé aux deux cinquièmes du segment $[KM]$ en partant de $K$.
\end{enumerate}
\end{exoresolu}

\begin{savaistu}[Thalès et la hauteur des pyramides]
\cle{Thalès de Milet} (vers $-625$ / $-547$) est l'un des premiers mathématiciens de l'histoire. Selon la légende rapportée par les historiens grecs, lors d'un voyage en Égypte, il aurait étonné le pharaon en mesurant la hauteur de la \textbf{grande pyramide de Gizeh}... sans y monter ! Son astuce : planter un bâton vertical dans le sable et attendre le moment précis où l'ombre du bâton est égale à sa hauteur. À ce même instant, l'ombre de la pyramide est égale à sa hauteur : il suffit alors de mesurer l'ombre au sol. C'est exactement le principe utilisé par Moussa pour son mât, et la configuration géométrique mise en jeu est celle que nous venons de découvrir. Plus de 2500 ans plus tard, ce théorème sert encore aux géomètres, aux architectes et aux ingénieurs du Cameroun pour mesurer des distances inaccessibles : largeur d'un fleuve, hauteur d'un pylône MTN, profondeur d'un puits.
\end{savaistu}

\begin{aretenir}[Ce qu'il faut retenir de cette section]
\begin{itemize}
\item Une \textbf{configuration de Thalès} est formée de deux droites parallèles $(MN)$ et $(BC)$ coupées par deux sécantes $(AB)$ et $(AC)$ concourantes en $A$.
\item Elle existe sous \textbf{deux formes} : le triangle emboîté ($M \in [AB]$, $N \in [AC]$) et le triangle papillon ($A$ entre $M$ et $B$, et entre $N$ et $C$).
\item L'activité de mesure conduit à la \textbf{conjecture} : $\dfrac{AM}{AB} = \dfrac{AN}{AC} = \dfrac{MN}{BC}$.
\item On dit que les parallèles déterminent sur les sécantes des \textbf{segments proportionnels}.
\item Tous les rapports partent du \textbf{même sommet} $A$.
\end{itemize}
Dans la prochaine section, nous démontrerons cette conjecture : ce sera la \emph{propriété directe de Thalès}, qui permettra enfin à Moussa de calculer la hauteur du mât de son collège.
\end{aretenir}

\coursec{Propriété directe de Thalès}

Dans la section précédente, nous avons découvert la \cle{configuration de Thalès} : un triangle $ABC$, un point $M$ sur le côté $[AB]$, un point $N$ sur le côté $[AC]$, et la droite $(MN)$ parallèle au troisième côté $(BC)$. Nous avons observé, en mesurant sur plusieurs figures, que les longueurs semblaient toujours \emph{proportionnelles}. Le moment est venu de transformer cette observation en un énoncé précis, que l'on pourra utiliser dans tous les exercices : c'est la \cle{propriété directe de Thalès}, l'un des résultats les plus importants de la géométrie de 3ème et un outil incontournable du BEPC.

\courssub{De l'observation à l'énoncé}

Reprenons l'expérience de la section précédente. Sur une figure où $AM = 2$ cm, $AB = 6$ cm, $AN = 3$ cm, $AC = 9$ cm, $MN = 2{,}5$ cm et $BC = 7{,}5$ cm, on calcule :
\[ \frac{AM}{AB} = \frac{2}{6} = \frac{1}{3} \qquad ; \qquad \frac{AN}{AC} = \frac{3}{9} = \frac{1}{3} \qquad ; \qquad \frac{MN}{BC} = \frac{2{,}5}{7{,}5} = \frac{1}{3}. \]
Les trois rapports sont égaux ! Ce n'est pas un hasard : dès que la droite $(MN)$ est parallèle à $(BC)$, le petit triangle $AMN$ est une \emph{réduction} du grand triangle $ABC$, avec un même \cle{coefficient de proportionnalité} (ici $\dfrac{1}{3}$). Toutes les longueurs du petit triangle s'obtiennent en multipliant celles du grand triangle par ce coefficient. C'est exactement ce que dit la propriété de Thalès.

\begin{propriete}[Propriété directe de Thalès (admise)]
Soit $ABC$ un triangle. Si $M$ est un point de la droite $(AB)$, $N$ un point de la droite $(AC)$, et si les droites $(MN)$ et $(BC)$ sont \cle{parallèles}, alors :
\[ \frac{AM}{AB} = \frac{AN}{AC} = \frac{MN}{BC}. \]
Autrement dit, les longueurs des côtés du triangle $AMN$ sont proportionnelles aux longueurs des côtés correspondants du triangle $ABC$.
\end{propriete}

\begin{attention}[Deux hypothèses indispensables]
Pour pouvoir appliquer la propriété de Thalès, il faut vérifier \textbf{deux} conditions :
\begin{enumerate}
\item les points $A$, $M$, $B$ sont alignés, et les points $A$, $N$, $C$ sont alignés (les points sont sur les côtés du triangle) ;
\item les droites $(MN)$ et $(BC)$ sont \textbf{parallèles}.
\end{enumerate}
Si l'une de ces conditions manque, on n'a pas le droit d'écrire l'égalité des rapports. Dans une rédaction, ces hypothèses doivent toujours être citées explicitement.
\end{attention}

\begin{popfigure}
\begin{center}
\begin{tikzpicture}[scale=1.1,line cap=round,line join=round]
\coordinate (A) at (0,4);
\coordinate (B) at (-2.6,0);
\coordinate (C) at (3.2,0);
\coordinate (M) at ($(A)!0.55!(B)$);
\coordinate (N) at ($(A)!0.55!(C)$);
\draw[thick,popdark] (A)--(B)--(C)--cycle;
\draw[very thick,poporange] (M)--(N);
\draw[very thick,popblue] (A)--(M);
\draw[very thick,popblue] (A)--(N);
\draw[very thick,popgreen] (B)--(C);
\fill[popblue] (A) circle (1.6pt) node[above]{$A$};
\fill[popdark] (B) circle (1.6pt) node[below left]{$B$};
\fill[popdark] (C) circle (1.6pt) node[below right]{$C$};
\fill[poporange] (M) circle (1.6pt) node[left]{$M$};
\fill[poporange] (N) circle (1.6pt) node[right]{$N$};
\node[popblue,left] at ($(A)!0.5!(M)$) {\small $AM$};
\node[popblue,right] at ($(A)!0.5!(N)$) {\small $AN$};
\node[poporange,above] at ($(M)!0.5!(N)$) {\small $MN$};
\node[popgreen,below] at ($(B)!0.5!(C)$) {\small $BC$};
\node[align=center,popdark] at (4.2,2.4) {$\dfrac{AM}{AB}=\dfrac{AN}{AC}=\dfrac{MN}{BC}$};
\end{tikzpicture}
\end{center}
\legende{Configuration de Thalès : $(MN) \parallel (BC)$, les longueurs des côtés des triangles $AMN$ et $ABC$ sont proportionnelles.}
\end{popfigure}

\courssub{Pourquoi est-ce vrai ? Une justification intuitive}

\begin{experience}[Une démonstration guidée par les aires (résultat admis)]
En 3ème, la démonstration complète de la propriété de Thalès est \textbf{admise} : on ne te demandera pas de la refaire au BEPC. Mais en voici l'idée, pour que ce résultat ne tombe pas du ciel.

\textbf{Étape 1 : deux triangles de même aire.} Les triangles $MBC$ et $NBC$ ont le même côté $[BC]$ et leurs sommets $M$ et $N$ sont sur une droite parallèle à $(BC)$ : ils ont donc la même hauteur issue de $M$ (ou de $N$) relative à $(BC)$. Ils ont donc la \emph{même aire}. En retirant à chacun l'aire du triangle $OBC$ (où $O$ est l'intersection de $(BN)$ et $(CM)$), on obtient que les triangles $MOB$ et $NOC$ ont la même aire.

\textbf{Étape 2 : des aires proportionnelles aux bases.} L'aire d'un triangle de hauteur $h$ et de base $b$ vaut $\dfrac{b \times h}{2}$. Les triangles $AMN$ et $MBN$ ont la même hauteur issue de $N$, donc le rapport de leurs aires est $\dfrac{AM}{MB}$. De même, les triangles $AMN$ et $MNC$ ont la même hauteur issue de $M$, donc le rapport de leurs aires est $\dfrac{AN}{NC}$. Grâce à l'égalité des aires de l'étape 1, on en tire $\dfrac{AM}{MB} = \dfrac{AN}{NC}$, puis, en ajoutant $1$ aux deux membres, $\dfrac{AM}{AB} = \dfrac{AN}{AC}$.

\textbf{Étape 3 : le troisième rapport.} En traçant par $M$ la parallèle à $(AC)$, on obtient un parallélogramme qui permet de montrer que $\dfrac{MN}{BC}$ est égal aux deux rapports précédents.

\textbf{Conclusion.} Ce raisonnement complet est délicat : on \textbf{admet} donc la propriété de Thalès, et on retient surtout l'idée intuitive : \emph{le parallélisme crée la proportionnalité}.
\end{experience}

\begin{savaistu}[Thalès de Milet, un mathématicien... et un homme d'affaires]
Thalès de Milet (vers 625 -- 547 av. J.-C.) est considéré comme l'un des premiers mathématiciens de l'histoire. La légende raconte qu'il aurait mesuré la hauteur de la pyramide de Khéops en Égypte grâce à l'ombre portée d'un bâton : en comparant l'ombre du bâton à celle de la pyramide, il appliquait déjà la proportionnalité des longueurs dans des triangles semblables ! Aujourd'hui encore, au Cameroun, les géomètres qui bornent les terrains et les techniciens de CAMTEL qui vérifient l'alignement de poteaux utilisent, sans le savoir, les idées de Thalès.
\end{savaistu}

\courssub{Bien écrire les rapports : le choix du sommet commun}

Pour écrire correctement l'égalité de Thalès, il y a une méthode infaillible.

\begin{methode}[Écrire l'égalité des rapports de Thalès]
\begin{enumerate}
\item Repérer le \cle{sommet commun} aux deux triangles (ici $A$) : c'est le point où se rejoignent les deux droites sécantes.
\item Écrire le nom des deux triangles dans le même ordre de sommets : le petit triangle $AMN$ et le grand triangle $ABC$ (le sommet $A$ d'abord, puis le point sur $(AB)$, puis le point sur $(AC)$).
\item Écrire les trois rapports \emph{dans le même ordre}, en mettant toujours le même triangle au numérateur :
\[ \underbrace{\frac{AM}{AB}}_{\text{côtés sur }(AB)} = \underbrace{\frac{AN}{AC}}_{\text{côtés sur }(AC)} = \underbrace{\frac{MN}{BC}}_{\text{côtés parallèles}}. \]
\end{enumerate}
On peut aussi choisir de mettre le grand triangle au numérateur : $\dfrac{AB}{AM} = \dfrac{AC}{AN} = \dfrac{BC}{MN}$. L'essentiel est de \textbf{ne pas mélanger} : si le petit triangle est en haut dans un rapport, il doit être en haut dans les trois.
\end{methode}

\begin{attention}[Erreur fréquente : les rapports mal choisis]
On voit souvent des élèves écrire $\dfrac{AM}{MB} = \dfrac{AN}{NC}$. Cette égalité est \emph{vraie} (on peut la démontrer), mais elle \textbf{ne fait pas partie de l'énoncé direct} de la propriété de Thalès vu en 3ème, et elle est dangereuse : on y confond facilement $MB$ et $AB$. Au BEPC, privilégie toujours l'écriture du cours :
\[ \frac{AM}{AB} = \frac{AN}{AC} = \frac{MN}{BC}, \]
où chaque rapport compare un côté du \emph{petit} triangle au côté correspondant du \emph{grand} triangle. Autre piège : ne jamais écrire $\dfrac{AM}{AN} = \dfrac{AB}{AC}$ en croyant appliquer Thalès sans vérifier : cette égalité mélange les deux triangles et n'est pas l'énoncé direct.
\end{attention}

\courssub{Calculer une longueur : la quatrième proportionnelle}

L'égalité de Thalès donne \emph{trois} rapports égaux. En pratique, on n'utilise que \emph{deux} de ces rapports à la fois : celui qui contient l'inconnue et celui qui est entièrement connu. On obtient alors une égalité du type $\dfrac{a}{b} = \dfrac{c}{d}$ avec une seule inconnue, que l'on résout grâce aux \cle{produits en croix}.

\begin{aretenir}[Produits en croix]
Si $\dfrac{a}{b} = \dfrac{c}{d}$ (avec $b \neq 0$ et $d \neq 0$), alors $a \times d = b \times c$. On en déduit la \cle{quatrième proportionnelle} : par exemple, $a = \dfrac{b \times c}{d}$.
\end{aretenir}

\begin{methode}[Calculer une longueur avec la propriété de Thalès]
\begin{enumerate}
\item \textbf{Vérifier les hypothèses} : les points sont-ils alignés comme il faut ? Les droites sont-elles parallèles ? Le citer dans la rédaction.
\item \textbf{Écrire l'égalité des rapports} de Thalès (petit triangle / grand triangle).
\item \textbf{Choisir les deux rapports utiles} : celui qui contient l'inconnue et celui dont on connaît les deux longueurs.
\item \textbf{Remplacer par les valeurs connues}, puis appliquer les \textbf{produits en croix} pour trouver l'inconnue.
\item \textbf{Conclure} par une phrase avec l'unité.
\end{enumerate}
\end{methode}

\begin{exemplebox}[Premier calcul : trouver $AC$]
Dans le triangle $ABC$, $M \in [AB]$, $N \in [AC]$ et $(MN) \parallel (BC)$. On donne $AM = 3$ cm, $AB = 5$ cm et $AN = 4$ cm. Calculons $AC$.

Les points $A$, $M$, $B$ d'une part et $A$, $N$, $C$ d'autre part sont alignés, et $(MN) \parallel (BC)$. D'après la propriété de Thalès :
\[ \frac{AM}{AB} = \frac{AN}{AC} \quad \text{donc} \quad \frac{3}{5} = \frac{4}{AC}. \]
Par les produits en croix : $3 \times AC = 5 \times 4 = 20$, d'où $AC = \dfrac{20}{3} \approx 6{,}7$ cm.
\end{exemplebox}

\begin{exemplebox}[Deuxième calcul : trouver $MN$]
Sur la même figure, on donne de plus $BC = 7{,}5$ cm. Calculons $MN$.

On utilise cette fois le rapport connu $\dfrac{AM}{AB} = \dfrac{3}{5}$ et le rapport des côtés parallèles :
\[ \frac{MN}{BC} = \frac{AM}{AB} \quad \text{donc} \quad \frac{MN}{7{,}5} = \frac{3}{5}. \]
Par les produits en croix : $5 \times MN = 3 \times 7{,}5 = 22{,}5$, d'où $MN = \dfrac{22{,}5}{5} = 4{,}5$ cm.
\end{exemplebox}

\begin{exoresolu}[Un calcul complet rédigé comme au BEPC]
\textbf{Énoncé.} Soit $ABC$ un triangle tel que $AB = 7$ cm, $AC = 9{,}3$ cm et $BC = 6$ cm. On place le point $M$ sur $[AB]$ tel que $AM = 3$ cm, puis on trace la parallèle à $(BC)$ passant par $M$ ; elle coupe $[AC]$ en $N$. Calculer $AN$ et $MN$ (on donnera une valeur approchée au dixième pour $AN$).
\tcblower
\textbf{Solution.} Les points $A$, $M$, $B$ sont alignés, les points $A$, $N$, $C$ sont alignés, et $(MN) \parallel (BC)$ par construction. On peut donc appliquer la propriété directe de Thalès :
\[ \frac{AM}{AB} = \frac{AN}{AC} = \frac{MN}{BC}. \]
\textbf{Calcul de $AN$.} On utilise $\dfrac{AM}{AB} = \dfrac{AN}{AC}$, c'est-à-dire $\dfrac{3}{7} = \dfrac{AN}{9{,}3}$. Par les produits en croix :
\[ AN = \frac{3 \times 9{,}3}{7} = \frac{27{,}9}{7} \approx 4{,}0 \text{ cm}. \]
\textbf{Calcul de $MN$.} On utilise $\dfrac{AM}{AB} = \dfrac{MN}{BC}$, c'est-à-dire $\dfrac{3}{7} = \dfrac{MN}{6}$. Par les produits en croix :
\[ MN = \frac{3 \times 6}{7} = \frac{18}{7} \approx 2{,}6 \text{ cm}. \]
\textbf{Conclusion.} $AN \approx 4{,}0$ cm et $MN \approx 2{,}6$ cm.
\end{exoresolu}

\courssub{Cas particulier : le théorème des milieux}

Que se passe-t-il lorsque $M$ est le \emph{milieu} de $[AB]$ ? Alors $\dfrac{AM}{AB} = \dfrac{1}{2}$. La propriété de Thalès donne aussitôt $\dfrac{AN}{AC} = \dfrac{1}{2}$, donc $N$ est le milieu de $[AC]$, et $\dfrac{MN}{BC} = \dfrac{1}{2}$, donc $MN = \dfrac{BC}{2}$. Ce résultat, très utile, porte un nom.

\begin{propriete}[Théorème des milieux]
Dans un triangle $ABC$, si $M$ est le milieu de $[AB]$ et $N$ le milieu de $[AC]$, alors :
\begin{itemize}
\item la droite $(MN)$ est \cle{parallèle} à $(BC)$ ;
\item $MN = \dfrac{BC}{2}$ : le segment joignant les milieux mesure la \emph{moitié} du troisième côté.
\end{itemize}
Ce théorème est une conséquence directe de la propriété de Thalès (le coefficient de proportionnalité vaut $\dfrac{1}{2}$).
\end{propriete}

\begin{popfigure}
\begin{center}
\begin{tikzpicture}[scale=1.1,line cap=round,line join=round]
\coordinate (A) at (0,3.6);
\coordinate (B) at (-2.8,0);
\coordinate (C) at (3.4,0);
\coordinate (M) at ($(A)!0.5!(B)$);
\coordinate (N) at ($(A)!0.5!(C)$);
\draw[thick,popdark] (A)--(B)--(C)--cycle;
\draw[very thick,poporange] (M)--(N);
\fill[popblue] (A) circle (1.6pt) node[above]{$A$};
\fill[popdark] (B) circle (1.6pt) node[below left]{$B$};
\fill[popdark] (C) circle (1.6pt) node[below right]{$C$};
\fill[poporange] (M) circle (1.6pt) node[left]{$M$};
\fill[poporange] (N) circle (1.6pt) node[right]{$N$};
% codages milieux
\draw[popblue,thick] ($(A)!0.5!(M)$) ++(-0.09,-0.12) -- ++(0.18,0.24);
\draw[popblue,thick] ($(M)!0.5!(B)$) ++(-0.09,-0.12) -- ++(0.18,0.24);
\draw[popgreen,thick] ($(A)!0.5!(N)$) ++(-0.12,-0.09) -- ++(0.24,0.18);
\draw[popgreen,thick] ($(A)!0.5!(N)$) ++(-0.10,-0.13) -- ++(0.24,0.18);
\draw[popgreen,thick] ($(N)!0.5!(C)$) ++(-0.12,-0.09) -- ++(0.24,0.18);
\draw[popgreen,thick] ($(N)!0.5!(C)$) ++(-0.10,-0.13) -- ++(0.24,0.18);
\node[poporange,above] at ($(M)!0.5!(N)$) {\small $MN=\dfrac{BC}{2}$};
\node[popdark,below] at ($(B)!0.5!(C)$) {\small $BC$};
\end{tikzpicture}
\end{center}
\legende{Théorème des milieux : $M$ et $N$ milieux (codages identiques), donc $(MN) \parallel (BC)$ et $MN = \dfrac{BC}{2}$.}
\end{popfigure}

\begin{exemplebox}[Application rapide du théorème des milieux]
Dans un triangle $ABC$ avec $BC = 8$ cm, $M$ est le milieu de $[AB]$ et $N$ le milieu de $[AC]$. Alors $(MN) \parallel (BC)$ et $MN = \dfrac{8}{2} = 4$ cm. Aucun produit en croix n'est nécessaire : le coefficient $\dfrac{1}{2}$ fait tout le travail.
\end{exemplebox}

\courssub{La rédaction type à connaître par cœur}

Au BEPC, la présentation compte autant que le résultat. Voici le modèle de rédaction attendu, en quatre temps :

\begin{aretenir}[Rédaction type d'un exercice de Thalès]
\begin{enumerate}
\item \textbf{Hypothèses} : « Les points $A$, $M$, $B$ sont alignés, les points $A$, $N$, $C$ sont alignés, et $(MN) \parallel (BC)$. »
\item \textbf{Propriété} : « D'après la propriété directe de Thalès : $\dfrac{AM}{AB} = \dfrac{AN}{AC} = \dfrac{MN}{BC}$. »
\item \textbf{Calcul} : on remplace par les valeurs connues et on utilise les produits en croix.
\item \textbf{Conclusion} : « Donc $AC = \ldots$ cm. »
\end{enumerate}
Oublier de citer les hypothèses (en particulier le parallélisme) fait perdre des points, même si le calcul est juste !
\end{aretenir}

\begin{exercice}[À toi de jouer]
Dans le triangle $ABC$, $M \in [AB]$, $N \in [AC]$ et $(MN) \parallel (BC)$. On donne $AM = 2$ cm, $AB = 6$ cm et $AC = 9$ cm.
\begin{enumerate}
\item Calculer $AN$.
\item Sachant que $BC = 12$ cm, calculer $MN$.
\end{enumerate}
\end{exercice}

\begin{corrige}[Exercice : À toi de jouer]
Les points $A$, $M$, $B$ sont alignés, les points $A$, $N$, $C$ sont alignés, et $(MN) \parallel (BC)$. D'après la propriété directe de Thalès : $\dfrac{AM}{AB} = \dfrac{AN}{AC} = \dfrac{MN}{BC}$.
\begin{enumerate}
\item $\dfrac{2}{6} = \dfrac{AN}{9}$, donc $AN = \dfrac{2 \times 9}{6} = 3$ cm.
\item $\dfrac{2}{6} = \dfrac{MN}{12}$, donc $MN = \dfrac{2 \times 12}{6} = 4$ cm.
\end{enumerate}
\end{corrige}

La propriété directe de Thalès est maintenant ton outil principal pour calculer des longueurs dans la configuration du triangle. Dans la section suivante, nous verrons comment cet outil permet de résoudre des problèmes concrets de la vie courante : hauteur d'un arbre, largeur d'une rivière, plans d'une maison à Yaoundé...

\coursec{Applications à la vie courante}

Dans la section précédente, nous avons découvert la \cle{propriété directe de Thalès} : lorsque deux droites sont parallèles dans une configuration de Thalès, les longueurs des segments sont proportionnelles. Mais à quoi cela sert-il concrètement ? Eh bien, cette propriété est l'un des outils les plus puissants de toute la géométrie : elle permet de \textbf{mesurer des distances inaccessibles} (la hauteur d'un mât, la largeur d'un fleuve) sans jamais les toucher, de \textbf{reproduire des figures à l'échelle} (plans, photos) et de \textbf{partager équitablement} des longueurs ou des sommes. Cette section te montre comment passer d'une situation de la vie courante à une configuration de Thalès, puis comment rédiger une solution complète, comme on te le demandera au BEPC.

\courssub{Modéliser une situation concrète : la démarche générale}

\textbf{Intuition.} Imagine que tu veux connaître la hauteur du mât du drapeau de ton collège. Impossible de grimper avec un mètre ruban ! Mais un jour de soleil, le mât projette une ombre au sol, et un simple bâton planté verticalement projette lui aussi une ombre. Les rayons du soleil sont parallèles entre eux : le mât et son ombre d'une part, le bâton et son ombre d'autre part, forment alors une \emph{configuration de Thalès}. Le sommet commun, c'est l'extrémité des ombres ; les deux droites parallèles, ce sont le mât et le bâton (tous deux verticaux). Il ne reste plus qu'à appliquer la proportionnalité.

Traduire ainsi une situation réelle en configuration de Thalès s'appelle \cle{modéliser}. C'est la compétence essentielle de cette leçon.

\begin{methode}[Modéliser une situation concrète par Thalès]
\begin{enumerate}
\item \textbf{Faire un schéma} de la situation, en codant les données (longueurs connues, angles droits, alignements).
\item \textbf{Repérer le sommet commun} (le point où se croisent les deux droites sécantes) et \textbf{les deux droites parallèles} de la configuration. Vérifier et justifier le parallélisme (par exemple : deux droites perpendiculaires à une même troisième droite sont parallèles).
\item \textbf{Appliquer la propriété directe de Thalès} : écrire l'égalité des trois rapports, puis isoler la longueur inconnue par un produit en croix.
\item \textbf{Interpréter le résultat} dans le contexte : donner la valeur avec son unité et rédiger une phrase de conclusion qui répond à la question posée.
\end{enumerate}
\end{methode}

\begin{attention}[Unités : le piège classique]
Avant tout calcul, \textbf{toutes les longueurs doivent être exprimées dans la même unité}. L'erreur la plus fréquente consiste à mélanger des mètres et des centimètres : par exemple, écrire $\dfrac{1{,}5}{180}$ alors que l'ombre mesure 1,8 m et la hauteur 1,5 m. Convertis d'abord : soit tout en mètres ($1{,}8$ m), soit tout en centimètres ($180$ cm). Le résultat sera alors dans l'unité choisie.
\end{attention}

\courssub{Le problème du mât : un exercice résolu complet}

\begin{exoresolu}[La hauteur du mât du collège]
Un jour ensoleillé, le mât du drapeau d'un collège de Douala projette au sol une ombre de $12$ m. Au même instant, un bâton vertical de $1{,}5$ m projette une ombre de $1{,}8$ m. Calculer la hauteur du mât.
\tcblower
\textbf{Étape 1 : schéma et hypothèses.} On note $S$ l'extrémité commune des deux ombres, $B$ le pied du mât, $T$ son sommet, $C$ le pied du bâton et $D$ son sommet. Les points $S$, $C$, $B$ sont alignés (au sol) ainsi que $S$, $D$, $T$ (le long du rayon de soleil). Le mât et le bâton sont verticaux, donc tous deux perpendiculaires au sol : $(BT) \perp (SB)$ et $(CD) \perp (SB)$, d'où $(CD) \parallel (BT)$.

\textbf{Étape 2 : configuration de Thalès.} Les triangles $SCD$ et $SBT$ sont en configuration de Thalès, de sommet commun $S$, avec $(CD) \parallel (BT)$.

\textbf{Étape 3 : application de la propriété directe de Thalès.}
\[ \frac{SC}{SB} = \frac{SD}{ST} = \frac{CD}{BT} \quad\text{soit}\quad \frac{1{,}8}{12} = \frac{1{,}5}{BT}. \]
Par le produit en croix :
\[ BT = \frac{1{,}5 \times 12}{1{,}8} = \frac{18}{1{,}8} = 10. \]

\textbf{Étape 4 : conclusion.} Le mât du collège mesure \textbf{10 mètres} de haut.
\end{exoresolu}

\begin{exemplebox}[Vérification rapide]
On peut contrôler la cohérence : le rapport de réduction est $\dfrac{1{,}8}{12} = 0{,}15$, donc le bâton est $0{,}15$ fois le mât : $10 \times 0{,}15 = 1{,}5$ m. C'est bien la hauteur du bâton : le résultat est cohérent.
\end{exemplebox}

\courssub{Application 1 : mesurer la largeur d'une rivière sans la traverser}

Comment mesurer la largeur du fleuve Wouri sans bateau ni pont ? Les géomètres utilisent une méthode classique basée sur Thalès : on construit sur la rive accessible un triangle rectangle semblable à celui formé par la largeur du fleuve.

\begin{popfigure}
\begin{center}
\begin{tikzpicture}[scale=0.9,>=Stealth]
% Rive opposée (hachurée)
\fill[popblueL] (-0.5,2.2) rectangle (12.5,3.8);
\draw[popblue,thick] (-0.5,2.2) -- (12.5,2.2);
\node[popblue] at (6,3.0) {\small Fleuve Wouri};
% Rive proche
\draw[popdark,very thick] (-0.5,0) -- (12.5,0);
\node[popdark,below left] at (-0.5,0) {\small Rive accessible};
% Points
\coordinate (A) at (2,2.2);    % Arbre rive opposée
\coordinate (B) at (2,0);      % Piquet face à A
\coordinate (C) at (6,0);      % Point C sur la rive (BC = 20 m)
\coordinate (D) at (10,0);     % Point D sur la rive (CD = 20 m)
\coordinate (M) at (10,-2.2);  % Point M sur la perpendiculaire en D, aligné A-C-M
% Perpendiculaires
\draw[poporange,very thick] (A) -- (B); % Largeur AB
\draw[popgreen,very thick] (D) -- (M);  % Segment DM mesuré
% Alignements
\draw[popink,thick] (A) -- (M);         % Visée A-C-M
\draw[popdark,thick] (B) -- (D);        % Rive B-C-D
% Codage parallélisme (AB // DM)
\draw[poporange] (1.8,2.35) -- (2.2,2.35);
\draw[poporange] (1.8,-0.15) -- (2.2,-0.15);
\draw[popgreen] (9.8,-2.35) -- (10.2,-2.35);
\draw[popgreen] (9.8,0.15) -- (10.2,0.15);
% Codage alignement B-C-D
\draw[popdark] (2,-0.3) -- (2.2,-0.3);
\draw[popdark] (6,-0.3) -- (6.2,-0.3);
\draw[popdark] (10,-0.3) -- (10.2,-0.3);
% Points et labels
\foreach \p/\pos/\nom in {A/above/$A$ (arbre), B/below/$B$, C/below/$C$, D/below/$D$, M/below/$M$}
  \fill[popink] (\p) circle (1.8pt) node[\pos] {\nom};
% Mesures
\draw[<->,poporange] (2.4,2.2) -- (2.4,0) node[midway,right] {\small $AB = ?$};
\draw[<->,popdark] (2,-0.5) -- (6,-0.5) node[midway,below] {\small $BC = 20$ m};
\draw[<->,popdark] (6,-0.5) -- (10,-0.5) node[midway,below] {\small $CD = 20$ m};
\draw[<->,popgreen] (10.4,0) -- (10.4,-2.2) node[midway,right] {\small $DM$ (mesuré)};
% Angle droit en B et D
\draw[popmuted] (1.85,0) -- (1.85,0.15) -- (2,0.15);
\draw[popmuted] (9.85,0) -- (9.85,0.15) -- (10,0.15);
\end{tikzpicture}
\end{center}
\legende{Mesure de la largeur $AB$ : on choisit $C$ et $D$ sur la rive tels que $BC = CD$. En $D$, on lève la perpendiculaire et on place $M$ tel que $A$, $C$, $M$ soient alignés. Alors $(AB) \parallel (DM)$ et Thalès au sommet $C$ donne $AB = DM$.}
\end{popfigure}

\begin{exemplebox}[Calcul de la largeur du fleuve]
\textbf{Construction :}
\begin{enumerate}
\item Planter un piquet $B$ sur la rive, face à l'arbre $A$ (largeur $AB$ inconnue).
\item Sur la rive, choisir $C$ et $D$ alignés avec $B$ de sorte que $BC = CD$ (par exemple $20$ m chacun).
\item En $D$, tracer la perpendiculaire à la rive (vers l'intérieur des terres).
\item Sur cette perpendiculaire, placer $M$ de manière que $A$, $C$, $M$ soient alignés (visée). Mesurer $DM$.
\end{enumerate}
\textbf{Justification :} $(AB)$ et $(DM)$ sont toutes deux perpendiculaires à la rive $(BD)$, donc $(AB) \parallel (DM)$. Les points $A$, $C$, $M$ sont alignés par construction ; $B$, $C$, $D$ sont alignés sur la rive. Les triangles $ABC$ et $MDC$ sont en configuration de Thalès de sommet $C$.
\textbf{Application de Thalès :}
\[ \frac{AB}{DM} = \frac{BC}{CD}. \]
Comme on a choisi $BC = CD$, le rapport vaut $1$, donc \textbf{$AB = DM$}.
Il suffit de mesurer $DM$ pour connaître la largeur du fleuve. Si $DM = 20$ m, le fleuve fait \textbf{20 m} de large.
\end{exemplebox}

\courssub{Application 2 : hauteur d'un arbre ou d'un immeuble à Yaoundé}

La méthode du mât se généralise : pour mesurer la hauteur d'un fromager dans la cour du collège ou d'un immeuble du centre-ville de Yaoundé, deux techniques sont au programme.

\textbf{Méthode de l'ombre.} C'est exactement le problème du mât résolu plus haut : on compare l'ombre de l'objet à celle d'un bâton de hauteur connue, au même instant (pour que les rayons du soleil soient les mêmes).

\textbf{Méthode du bâton (visée à hauteur d'œil).} On place un bâton vertical entre l'œil et l'immeuble, de sorte que le sommet du bâton, l'œil et le sommet de l'immeuble soient alignés. L'œil est le sommet commun de la configuration de Thalès ; le bâton et l'immeuble, tous deux verticaux, sont les deux parallèles. En mesurant la distance œil--bâton et la distance œil--immeuble, on calcule la hauteur cherchée par proportionnalité.

\begin{attention}[Conditions d'application]
Pour que la méthode de l'ombre soit valable, il faut : (1) prendre les deux mesures \emph{au même instant} (le soleil tourne !) ; (2) que le sol soit plat et horizontal ; (3) que l'objet et le bâton soient bien verticaux, donc parallèles entre eux. Sans parallélisme justifié, \textbf{pas de Thalès}.
\end{attention}

\courssub{Application 3 : agrandissement et réduction de plans}

Quand on photocopie un document à 150 \% ou qu'on dessine le plan d'une parcelle à l'échelle $\dfrac{1}{200}$, on effectue un \cle{agrandissement} ou une \cle{réduction} : toutes les longueurs sont multipliées par un même nombre $k$, appelé \cle{rapport} (ou échelle). Si $k > 1$, c'est un agrandissement ; si $0 < k < 1$, une réduction. Les angles, eux, ne changent pas : la figure garde la même forme. C'est encore Thalès qui se cache derrière : les longueurs correspondantes sont proportionnelles.

\begin{popfigure}
\begin{center}
\begin{tikzpicture}[scale=1.0,>=Stealth]
% Photo initiale
\draw[popblue,very thick,fill=popblueL] (0,0) rectangle (2,1.4);
\node[popblue] at (1,0.7) {\small photo};
\draw[<->,popblue] (0,-0.3) -- (2,-0.3) node[midway,below] {\small 10 cm};
\draw[<->,popblue] (-0.3,0) -- (-0.3,1.4) node[midway,left] {\small 7 cm};
% Flèche agrandissement
\draw[->,poporange,very thick] (2.4,0.7) -- (3.8,0.7) node[midway,above] {\small $k = 1{,}5$};
% Photo agrandie
\draw[poporange,very thick,fill=poporangeL] (4.2,0) rectangle (7.2,2.1);
\node[poporange] at (5.7,1.05) {\small agrandie};
\draw[<->,poporange] (4.2,-0.3) -- (7.2,-0.3) node[midway,below] {\small 15 cm};
\draw[<->,poporange] (7.5,0) -- (7.5,2.1) node[midway,right] {\small 10{,}5 cm};
\end{tikzpicture}
\end{center}
\legende{Agrandissement d'une photo rectangulaire de rapport $k = 1{,}5$ : chaque longueur est multipliée par $1{,}5$ ($10 \times 1{,}5 = 15$ cm et $7 \times 1{,}5 = 10{,}5$ cm).}
\end{popfigure}

\begin{exemplebox}[Plan d'une parcelle à l'échelle]
Sur le plan d'une parcelle à Nkongsamba, dessiné à l'échelle $\dfrac{1}{500}$, le terrain mesure $8$ cm de long. La longueur réelle est $8 \times 500 = 4000$ cm $= 40$ m. Inversement, une façade réelle de $25$ m ($2500$ cm) est représentée par $2500 \div 500 = 5$ cm sur le plan.
\end{exemplebox}

\courssub{Application 4 : partage d'un segment en parties proportionnelles}

\textbf{Situation.} Trois frères veulent partager un terrain rectangulaire bordé par une route en trois lots de même largeur de façade. Comment tracer les limites au cordeau, sans instrument de mesure sophistiqué ? La réponse est un classique de la géométrie : le \textbf{partage d'un segment en parties égales} grâce à des droites parallèles.

\begin{methode}[{Partager un segment $[AB]$ en 3 parties égales}]
\begin{enumerate}
\item Tracer une demi-droite $[Ax)$ quelconque (différente de $(AB)$).
\item Reporter sur $[Ax)$, au compas, trois fois la même longueur : $AM_1 = M_1M_2 = M_2M_3$.
\item Tracer la droite $(BM_3)$.
\item Tracer les parallèles à $(BM_3)$ passant par $M_1$ et $M_2$ : elles coupent $[AB]$ en $P$ et $Q$.
\item Alors $AP = PQ = QB$ : le segment est partagé en trois parties égales.
\end{enumerate}
La justification est la propriété directe de Thalès appliquée dans le triangle $ABM_3$ : le parallélisme des droites transmet la proportionnalité des partages.
\end{methode}

\begin{popfigure}
\begin{center}
\begin{tikzpicture}[scale=1.0,>=Stealth]
\coordinate (A) at (0,0);
\coordinate (B) at (6,0);
\coordinate (M1) at (1.1,1.1);
\coordinate (M2) at (2.2,2.2);
\coordinate (M3) at (3.3,3.3);
% Demi-droite [Ax)
\draw[popdark,thick] (A) -- (4,3.85) node[above left] {\small $x$};
% Segment AB
\draw[popink,very thick] (A) -- (B);
% Points de report
\foreach \p/\nom in {M1/$M_1$, M2/$M_2$, M3/$M_3$}
  \fill[popblue] (\p) circle (1.8pt) node[above left] {\nom};
% Codage des reports égaux
\foreach \a/\b in {A/M1, M1/M2, M2/M3}
  \draw[popblue] ($(\a)!0.5!(\b)+(0.09,-0.09)$) -- ($(\a)!0.5!(\b)+(-0.09,0.09)$);
% Droite (BM3) et parallèles
\draw[poporange,thick] (B) -- (M3);
\coordinate (P) at (2,0);
\coordinate (Q) at (4,0);
\draw[poporange,thick,dashed] (P) -- (M1);
\draw[poporange,thick,dashed] (Q) -- (M2);
% Points P, Q, A, B
\fill[popink] (A) circle (1.8pt) node[below left] {$A$};
\fill[popink] (B) circle (1.8pt) node[below right] {$B$};
\fill[popgreen] (P) circle (1.8pt) node[below] {$P$};
\fill[popgreen] (Q) circle (1.8pt) node[below] {$Q$};
% Accolades de partage
\draw[<->,popgreen] (0,-0.55) -- (2,-0.55) node[midway,below] {\small $AP$};
\draw[<->,popgreen] (2,-0.55) -- (4,-0.55) node[midway,below] {\small $PQ$};
\draw[<->,popgreen] (4,-0.55) -- (6,-0.55) node[midway,below] {\small $QB$};
\end{tikzpicture}
\end{center}
\legende{Partage de $[AB]$ en trois parties égales : les droites $(M_1P)$, $(M_2Q)$ et $(M_3B)$ sont parallèles, donc $AP = PQ = QB$.}
\end{popfigure}

\begin{exemplebox}[Partage d'une somme d'argent]
Le même principe de proportionnalité permet de partager une somme de \num{60000} FCFA entre trois personnes proportionnellement aux nombres 2, 3 et 5 (par exemple selon les jours travaillés). Le total des parts est $2+3+5 = 10$. Chaque personne reçoit :
\[ \frac{2}{10}\times 60000 = 12000 \text{ F} ;\quad \frac{3}{10}\times 60000 = 18000 \text{ F} ;\quad \frac{5}{10}\times 60000 = 30000 \text{ F}. \]
Vérification : $12000 + 18000 + 30000 = 60000$ F. Le partage est correct.
\end{exemplebox}

\courssub{Rédiger une solution complète : le modèle BEPC}

Au BEPC, une question de géométrie utilisant Thalès est notée sur la \textbf{rigueur de la rédaction} autant que sur le calcul. Voici la structure attendue, que tu dois reproduire à chaque fois :

\begin{aretenir}[Les 5 éléments d'une solution complète]
\begin{enumerate}
\item Un \textbf{schéma codé} : points nommés, longueurs connues reportées, parallélisme codé par des flèches.
\item Les \textbf{hypothèses} énoncées : alignement des points, parallélisme des droites (avec sa justification).
\item La \textbf{propriété citée} : « D'après la propriété directe de Thalès... » suivie de l'égalité des rapports.
\item Le \textbf{calcul} : produit en croix, avec des longueurs dans la même unité.
\item La \textbf{réponse} : valeur avec unité et phrase de conclusion répondant à la question posée.
\end{enumerate}
\end{aretenir}

\begin{attention}[Erreurs fréquentes à éviter]
\begin{itemize}
\item Appliquer Thalès sans avoir \emph{justifié} le parallélisme des deux droites : la propriété ne s'applique pas !
\item Écrire les rapports « en croix » de travers : vérifie que chaque rapport compare des côtés \emph{correspondants} (ceux qui se font face dans les deux triangles).
\item Oublier l'unité dans la réponse finale, ou donner une réponse sans phrase de conclusion.
\item Mélanger les unités dans le calcul (m et cm) : convertis tout avant de poser l'égalité des rapports.
\end{itemize}
\end{attention}

\begin{savaistu}[Thalès, métier : géomètre-expert]
Le principe que tu viens d'apprendre est vieux de plus de 2500 ans : la légende raconte que Thalès de Milet mesura la hauteur de la pyramide de Khéops en Égypte grâce à son ombre, exactement comme notre mât. Aujourd'hui encore, au Cameroun, les \textbf{géomètres-experts} utilisent ce principe de proportionnalité pour le bornage des terrains, la topographie des parcelles et l'établissement des titres fonciers. Les instruments ont changé (théodolites, GPS), mais les mathématiques derrière sont celles de ton cours de 3ème !
\end{savaistu}

\begin{exercice}[À toi de jouer]
Un poteau électrique de la campagne de Bafoussam projette une ombre de $7{,}5$ m. Au même instant, un piquet vertical de $1{,}2$ m projette une ombre de $1{,}5$ m. Calcule la hauteur du poteau, en rédigeant ta solution selon les 5 éléments du modèle BEPC.
\end{exercice}

\coursec{Propriété réciproque de Thalès}

\courssub{De l'application au raisonnement inverse}

Dans la section précédente, nous avons vu comment la propriété \emph{directe} de Thalès permet de \textbf{calculer une longueur inconnue} lorsque l'on sait déjà que deux droites sont parallèles. Mais dans la vie courante, comme en géométrie, on se trouve souvent face à la situation inverse : \textbf{on connaît des mesures, et l'on veut prouver que deux droites sont parallèles}.

Imagine un menuisier à Yaoundé qui pose une étagère. Il a mesuré les équerres : $AM = \SI{40}{cm}$, $AB = \SI{60}{cm}$ d'un côté, et $AN = \SI{50}{cm}$, $AC = \SI{75}{cm}$ de l'autre. Avant de visser définitivement, il veut être \emph{sûr} que l'étagère $(MN)$ est bien parallèle au sol $(BC)$. Comment faire sans équerre de menuisier ni niveau à bulle ? C'est exactement le rôle de la \textbf{propriété réciproque de Thalès} : transformer une égalité de rapports en une certitude de parallélisme.

\courssub{Énoncé de la propriété réciproque}

\begin{propriete}[Réciproque de Thalès — Admise en classe de 3ème]
Soit un triangle $ABC$. Soient $M$ un point du segment $[AB]$ et $N$ un point du segment $[AC]$, \textbf{distincts de $A$}.
Si les points $A$, $M$, $B$ sont alignés \textbf{dans cet ordre} et les points $A$, $N$, $C$ sont alignés \textbf{dans cet ordre}, et si l'égalité des rapports est vérifiée :
\[
\frac{AM}{AB} = \frac{AN}{AC}
\]
alors les droites $(MN)$ et $(BC)$ sont \textbf{parallèles}.
\end{propriete}

\begin{attention}[L'ordre des points : une condition \textbf{indispensable}]
L'hypothèse « $A$, $M$, $B$ alignés \textbf{dans cet ordre} » (et de même pour $A$, $N$, $C$) signifie que $M$ appartient au \textbf{segment} $[AB]$ et $N$ au \textbf{segment} $[AC]$. Autrement dit, $M$ est \emph{entre} $A$ et $B$, et $N$ est \emph{entre} $A$ et $C$.

Si cet ordre n'est pas respecté (par exemple si $M$ est sur la droite $(AB)$ mais \emph{en dehors} du segment $[AB]$), l'égalité des rapports \textbf{ne garantit plus} le parallélisme. C'est l'erreur la plus fréquente au BEPC : calculer les rapports, trouver l'égalité, et conclure trop vite sans vérifier où se placent réellement les points.
\end{attention}

\courssub{Contre-exemple : quand l'ordre n'est pas respecté}

\begin{popfigure}
\begin{tikzpicture}[scale=1.5, >=Stealth]
% Triangle ABC : A(0,0), B(4,0), C(0,3)
\coordinate (A) at (0,0);
\coordinate (B) at (4,0);
\coordinate (C) at (0,3);
% M on segment [AB] : AM=1, AB=4 -> ratio 1/4
\coordinate (M) at (1,0);
% N on line AC but BEYOND A (order N-A-C). AN=0.75, AC=3 -> ratio 1/4
\coordinate (N) at (0,-0.75);

% Triangle ABC
\draw[popdark, thick] (A) -- (B) -- (C) -- cycle;
% Line AC extended (dashed)
\draw[popmuted, dashed] (N) -- (C);
% Segment MN
\draw[poporange, thick] (M) -- (N);

% Points
\fill[popblue] (A) circle (2pt) node[below left] {$A$};
\fill[popblue] (B) circle (2pt) node[below right] {$B$};
\fill[popblue] (C) circle (2pt) node[above left] {$C$};
\fill[poporange] (M) circle (2pt) node[below] {$M$};
\fill[poporange] (N) circle (2pt) node[left] {$N$};

% Ratio labels
\node[popblue, align=center, font=\small] at (2, -0.7) {$\dfrac{AM}{AB} = \dfrac{1}{4}$};
\node[popblue, align=center, font=\small] at (-1.3, 1.2) {$\dfrac{AN}{AC} = \dfrac{0,75}{3} = \dfrac{1}{4}$};
\node[popred, align=center, font=\bfseries\small] at (2.5, 1.5) {Rapports égaux \\ mais $(MN) \not\parallel (BC)$ \\ (ordre $N$-$A$-$C$ non respecté)};
\end{tikzpicture}
\legende{Contre-exemple : $M \in [AB]$ mais $N \notin [AC]$ (ordre $N$-$A$-$C$). Les rapports sont égaux ($\frac{1}{4}$), pourtant $(MN)$ coupe $(BC)$ : elles ne sont pas parallèles. L'hypothèse d'ordre est violée.}
\end{popfigure}

\courssub{Pourquoi cela marche-t-il ? (Démonstration guidée, admise en 3ème)}

Bien que la démonstration rigoureuse fasse appel à des axiomes du niveau supérieur, on peut en donner une version intuitive et convaincante, admise au BEPC, qui repose sur \textbf{l'unicité de la parallèle}.

\begin{experience}[Démonstration guidée de la réciproque]
On considère un triangle $ABC$ et des points $M \in [AB]$, $N \in [AC]$ tels que $\frac{AM}{AB} = \frac{AN}{AC}$.

\begin{enumerate}
    \item \textbf{Construction :} Par le point $M$, on trace la parallèle à $(BC)$. Elle coupe $(AC)$ en un point que l'on note $N'$.
    \item \textbf{Application de la directe :} Dans le triangle $ABC$, comme $(MN') \parallel (BC)$ et $M \in [AB]$, la propriété \emph{directe} de Thalès nous donne :
    \[
    \frac{AM}{AB} = \frac{AN'}{AC}
    \]
    \item \textbf{Comparaison :} Par hypothèse, on a $\frac{AM}{AB} = \frac{AN}{AC}$. Donc :
    \[
    \frac{AN'}{AC} = \frac{AN}{AC} \quad \Rightarrow \quad AN' = AN
    \]
    \item \textbf{Conclusion :} Les points $N$ et $N'$ appartiennent tous les deux à la demi-droite $[AC)$ (car $A,N,C$ et $A,N',C$ sont alignés dans le même ordre) et sont à la même distance de $A$. \textbf{Ils coïncident donc : $N = N'$.}
    \item Comme $(MN') \parallel (BC)$ par construction et $N=N'$, la droite $(MN)$ est bien parallèle à $(BC)$.
\end{enumerate}
\emph{C'est cette logique — unicité du point sur une demi-droite à une distance donnée — qui valide la réciproque. En 3ème, on admet le résultat en citant le théorème.}
\end{experience}

\courssub{Méthode : Démontrer que deux droites sont parallèles}

\begin{methode}[Vérifier le parallélisme par la réciproque de Thalès]
Pour prouver que $(MN) \parallel (BC)$ dans un triangle $ABC$ avec $M \in (AB)$ et $N \in (AC)$ :
\begin{enumerate}
    \item \textbf{Vérifier l'alignement et l'ordre :} Écrire explicitement : « $A$, $M$, $B$ sont alignés dans cet ordre » et « $A$, $N$, $C$ sont alignés dans cet ordre ». (Cela veut dire $M \in [AB]$ et $N \in [AC]$).
    \item \textbf{Calculer les deux rapports séparément :} Calculer $\frac{AM}{AB}$ et $\frac{AN}{AC}$. Simplifier les fractions si possible.
    \item \textbf{Comparer :} Montrer que les deux fractions sont égales (même valeur décimale, ou même fraction irréductible, ou égalité des produits en croix : $AM \times AC = AN \times AB$).
    \item \textbf{Conclure :} Écrire : « D'après la réciproque de Thalès, comme les points sont alignés dans le même ordre et $\frac{AM}{AB} = \frac{AN}{AC}$, les droites $(MN)$ et $(BC)$ sont parallèles. »
\end{enumerate}
\end{methode}

\courssub{Exemple chiffré détaillé}

\begin{exemplebox}[Vérification de parallélisme]
Soit un triangle $ABC$. On place $M$ sur $[AB]$ et $N$ sur $[AC]$ tels que :
\[
AM = \SI{2}{cm},\quad AB = \SI{5}{cm},\quad AN = \SI{3}{cm},\quad AC = \SI{7,5}{cm}.
\]
Montrer que $(MN) \parallel (BC)$.
\end{exemplebox}

\textbf{Résolution rédigée (modèle BEPC) :}

\begin{enumerate}
    \item \textbf{Alignement et ordre :} Par construction, $M$ appartient au segment $[AB]$ et $N$ appartient au segment $[AC]$. Donc $A$, $M$, $B$ sont alignés dans cet ordre, et $A$, $N$, $C$ sont alignés dans cet ordre.
    \item \textbf{Calcul des rapports :}
    \[
    \frac{AM}{AB} = \frac{2}{5} = 0,4
    \]
    \[
    \frac{AN}{AC} = \frac{3}{7,5} = \frac{30}{75} = \frac{2}{5} = 0,4
    \]
    \item \textbf{Comparaison :} On trouve $\frac{AM}{AB} = \frac{AN}{AC} = 0,4$ (ou $\frac{2}{5}$).
    \item \textbf{Conclusion :} D'après la \textbf{réciproque du théorème de Thalès}, les droites $(MN)$ et $(BC)$ sont \textbf{parallèles}.
\end{enumerate}

\begin{popfigure}
\begin{tikzpicture}[scale=1.5, >=Stealth]
% Triangle ABC
\coordinate (A) at (0,0);
\coordinate (B) at (5,0);
\coordinate (C) at (2, 6);
% M on AB at 2/5
\coordinate (M) at (2,0);
% N on AC at 3/7.5 = 0.4
\coordinate (N) at (0.8, 2.4);

\draw[popdark, thick] (A) -- (B) -- (C) -- cycle;
\draw[popblue, thick] (M) -- (N);

% Points
\fill[popdark] (A) circle (1.5pt) node[below left] {$A$};
\fill[popdark] (B) circle (1.5pt) node[below right] {$B$};
\fill[popdark] (C) circle (1.5pt) node[above] {$C$};
\fill[popblue] (M) circle (1.5pt) node[below] {$M$};
\fill[popblue] (N) circle (1.5pt) node[left] {$N$};

% Parallel marks (double arrows)
\draw[popgreen, double distance=1.5pt, -{Stealth[length=2mm]}] ($(B)!0.25!(C)$) -- ($(B)!0.45!(C)$);
\draw[popgreen, double distance=1.5pt, -{Stealth[length=2mm]}] ($(B)!0.55!(C)$) -- ($(B)!0.75!(C)$);
\draw[popgreen, double distance=1.5pt, -{Stealth[length=2mm]}] ($(M)!0.25!(N)$) -- ($(M)!0.45!(N)$);
\draw[popgreen, double distance=1.5pt, -{Stealth[length=2mm]}] ($(M)!0.55!(N)$) -- ($(M)!0.75!(N)$);

% Length labels
\node[popmuted, font=\small] at (1, -0.35) {$AM=\SI{2}{cm}$};
\node[popmuted, font=\small] at (3.5, -0.35) {$MB=\SI{3}{cm}$};
\node[popmuted, font=\small, rotate=71.5] at (0.4, 1.2) {$AN=\SI{3}{cm}$};
\node[popmuted, font=\small, rotate=71.5] at (1.4, 4.2) {$NC=\SI{4,5}{cm}$};

% Ratio boxes
\node[draw=popblue, fill=popblueL, rounded corners, font=\small, align=center] at (4.5, 2) {$\dfrac{AM}{AB} = \dfrac{2}{5}$\\[2pt] $\dfrac{AN}{AC} = \dfrac{3}{7,5} = \dfrac{2}{5}$};
\node[draw=popgreen, fill=popgreenL, rounded corners, font=\small\bfseries, align=center] at (4.5, 0.8) {$\dfrac{AM}{AB} = \dfrac{AN}{AC}$\\[2pt] $\Rightarrow (MN) \parallel (BC)$};
\end{tikzpicture}
\legende{Configuration de Thalès : les rapports sont égaux ($\frac{2}{5}$), les points sont dans l'ordre sur les segments, donc $(MN) \parallel (BC)$.}
\end{popfigure}

\courssub{Autre exemple : fractions simplifiées}

\begin{exoresolu}[Parallélisme et fractions irréductibles]
Dans le triangle $RST$, les points $U$ et $V$ sont placés sur les segments $[RS]$ et $[RT]$ respectivement. On donne : $RU = \SI{4}{cm}$, $RS = \SI{6}{cm}$, $RV = \SI{6}{cm}$, $RT = \SI{9}{cm}$.
Montrer que $(UV) \parallel (ST)$.
\tcblower
\textbf{Solution :}
\begin{enumerate}
    \item $U \in [RS]$ et $V \in [RT]$, donc $R, U, S$ alignés dans cet ordre et $R, V, T$ alignés dans cet ordre.
    \item $\dfrac{RU}{RS} = \dfrac{4}{6} = \dfrac{2}{3}$.
    \item $\dfrac{RV}{RT} = \dfrac{6}{9} = \dfrac{2}{3}$.
    \item On a $\dfrac{RU}{RS} = \dfrac{RV}{RT} = \dfrac{2}{3}$.
    \item Par la réciproque de Thalès, $(UV) \parallel (ST)$.
\end{enumerate}
\end{exoresolu}

\courssub{Erreurs classiques à ne plus commettre}

\begin{attention}[Pièges à éviter absolument]
\begin{itemize}
    \item \textbf{Oublier la phrase d'alignement/ordre.} Écrire seulement « $\frac{AM}{AB} = \frac{AN}{AC} $, donc $(MN) \parallel (BC)$ » fait perdre \textbf{la moitié des points} (ou tous les points) au BEPC. La condition d'ordre est une \emph{hypothèse} du théorème, elle doit être citée.
    \item \textbf{Confondre directe et réciproque.}
    \begin{itemize}
        \item \emph{Directe :} « $(MN) \parallel (BC)$ \textbf{donc} $\frac{AM}{AB} = \frac{AN}{AC}$ » (on \textbf{cherche une longueur}).
        \item \emph{Réciproque :} « $\frac{AM}{AB} = \frac{AN}{AC}$ \textbf{donc} $(MN) \parallel (BC)$ » (on \textbf{prouve un parallélisme}).
    \end{itemize}
    \item \textbf{Utiliser les mauvais segments.} Le rapport est toujours $\frac{\text{côté du petit triangle}}{\text{côté correspondant du grand triangle}}$ (ex: $\frac{AM}{AB}$, pas $\frac{AM}{MB}$). Attention : $\frac{AM}{MB} = \frac{AN}{NC}$ est \textbf{aussi vrai} si $(MN) \parallel (BC)$, mais ce n'est \textbf{pas} l'énoncé standard du théorème de Thalès au programme de 3ème. Reste sur $\frac{AM}{AB} = \frac{AN}{AC} = \frac{MN}{BC}$.
    \item \textbf{Calculer un seul rapport.} Il faut \emph{les deux} calculs distincts, puis la comparaison. Dire « $2/5 = 3/7,5$ » sans montrer $2/5=0,4$ et $3/7,5=0,4$ est insuffisant.
\end{itemize}
\end{attention}

\courssub{Exercices d'application immédiate}

\begin{exercice}[Entraînement : Vérifier le parallélisme]
Dans chaque figure ci-dessous, le triangle est nommé $ABC$, avec $M \in [AB]$ et $N \in [AC]$. Les longueurs sont en centimètres. Déterminer, en justifiant soigneusement, si $(MN) \parallel (BC)$.

\begin{enumerate}
    \item $AM = 3,\; AB = 7,\; AN = 4,5,\; AC = 10,5$.
    \item $AM = 5,\; AB = 8,\; AN = 6,\; AC = 9$.
    \item $AM = 2,4,\; AB = 6,\; AN = 3,2,\; AC = 8$.
    \item $AM = 4,\; MB = 6,\; AN = 5,\; NC = 7,5$. \emph{(Attention : on donne $MB$ et $NC$, pas $AB$ et $AC$ directement !)}
\end{enumerate}
\end{exercice}

\begin{corrige}[Exercice : Vérifier le parallélisme]
\begin{enumerate}
    \item \textbf{Alignement :} $M \in [AB], N \in [AC]$. \\
    $\frac{AM}{AB} = \frac{3}{7}$. $\frac{AN}{AC} = \frac{4,5}{10,5} = \frac{45}{105} = \frac{3}{7}$. \\
    Égalité vérifiée $\Rightarrow$ \textbf{$(MN) \parallel (BC)$}.
    \item \textbf{Alignement :} $M \in [AB], N \in [AC]$. \\
    $\frac{AM}{AB} = \frac{5}{8} = 0,625$. $\frac{AN}{AC} = \frac{6}{9} = \frac{2}{3} \approx 0,666$. \\
    Rapports différents $\Rightarrow$ \textbf{$(MN)$ n'est pas parallèle à $(BC)$}.
    \item \textbf{Alignement :} $M \in [AB], N \in [AC]$. \\
    $\frac{AM}{AB} = \frac{2,4}{6} = \frac{24}{60} = \frac{2}{5} = 0,4$. $\frac{AN}{AC} = \frac{3,2}{8} = \frac{32}{80} = \frac{2}{5} = 0,4$. \\
    Égalité vérifiée $\Rightarrow$ \textbf{$(MN) \parallel (BC)$}.
    \item \textbf{Piège !} On a $MB=6$ et $NC=7,5$. Il faut d'abord retrouver $AB$ et $AC$. \\
    $AB = AM + MB = 4 + 6 = 10$. $AC = AN + NC = 5 + 7,5 = 12,5$. \\
    $\frac{AM}{AB} = \frac{4}{10} = \frac{2}{5} = 0,4$. $\frac{AN}{AC} = \frac{5}{12,5} = \frac{50}{125} = \frac{2}{5} = 0,4$. \\
    Égalité vérifiée $\Rightarrow$ \textbf{$(MN) \parallel (BC)$}.
\end{enumerate}
\end{corrige}

\courssub{Synthèse : Directe vs Réciproque}

\begin{aretenir}[Résumé visuel : Les deux visages de Thalès]
\begin{center}
\begin{tabularx}{\linewidth}{|>{\bfseries}l|X|X|}
\hline
\rowcolor{popblueL}
& \textbf{Propriété DIRECTE} & \textbf{Propriété RÉCIPROQUE} \\
\hline
\textbf{Point de départ (Hypothèses)} & $(MN) \parallel (BC)$ \newline $A,M,B$ et $A,N,C$ alignés même ordre & $\frac{AM}{AB} = \frac{AN}{AC}$ \newline $A,M,B$ et $A,N,C$ alignés même ordre \\
\hline
\textbf{Conclusion (Thèse)} & $\frac{AM}{AB} = \frac{AN}{AC} = \frac{MN}{BC}$ & $(MN) \parallel (BC)$ \\
\hline
\textbf{Question type} & « Calculer $MN$ », « Calculer $AC$ » & « Prouver que $(MN) \parallel (BC)$ », « Vérifier si... » \\
\hline
\textbf{Mots-clés rédaction} & « D'après Thalès \textbf{dans le triangle} $ABC$... » & « D'après la \textbf{réciproque} de Thalès... » \\
\hline
\end{tabularx}
\end{center}
\end{aretenir}

\begin{savaistu}[Thalès au champ : l'alignement des cultures]
Au Cameroun, comme partout dans les zones agricoles (Dschang, Ngaoundéré, Maroua), les agriculteurs utilisent intuitivement la réciproque de Thalès pour planter en lignes parallèles. Pour tracer des rangées de maïs ou de haricots parfaitement parallèles sans GPS, on plante deux piquets $A$ et $B$ pour la première rangée. On choisit un point $M$ sur la ligne $AB$. On mesure $AM$ et $AB$. On reporte le même rapport sur une ligne de départ $AC$ pour trouver $N$ tel que $\frac{AN}{AC} = \frac{AM}{AB}$. La ligne $MN$ tracée à la corde sera \emph{mathématiquement} parallèle à $AB$. C'est de la géométrie vivante, sans théodolite ni laser !
\end{savaistu}

\coursec{Contraposée : prouver que des droites ne sont pas parallèles}

Dans la section précédente, nous avons vu comment la \emph{réciproque} de Thalès permet de \textbf{prouver que deux droites sont parallèles} en vérifiant l'égalité des rapports. Mais que se passe-t-il si les rapports \textbf{ne sont pas égaux} ? C'est précisément la question à laquelle répond la \textbf{contraposée}. C'est un outil logique puissant, souvent oublié, mais indispensable pour conclure \emph{négativement} : « ces droites ne sont pas parallèles ».

\courssub{Rappel logique : la contraposée}

Rappelons un principe fondamental de la logique mathématique, valable pour \textbf{toute} implication vraie.

\begin{definition}[Contraposée d'une implication]
Soit une implication de la forme \og Si \textbf{P}, alors \textbf{Q} \fg (notée $P \Rightarrow Q$).
Sa \textbf{contraposée} est l'implication : \og Si \textbf{non Q}, alors \textbf{non P} \fg (notée $\lnot Q \Rightarrow \lnot P$).
\end{definition}

\begin{propriete}[Équivalence logique]
Une implication et sa contraposée sont \textbf{logiquement équivalentes} : elles sont \emph{simultanément vraies ou simultanément fausses}.
\[
(P \Rightarrow Q) \quad \Longleftrightarrow \quad (\lnot Q \Rightarrow \lnot P)
\]
Si une propriété (théorème) est vraie, sa contraposée l'est \textbf{automatiquement}, sans nouvelle démonstration.
\end{propriete}

\emph{Intuition :} Si \og S'il pleut, le sol est mouillé \fg est vrai, alors \og Si le sol n'est \textbf{pas} mouillé, il n'a \textbf{pas} plu \fg est aussi vrai. En revanche, l'inverse (\og Si le sol est mouillé, il a plu \fg) et la réciproque peuvent être faux (arrosage, fuite de tuyau).

Appliquons cela à la propriété directe de Thalès.

\courssub{Énoncé de la contraposée de Thalès}

La propriété directe s'écrit :
\[
\textbf{P :} \left\{ \begin{array}{l}
A, M, B \text{ alignés dans cet ordre} \\
A, N, C \text{ alignés dans cet ordre} \\
(MN) \parallel (BC)
\end{array} \right.
\quad \Rightarrow \quad
\textbf{Q :} \quad \frac{AM}{AB} = \frac{AN}{AC} = \frac{MN}{BC}
\]

Sa contraposée (\og Si non Q, alors non P \fg) donne :

\begin{propriete}[Contraposée de la propriété de Thalès]
Soit un triangle $ABC$. Soient $M \in (AB)$ et $N \in (AC)$, les points étant dans le \textbf{même ordre} ($A-M-B$ et $A-N-C$).
Si les rapports ne sont \textbf{pas égaux} :
\[
\frac{AM}{AB} \neq \frac{AN}{AC}
\]
alors les droites $(MN)$ et $(BC)$ ne sont \textbf{pas parallèles} (elles sont sécantes).
\end{propriete}

\begin{attention}[Conditions d'application indispensables]
Pour utiliser la contraposée, il faut \textbf{impérativement} vérifier \textbf{avant tout calcul} :
\begin{enumerate}
    \item $M$ appartient bien à la droite $(AB)$ et $N$ à la droite $(AC)$ (alignement).
    \item Les points sont dans le \textbf{même ordre} : $A-M-B$ et $A-N-C$ (ou $M-A-B$ et $N-A-C$, etc.). L'ordre détermine le sens des segments et la validité des rapports.
\end{enumerate}
Si l'alignement ou l'ordre n'est pas respecté, la contraposée (comme la directe et la réciproque) \textbf{ne s'applique pas}. On ne peut rien conclure sur le parallélisme avec Thalès.
\end{attention}

\courssub{Distinction claire entre les trois usages}

C'est le cœur de la maîtrise de Thalès en 3ème : savoir \textbf{quelle propriété invoquer} selon la question posée.

\begin{methode}[Choisir la bonne propriété de Thalès]
\begin{enumerate}
    \item \textbf{On cherche une longueur inconnue} (ou on vérifie une égalité de longueurs) \\
    $\rightarrow$ \textbf{Propriété DIRECTE} \\ 
    Hypothèses : Parallélisme connu + Alignements + Ordre. \\
    Conclusion : Égalité des rapports $\Rightarrow$ Calcul de longueur.
    
    \item \textbf{On veut PROUVER que deux droites SONT parallèles} \\
    $\rightarrow$ \textbf{Propriété RÉCIPROQUE} \\
    Hypothèses : Alignements + Ordre + Égalité des rapports (vérifiée par le calcul). \\
    Conclusion : Parallélisme.
    
    \item \textbf{On veut PROUVER que deux droites NE SONT PAS parallèles} \\
    $\rightarrow$ \textbf{Propriété CONTRAPOSÉE} \\
    Hypothèses : Alignements + Ordre + \textbf{Inégalité des rapports} (vérifiée par le calcul). \\
    Conclusion : Non-parallélisme (droites sécantes).
\end{enumerate}
\end{methode}

Voici un tableau récapitulatif visuel pour ne plus jamais confondre :

\begin{popfigure}
\centering
\begin{tikzpicture}[
    cell/.style={rectangle, draw=popdark, minimum height=1.1cm, align=center, text width=3.2cm, font=\small},
    header/.style={cell, fill=popblue, text=white, font=\small\bfseries},
    title/.style={cell, fill=popdark, text=white, font=\normalsize\bfseries, text width=10.6cm},
    row1/.style={fill=popblueL},
    row2/.style={fill=poporangeL},
    row3/.style={fill=popgreenL},
]
% Title
\node[title] (t) at (0,0) {TABLEAU RÉCAPITULATIF : LES TROIS PROPRIÉTÉS DE THALÈS};
% Headers
\node[header, below=0pt of t.south west, anchor=north west] (h1) {Propriété};
\node[header, right=0pt of h1] (h2) {Hypothèses (Ce qu'on sait / vérifie)};
\node[header, right=0pt of h2] (h3) {Conclusion (Ce qu'on obtient)};
\node[header, right=0pt of h3] (h4) {Usage principal};
% Row 1 : Directe
\node[cell, row1, below=0pt of h1.south west, anchor=north west] (r1c1) {\textbf{Directe}};
\node[cell, row1, right=0pt of r1c1] (r1c2) {
    \begin{itemize}\itemsep0pt \parskip0pt
    \item $(MN) \parallel (BC)$
    \item $A,M,B$ alignés (ordre)
    \item $A,N,C$ alignés (ordre)
    \end{itemize}
};
\node[cell, row1, right=0pt of r1c2] (r1c3) {$\dfrac{AM}{AB} = \dfrac{AN}{AC} = \dfrac{MN}{BC}$};
\node[cell, row1, right=0pt of r1c3] (r1c4) {Calculer une longueur};
% Row 2 : Réciproque
\node[cell, row2, below=0pt of r1c1.south west, anchor=north west] (r2c1) {\textbf{Réciproque}};
\node[cell, row2, right=0pt of r2c1] (r2c2) {
    \begin{itemize}\itemsep0pt \parskip0pt
    \item $A,M,B$ alignés (ordre)
    \item $A,N,C$ alignés (ordre)
    \item $\dfrac{AM}{AB} = \dfrac{AN}{AC}$
    \end{itemize}
};
\node[cell, row2, right=0pt of r2c2] (r2c3) {$(MN) \parallel (BC)$};
\node[cell, row2, right=0pt of r2c3] (r2c4) {Prouver le parallélisme};
% Row 3 : Contraposée
\node[cell, row3, below=0pt of r2c1.south west, anchor=north west] (r3c1) {\textbf{Contraposée}};
\node[cell, row3, right=0pt of r3c1] (r3c2) {
    \begin{itemize}\itemsep0pt \parskip0pt
    \item $A,M,B$ alignés (ordre)
    \item $A,N,C$ alignés (ordre)
    \item $\dfrac{AM}{AB} \neq \dfrac{AN}{AC}$
    \end{itemize}
};
\node[cell, row3, right=0pt of r3c2] (r3c3) {$(MN)$ \textbf{n'est pas} $\parallel$ à $(BC)$};
\node[cell, row3, right=0pt of r3c3] (r3c4) {Prouver le non-parallélisme};
\end{tikzpicture}
\legende{Les trois visages de Thalès : un même triangle, trois questions différentes.}
\end{popfigure}

\courssub{Exemple chiffré détaillé}

Appliquons la méthode sur un cas concret.

\begin{exoresolu}[Utiliser la contraposée pour prouver le non-parallélisme]
Dans le triangle $ABC$, on place $M \in [AB]$ et $N \in [AC]$ tels que :
$AM = 2$ cm, $AB = 5$ cm, $AN = 3$ cm, $AC = 8$ cm.
Montrer que les droites $(MN)$ et $(BC)$ ne sont pas parallèles.
\tcblower
\textbf{Solution rédigée :}

\textbf{1. Vérification des hypothèses (alignement et ordre).}
D'après l'énoncé, $M \in [AB]$ et $N \in [AC]$.
Les points sont donc alignés dans l'ordre $A-M-B$ et $A-N-C$. Les hypothèses de configuration sont satisfaites.

\textbf{2. Calcul des rapports.}
\[
\frac{AM}{AB} = \frac{2}{5} = 0,4
\]
\[
\frac{AN}{AC} = \frac{3}{8} = 0,375
\]

\textbf{3. Comparaison.}
On observe que $0,4 \neq 0,375$, donc :
\[
\frac{AM}{AB} \neq \frac{AN}{AC}
\]

\textbf{4. Conclusion par la contraposée.}
Puisque les rapports sont inégaux, d'après la \textbf{contraposée de la propriété de Thalès}, les droites $(MN)$ et $(BC)$ \textbf{ne sont pas parallèles}. Elles sont sécantes.
\end{exoresolu}

\begin{exemplebox}[Autre exemple numérique]
Soit un triangle $ABC$ avec $M \in (AB)$, $N \in (AC)$ (ordre $A-M-B$, $A-N-C$).
Données : $AM = 5$, $AB = 12$, $AN = 7$, $AC = 17$.
Calculons :
$\dfrac{AM}{AB} = \dfrac{5}{12} \approx 0,4167$
$\dfrac{AN}{AC} = \dfrac{7}{17} \approx 0,4118$
Les rapports sont \textbf{différents} (même s'ils sont proches !).
$\Rightarrow$ $(MN)$ et $(BC)$ ne sont \textbf{pas parallèles}.
\end{exemplebox}

\courssub{Applications pratiques : vérifier la réalité}

La contraposée n'est pas qu'un exercice scolaire : c'est l'outil du \textbf{contrôle qualité} et de la \textbf{vérification sur le terrain}.

\begin{savaistu}[Le tailleur et le tissu pagne]
Au marché central de Yaoundé ou de Douala, un tailleur coupe un grand rectangle de tissu pagne pour faire une chemise. Il doit vérifier que les deux grands bords (les lisières) sont bien parallèles avant de couper, sinon la chemise sera de travers.
Il ne peut pas mesurer l'angle (pas de grand rapporteur). Il utilise la contraposée de Thalès :
\begin{enumerate}
    \item Il choisit un coin $A$ du tissu.
    \item Sur un bord, il marque un point $M$ à $50$ cm. Sur l'autre bord, un point $N$ à $80$ cm.
    \item Il mesure la distance $MN$ (diagonale) et la compare à ce qu'elle devrait être si c'était parallèle (calcul via Thalès direct ou mesure de l'autre diagonale).
    \item \textbf{Si les rapports ne collent pas} ($AM/AB \neq AN/AC$), la contraposée lui dit formellement : \og \textbf{Attention !} Ces bords ne sont pas parallèles, le tissu est déformé. \fg
\end{enumerate}
C'est une application directe de \og non Q $\Rightarrow$ non P \fg : rapports inégaux $\Rightarrow$ pas de parallélisme.
\end{savaistu}

\begin{savaistu}[Électricien / Maçon : verticalité d'un mur]
Pour vérifier qu'un mur est vertical (parallèle à la verticale), on peut utiliser un fil à plomb (verticale de référence) et mesurer des écartements à deux hauteurs différentes ($A$ en bas, $B$ en haut). Si les rapports des écartements horizontaux ne sont pas égaux, le mur n'est pas vertical (il penche). La contraposée valide le diagnostic.
\end{savaistu}

\courssub{Illustration géométrique : rapports inégaux $\Rightarrow$ droites sécantes}

La figure ci-dessous montre concrètement ce qui se passe quand les rapports diffèrent : la droite $(MN)$ « tourne » par rapport à $(BC)$ et la coupe nécessairement.

\begin{popfigure}
\centering
\begin{tikzpicture}[scale=1.2, >=Stealth]
% Points du grand triangle
\coordinate (A) at (0,0);
\coordinate (B) at (6,0);
\coordinate (C) at (1.5,4.5);

% Points M, N sur les côtés (rapports inégaux)
% AM/AB = 2/5 = 0.4  -> M à 40% de AB
% AN/AC = 3/8 = 0.375 -> N à 37.5% de AC
\coordinate (M) at ($(A)!0.4!(B)$); % 2.4, 0
\coordinate (N) at ($(A)!0.375!(C)$); % 0.5625, 1.6875

% Intersection de (MN) et (BC) -> point I
\coordinate (I) at (intersection of M--N and B--C);

% Dessin triangle principal
\draw[popdark, thick] (A) -- (B) -- (C) -- cycle;
\node[below left] at (A) {$A$};
\node[below right] at (B) {$B$};
\node[above left] at (C) {$C$};

% Points M, N
\fill[popblue] (M) circle (2pt) node[below] {$M$};
\fill[popblue] (N) circle (2pt) node[left] {$N$};

% Droite MN (prolongée pour voir intersection)
\draw[poporange, thick, dashed] ($(M)!-1.5!(N)$) -- ($(N)!-0.5!(M)$);
% Droite BC
\draw[popgreen, thick] (B) -- (C);

% Intersection I
\fill[poppink] (I) circle (2.5pt) node[above right] {$I$};

% Annotations des rapports
% Côté AB
\draw[popmuted, <->] ($(A)!-0.3!(B)$) -- ($(M)!-0.3!(B)$) node[midway, left=2mm, popblue] {$AM=2$};
\draw[popmuted, <->] ($(M)!-0.3!(B)$) -- ($(B)!-0.3!(A)$) node[midway, left=2mm, popmuted] {$MB=3$};
\draw[popmuted, <->] ($(A)!-0.6!(B)$) -- ($(B)!-0.6!(A)$) node[midway, left=4mm, popdark] {$AB=5$};

% Côté AC
\draw[popmuted, <->] ($(A)!-0.3!(C)$) -- ($(N)!-0.3!(C)$) node[midway, left=2mm, popblue] {$AN=3$};
\draw[popmuted, <->] ($(N)!-0.3!(C)$) -- ($(C)!-0.3!(A)$) node[midway, left=2mm, popmuted] {$NC=5$};
\draw[popmuted, <->] ($(A)!-0.6!(C)$) -- ($(C)!-0.6!(A)$) node[midway, left=4mm, popdark] {$AC=8$};

% Boîte de conclusion
\node[draw=poppink, fill=poppink!10, rounded corners, align=center, font=\small\bfseries, anchor=north] at (3.5, -1.2) {
    $\dfrac{AM}{AB} = \dfrac{2}{5} = \num{0,4}$ \quad $\neq$ \quad $\dfrac{AN}{AC} = \dfrac{3}{8} = \num{0,375}$ \\
    $\Rightarrow$ $(MN)$ et $(BC)$ sont \textcolor{poppink}{\textbf{sécantes}} en $I$.
};
\end{tikzpicture}
\legende{Configuration de la contraposée : $AM/AB \neq AN/AC \Rightarrow (MN)$ coupe $(BC)$ en $I$.}
\end{popfigure}

\courssub{Synthèse et pièges à éviter avant les exercices}

\begin{aretenir}[Résumé : Les 3 cartes à jouer avec Thalès]
\begin{center}
\begin{tabular}{|c|c|c|c|}
\hline
\textbf{Question posée} & \textbf{Propriété} & \textbf{On vérifie...} & \textbf{On conclut...} \\
\hline
Calculer $MN$, $AB$, $AC$... & \textbf{Directe} & Parallélisme + Alignements & Égalité des rapports $\rightarrow$ Longueur \\
\hline
Prouver $(MN) \parallel (BC)$ & \textbf{Réciproque} & Alignements + $\frac{AM}{AB}=\frac{AN}{AC}$ & Parallélisme \\
\hline
Prouver $(MN)$ \textbf{non} $\parallel$ $(BC)$ & \textbf{Contraposée} & Alignements + $\frac{AM}{AB}\neq\frac{AN}{AC}$ & Non-parallélisme (sécantes) \\
\hline
\end{tabular}
\end{center}
\end{aretenir}

\begin{attention}[Les 3 erreurs fatales à ne plus commettre]
\begin{enumerate}
    \item \textbf{Oublier de vérifier l'alignement et l'ordre.} Écrire « D'après la contraposée de Thalès... » sans avoir écrit « $M \in (AB)$, $N \in (AC)$ dans l'ordre $A-M-B$ et $A-N-C$ » fait perdre \textbf{tous les points} de la justification au BEPC.
    \item \textbf{Confondre Réciproque et Contraposée.}
    \begin{itemize}
        \item Rapports \textbf{égaux} $\rightarrow$ Réciproque $\rightarrow$ \textbf{SONT} parallèles.
        \item Rapports \textbf{inégaux} $\rightarrow$ Contraposée $\rightarrow$ \textbf{NE SONT PAS} parallèles.
    \end{itemize}
    \item \textbf{Arrondir trop vite.} $\frac{5}{12} \approx 0,417$ et $\frac{7}{17} \approx 0,412$. Si tu arrondis à $0,42$ pour les deux, tu crois à tort qu'ils sont égaux. \textbf{Garde les fractions} ou calcule en produit en croix ($5 \times 17 = 85$ vs $7 \times 12 = 84$) pour être exact.
\end{enumerate}
\end{attention}

\begin{exoresolu}[Examen type BEPC : Argumenter avec la contraposée]
Dans la figure ci-contre (non fournie, se fier aux données), $A, M, B$ sont alignés dans cet ordre et $A, N, C$ sont alignés dans cet ordre.
On donne : $AM = 4,5$ cm, $MB = 3,5$ cm, $AN = 6$ cm, $NC = 4$ cm.
Les droites $(MN)$ et $(BC)$ sont-elles parallèles ? Justifier.
\tcblower
\textbf{Rédaction attendue :}

\begin{enumerate}
    \item \textbf{Calcul des longueurs totales :}
    $AB = AM + MB = 4,5 + 3,5 = 8$ cm.
    $AC = AN + NC = 6 + 4 = 10$ cm.
    
    \item \textbf{Calcul des rapports :}
    $\dfrac{AM}{AB} = \dfrac{4,5}{8} = \dfrac{45}{80} = \dfrac{9}{16}$.
    $\dfrac{AN}{AC} = \dfrac{6}{10} = \dfrac{3}{5} = \dfrac{9,6}{16}$.
    
    \item \textbf{Comparaison (produit en croix plus sûr) :}
    $AM \times AC = 4,5 \times 10 = 45$.
    $AN \times AB = 6 \times 8 = 48$.
    Comme $45 \neq 48$, on a $\dfrac{AM}{AB} \neq \dfrac{AN}{AC}$.
    
    \item \textbf{Conclusion :}
    Les points $M$ et $N$ sont sur les côtés $[AB]$ et $[AC]$ (alignements et ordre vérifiés).
    Comme les rapports ne sont pas égaux, d'après la \textbf{contraposée du théorème de Thalès}, les droites $(MN)$ et $(BC)$ \textbf{ne sont pas parallèles}.
\end{enumerate}
\end{exoresolu}

\courssub{Exercices d'entraînement immédiat}

\begin{exercice}[Choix de la propriété]
Pour chaque situation, indique quelle propriété de Thalès utiliser (Directe, Réciproque ou Contraposée) et écris la conclusion attendue.
\begin{enumerate}
    \item Dans $\triangle ABC$, $M \in [AB], N \in [AC]$, $(MN) \parallel (BC)$. On connaît $AM, AB, AC$. On cherche $AN$.
    \item Dans $\triangle DEF$, $P \in [DE], Q \in [DF]$. On vérifie que $\frac{DP}{DE} = \frac{DQ}{DF}$. On veut savoir si $(PQ) \parallel (EF)$.
    \item Dans $\triangle GHI$, $R \in [GH], S \in [GI]$. On calcule $\frac{GR}{GH} = 0,6$ et $\frac{GS}{GI} = 0,55$. On veut savoir si $(RS) \parallel (HI)$.
    \item Un menuisier vérifie l'équerrage d'une porte rectangulaire. Il mesure les diagonales $AC$ et $BD$ qui se coupent en $O$. Il compare $AO/AC$ et $BO/BD$. Quelle propriété utilise-t-il implicitement ?
\end{enumerate}
\end{exercice}

\begin{exercice}[Preuve de non-parallélisme]
Soit un triangle $JKL$. Les points $M$ et $N$ sont sur les demi-droites $[JK)$ et $[JL)$ (ordre $J-M-K$ et $J-N-L$).
On a : $JM = 12$, $JK = 20$, $JN = 15$, $JL = 24$.
Montrer que $(MN)$ et $(KL)$ ne sont pas parallèles.
\end{exercice}

\begin{exercice}[Cas limite : ordre des points]
Dans un triangle $OPQ$, $R \in (OP)$ et $S \in (OQ)$.
On a $OR = 6$, $OP = 10$, $OS = 9$, $OQ = 15$.
\begin{enumerate}
    \item Vérifier si $\frac{OR}{OP} = \frac{OS}{OQ}$.
    \item Si l'ordre est $O-R-P$ et $O-S-Q$, que conclure sur $(RS)$ et $(PQ)$ ?
    \item Si l'ordre est $R-O-P$ et $O-S-Q$ ($R$ de l'autre côté de $O$), peut-on conclure avec Thalès ? Pourquoi ?
\end{enumerate}
\end{exercice}

\begin{corrige}[Exercices d'entraînement immédiat]
\textbf{Exercice 1 :}
\begin{enumerate}
    \item \textbf{Directe.} Conclusion : $AN = \frac{AM \times AC}{AB}$.
    \item \textbf{Réciproque.} Conclusion : $(PQ) \parallel (EF)$.
    \item \textbf{Contraposée.} Conclusion : $(RS)$ n'est pas parallèle à $(HI)$ (elles sont sécantes).
    \item Dans un rectangle, les diagonales ont même milieu $O$, donc $AO=OC$ et $BO=OD$. Les rapports $AO/AC = 1/2$ et $BO/BD = 1/2$ sont égaux. C'est la \textbf{Réciproque} de Thalès appliquée dans les triangles $AOB$ et $COD$ (ou $AOD$ et $BOC$) qui garantit que les côtés opposés sont parallèles (parallélogramme) et l'égalité des diagonales donne le rectangle. Ici, le menuisier vérifie l'égalité des diagonales (propriété du rectangle), ce qui revient à vérifier l'égalité des rapports des demi-diagonales (Réciproque de Thalès).
\end{enumerate}

\textbf{Exercice 2 :}
$JK = 20 \Rightarrow \frac{JM}{JK} = \frac{12}{20} = \frac{3}{5} = 0,6$.
$JL = 24 \Rightarrow \frac{JN}{JL} = \frac{15}{24} = \frac{5}{8} = 0,625$.
$0,6 \neq 0,625$.
Alignements et ordre vérifiés ($J-M-K$, $J-N-L$).
Par la contraposée, $(MN)$ et $(KL)$ ne sont pas parallèles.

\textbf{Exercice 3 :}
\begin{enumerate}
    \item $\frac{OR}{OP} = \frac{6}{10} = 0,6$. $\frac{OS}{OQ} = \frac{9}{15} = 0,6$. Les rapports \textbf{sont égaux}.
    \item Ordre $O-R-P$, $O-S-Q$ + rapports égaux $\Rightarrow$ \textbf{Réciproque} : $(RS) \parallel (PQ)$.
    \item Ordre $R-O-P$ (donc $O$ est entre $R$ et $P$) et $O-S-Q$. Les points ne sont \textbf{pas dans le même ordre} sur les deux droites (sur $OP$, $O$ sépare $R$ et $P$ ; sur $OQ$, $O$ est à l'extrémité).
    \textbf{Thalès ne s'applique pas} (configuration « en X » ou papillon, pas configuration « en V » ou triangle). On ne peut rien conclure avec Thalès.
\end{enumerate}
\end{corrige}

\emph{Tu as maintenant toutes les cartes en main : Directe pour calculer, Réciproque pour valider le parallélisme, Contraposée pour l'invalider. La semaine prochaine, nous synthétiserons tout le chapitre et verrons des problèmes complexes mélangeant ces trois outils.}

\coursec{Synthèse et prolongements}

\courssub{Bilan des deux propriétés de Thalès}

Après avoir étudié la propriété directe, la réciproque et la contraposée, il est temps de faire le point. Rappelle‑toi que \textbf{toute utilisation du théorème de Thalès repose sur une configuration précise} : deux droites sécantes coupées par au moins deux droites parallèles.

\begin{definition}[Configuration de Thalès]
On appelle \textbf{configuration de Thalès} la figure formée par deux droites $(AB)$ et $(AC)$ sécantes en $A$, deux points $M \in (AB)$ et $N \in (AC)$ distincts de $A$, tels que les droites $(MN)$ et $(BC)$ sont parallèles. Les points $A, M, B$ sont alignés (de même $A, N, C$) ; ils peuvent être dans l'ordre $A-M-B$ ou $A-B-M$ (de même pour l'autre sécante). On dit que les triangles $AMN$ et $ABC$ sont en \emph{configuration de Thalès}.
\end{definition}

\begin{propriete}[Propriété directe]
Si $(MN)\parallel(BC)$, alors les rapports de longueurs sont égaux :
\[
\frac{AM}{AB}=\frac{AN}{AC}=\frac{MN}{BC}.
\]
Cette égalité permet de calculer une longueur inconnue dès qu'on connaît les trois autres.
\end{propriete}

\begin{propriete}[Propriété réciproque]
Si les points $A,M,B$ et $A,N,C$ sont alignés dans le \textbf{même ordre} sur chaque sécante et si
\[
\frac{AM}{AB}=\frac{AN}{AC},
\]
alors $(MN)\parallel(BC)$.
\end{propriete}

\begin{attention}[Conditions d'application souvent oubliées]
\begin{itemize}
\item \textbf{Directe :} il faut \emph{absolument} que $(MN)\parallel(BC)$ soit donné ou prouvé avant d'écrire les égalités de rapports.
\item \textbf{Réciproque :} il faut vérifier l'égalité des deux rapports \textbf{et} que les points sont dans le même ordre sur chaque sécante (pas de points « croisés »).
\item \textbf{Contraposée :} pour montrer que deux droites ne sont pas parallèles, on montre que les rapports sont \emph{différents}.
\end{itemize}
\end{attention}

\courssub{Fiche de rédaction modèle pour l'examen (BEPC)}

Une rédaction soignée vaut des points. Voici la structure type que tu dois reproduire à chaque exercice.

\begin{methode}[Rédaction type examen — Thalès]
\begin{enumerate}
\item \textbf{Citer les hypothèses} : nommer les droites, les points, l'alignement et la parallélité (ou l'égalité des rapports).
\item \textbf{Nommer la propriété} : « D'après la propriété directe (ou réciproque) de Thalès dans le triangle $ABC$… »
\item \textbf{Écrire l'égalité des rapports} : $\dfrac{AM}{AB}=\dfrac{AN}{AC}$ (ou avec $MN/BC$).
\item \textbf{Effectuer le calcul} : substituer les valeurs connues, simplifier, résoudre l'inconnue.
\item \textbf{Conclure avec unité} : « Donc $NC = 6\,\text{cm}$ » ou « Donc $(MN)\parallel(BC)$ ».
\end{enumerate}
\end{methode}

\courssub{Lien avec les agrandissements / réductions : le coefficient $k$}

La configuration de Thalès est exactement celle d'un \emph{agrandissement} (ou d'une \emph{réduction}) du triangle $AMN$ en $ABC$ (ou l'inverse). Le rapport
\[
k = \frac{AM}{AB} = \frac{AN}{AC} = \frac{MN}{BC}
\]
est le \textbf{coefficient de réduction} (si $k<1$) ou d'\textbf{agrandissement} (si $k>1$) qui transforme $ABC$ en $AMN$ (centre $A$).

\begin{itemize}
\item Si $0<k<1$, $AMN$ est une réduction de $ABC$ (le triangle initial est $ABC$).
\item Si $k>1$, $AMN$ est un agrandissement de $ABC$.
\item Si $k=1$, les deux triangles sont confondus.
\end{itemize}

C'est pourquoi, dans les exercices de « mise à l'échelle » (plans, cartes, maquettes), on utilise directement Thalès pour trouver les dimensions réelles à partir d'un dessin réduit.

\courssub{Lien avec les triangles semblables}

Deux triangles en configuration de Thalès ont \textbf{trois angles respectivement égaux} (angles correspondants formés par des parallèles) et \textbf{trois côtés proportionnels}. Ils sont donc \emph{semblables} au sens de la géométrie.

\begin{propriete}[Triangles semblables et Thalès]
Dans la configuration de Thalès, $\triangle AMN \sim \triangle ABC$. Par conséquent :
\[
\frac{AM}{AB}=\frac{AN}{AC}=\frac{MN}{BC},\qquad
\widehat{MAN}=\widehat{BAC},\;
\widehat{AMN}=\widehat{ABC},\;
\widehat{ANM}=\widehat{ACB}.
\]
\end{propriete}


\courssub{Erreurs classiques récapitulées et auto‑évaluation (quiz 5 questions)}

\begin{attention}[Pièges à éviter — checklist avant de rendre ta copie]
\begin{enumerate}
\item \textbf{Oublier la parallélité} pour la directe : on n'écrit pas les rapports si $(MN)\parallel(BC)$ n'est pas assuré.
\item \textbf{Inverser les rapports} : $\frac{AM}{AB}\neq\frac{AB}{AM}$ sauf si $k=1$.
\item \textbf{Mélanger l'ordre des points} sur une sécante (ex. $A,B,M$ au lieu de $A,M,B$) : la réciproque devient fausse.
\item \textbf{Confondre $MN$ et $BC$} dans le calcul : le rapport $\frac{MN}{BC}$ n'est utile que si on cherche $MN$ ou $BC$.
\item \textbf{Négliger l'unité} dans la réponse finale (cm, m, km…).
\end{enumerate}
\end{attention}

\begin{exoresolu}[Quiz d'auto‑évaluation — 5 questions rapides]
\textbf{Énoncé :} Dans le triangle $ABC$, les points $M\in[AB]$ et $N\in[AC]$ sont tels que $(MN)\parallel(BC)$. On donne $AM=6\,\text{cm}$, $MB=4\,\text{cm}$, $AN=9\,\text{cm}$. 
\begin{enumerate}
\item Calcule $AB$ et $AC$.
\item Détermine le coefficient $k=\frac{AM}{AB}$.
\item Trouve $NC$.
\item Vérifie que $\frac{MN}{BC}=k$ si $BC=12\,\text{cm}$.
\item Si on inverse les rôles (on connaît $NC=6$), comment retrouver $AN$ ?
\end{enumerate}

\tcblower

\textbf{Solution détaillée :}
\begin{enumerate}
\item $AB = AM+MB = 6+4 = 10\,\text{cm}$. \\
      Comme $(MN)\parallel(BC)$, $\frac{AM}{AB}=\frac{AN}{AC}\Rightarrow \frac{6}{10}=\frac{9}{AC}\Rightarrow AC = \frac{9\times10}{6}=15\,\text{cm}$.
\item $k = \frac{AM}{AB} = \frac{6}{10}=0,6$ (réduction).
\item $NC = AC-AN = 15-9 = 6\,\text{cm}$.
\item $\frac{MN}{BC}=k=0,6 \Rightarrow MN = 0,6\times12 = 7,2\,\text{cm}$.
\item On utilise la réciproque : $\frac{AN}{AC}=\frac{AM}{AB}=0,6$. Sachant $NC=6$, $AC=AN+6$. L'équation $\frac{AN}{AN+6}=0,6$ donne $AN=9\,\text{cm}$ (même résultat).
\end{enumerate}
\end{exoresolu}

\begin{exercice}[Entraînement — à faire seul]
Dans un triangle $XYZ$, $P\in[XY]$, $Q\in[XZ]$ avec $(PQ)\parallel(YZ)$. On a $XP=8\,\text{cm}$, $PY=4\,\text{cm}$, $XQ=12\,\text{cm}$. 
\begin{enumerate}
\item Calcule $XY$, $XZ$ et $QZ$.
\item Quel est le coefficient d'agrandissement de $XPQ$ vers $XYZ$ ?
\item Si $YZ=18\,\text{cm}$, quelle est la longueur de $PQ$ ?
\end{enumerate}
\end{exercice}

\begin{corrige}[Exercice entraînement]
\begin{enumerate}
\item $XY=12\,\text{cm}$, $\frac{XP}{XY}=\frac{XQ}{XZ}\Rightarrow \frac{8}{12}=\frac{12}{XZ}\Rightarrow XZ=18\,\text{cm}$, $QZ=6\,\text{cm}$.
\item $k=\frac{XP}{XY}=\frac{8}{12}=\frac{2}{3}$ (réduction).
\item $PQ = k\times YZ = \frac{2}{3}\times18 = 12\,\text{cm}$.
\end{enumerate}
\end{corrige}

\courssub{Synthèse finale}

Tu disposes maintenant d'une \textbf{boîte à outils complète} :
\begin{itemize}
\item La propriété directe pour \textbf{calculer} une longueur.
\item La propriété réciproque pour \textbf{prouver} une parallélité.
\item La contraposée pour \textbf{prouver} que deux droites ne sont pas parallèles.
\item Le lien avec l'homothétie (coefficient $k$) et les triangles semblables.
\item Une méthode de rédaction rigoureuse, exigible au BEPC.
\end{itemize}

Entraîne‑toi sur des figures variées (triangle, trapèze, losange, problèmes concrets : ombre d'un poteau, plan de ville, maquette d'architecture). La maîtrise de Thalès est un socle pour la géométrie du lycée (trigonométrie, vecteurs, géométrie analytique). Bonne révision et bon courage pour l'examen !

\newpage

\coursec{Activité d'intégration}

\coursec{Thalès dans le triangle : Activité d'intégration}

\courssub{Objectifs de l'activité}
Cette activité d'intégration vise à mobiliser l'ensemble des compétences travaillées dans la leçon sur Thalès dans un contexte authentique de \cle{modélisation géométrique} et de \cle{résolution de problèmes concrets}. À l'issue de cette séance, vous devez être capable de :
\begin{itemize}
    \item Reconnaître une configuration de Thalès (en \og X \fg{} ou en \og \textasciicircum \fg{}) dans une situation réelle et la modéliser par un schéma codé rigoureux.
    \item Appliquer la \cle{propriété directe} pour calculer une longueur inaccessible (hauteur d'un mât) à partir de mesures accessibles (ombres).
    \item Appliquer la \cle{propriété réciproque} ou sa \cle{contraposée} pour décider du parallélisme de deux droites à partir de mesures de longueurs sur deux sécantes.
    \item Rédiger un \cle{rapport mathématique} complet : hypothèses, propriété invoquée (avec sa condition), calculs détaillés, conclusion précise et réponse à la question posée.
    \item Travailler en équipe : répartir les tâches (mesures, schéma, calculs, rédaction, présentation orale).
\end{itemize}

\courssub{Situation : Le géomètre du quartier}
\noindent
Le Collège Bilingue de \textbf{Nkolbisson} (Yaoundé) organise sa journée \og Portes Ouvertes \fg{}. L'administration a confié à votre groupe de quatre élèves, l'équipe \og \emph{Les Géomètres de Nkolbisson} \fg{}, deux missions techniques à réaliser dans la cour de récréation, à l'aide uniquement de matériel simple : un mètre ruban de \SI{10}{m}, une règle de \SI{1}{m}, un bâton de bois droit de \textbf{\SI{1,50}{m}} de long, du crayon, du papier et une calculatrice. Aucun appareil électronique de mesure laser n'est autorisé.

\vspace{0,3cm}
\noindent
\textbf{Mission 1 : La hauteur du mât du drapeau.}
Le mât du drapeau national, planté verticalement au centre de la cour, doit être repeint. Le fournisseur de peinture demande la hauteur exacte pour établir le devis. Il est impossible de monter au sommet. En ce milieu de matinée (vers 10\,h), le soleil projette des ombres nettes. Vous plantez le bâton de \SI{1,50}{m} verticalement sur le sol horizontal, à proximité du mât. Vous mesurez :
\begin{itemize}
    \item L'ombre du bâton : \textbf{\SI{1,20}{m}}.
    \item L'ombre du mât : \textbf{\SI{8,40}{m}}.
\end{itemize}
Le sol est considéré comme horizontal et le soleil assez loin pour que ses rayons soient parallèles.

\vspace{0,3cm}
\noindent
\textbf{Mission 2 : Vérification des lignes de handball.}
Deux lignes blanches, notées $(d_1)$ et $(d_2)$, ont été tracées au sol pour délimiter la zone de but d'un terrain de handball. Elles doivent être \textbf{parallèles} pour respecter la réglementation. Pour les vérifier sans équerre géante, vous choisissez deux points $A$ et $B$ sur $(d_1)$ et deux points $C$ et $D$ sur $(d_2)$ de sorte que les droites $(AC)$ et $(BD)$ se coupent en un point $O$ accessible (au niveau d'un poteau de but). Vous effectuez les mesures suivantes (en mètres) :
\[
OA = 3,6 \quad ; \quad OC = 2,7 \quad ; \quad OB = 4,8 \quad ; \quad OD = 3,5.
\]
Les points sont dans l'ordre $O$--$A$--$B$ sur $(AC)$ et $O$--$C$--$D$ sur $(BD)$ (configuration en \og X \fg{}).

\vspace{0,3cm}
\noindent
Chaque groupe doit produire un \textbf{rapport écrit} (une page recto maximum) contenant : un schéma codé pour chaque mission, l'énoncé des hypothèses, la propriété utilisée (directe, réciproque ou contraposée), les calculs, la conclusion numérique ou logique, et la réponse finale formulée dans le langage de la situation. Une présentation orale de 3 minutes par groupe clôturera l'activité.

\courssub{Consignes}
\noindent
Travailler par groupe de 4. Répartir les rôles : \og Mesureur \fg{}, \og Schématiseur \fg{}, \og Calculateur \fg{}, \og Rédacteur \fg{} (le Rédacteur finalise le rapport, les autres relisent).

\begin{enumerate}
    \item \textbf{Modélisation Mission 1 :} Dessiner un schéma clair de la situation (mât, bâton, ombres, rayons solaires). Coder les longueurs connues et l'inconnue $h$ (hauteur du mât). Justifier en une phrase pourquoi les droites (mât) et (bâton) sont parallèles et pourquoi les rayons solaires forment des sécantes.
    \item \textbf{Application directe :} Écrire l'égalité des rapports issue de la propriété directe de Thalès dans le triangle formé. Calculer $h$ en montrant les étapes (produit en croix, division). Donner la hauteur du mât en mètres, arrondie au centimètre.
    \item \textbf{Modélisation Mission 2 :} Dessiner la configuration en \og X \fg{} avec les points $O, A, B, C, D$. Reporter les quatre mesures. Identifier les deux triangles concernés.
    \item \textbf{Test de parallélisme :} Calculer les deux rapports $\dfrac{OA}{OC}$ et $\dfrac{OB}{OD}$. Les comparer. En déduire, en invoquant explicitement la \textbf{réciproque} ou la \textbf{contraposée} de Thalès, si les droites $(d_1)$ et $(d_2)$ (support de $(AB)$ et $(CD)$) sont parallèles ou non.
    \item \textbf{Rédaction finale :} Rédiger le rapport complet pour les deux missions. Vérifier l'orthographe mathématique (notions de \og triangle \fg{}, \og sécantes \fg{}, \og parallèles \fg{}, \og rapports \fg{}, \og produit en croix \fg{}).
    \item \textbf{Analyse critique (bonus) :} Si le bâton n'avait pas été planté parfaitement vertical, ou si le sol n'était pas horizontal, la propriété directe serait-elle encore applicable ? Justifier brièvement.
\end{enumerate}

\courssub{Ressources à mobiliser}
\begin{propriete}[Propriété directe de Thalès]
Soit un triangle $XYZ$. Si une droite parallèle à $(YZ)$ coupe $(XY)$ en $M$ et $(XZ)$ en $N$, alors les rapports des longueurs des côtés homologues sont égaux :
\[
\frac{XM}{XY} = \frac{XN}{XZ} = \frac{MN}{YZ}.
\]
\emph{Au BEPC : on sait écrire l'égalité des deux rapports utiles (ex: $\frac{XM}{XY}=\frac{XN}{XZ}$) et conclure par le produit en croix.}
\end{propriete}

\begin{propriete}[Propriété réciproque de Thalès]
Soient deux droites $(d_1)$ et $(d_2)$ coupées par deux sécantes en $A,B$ sur $(d_1)$ et $C,D$ sur $(d_2)$, les points étant dans l'ordre $A,B$ et $C,D$ (ou $B,A$ et $D,C$) sur les sécantes. Si $\frac{OA}{OC} = \frac{OB}{OD}$ (ou $\frac{AB}{CD} = \frac{OA}{OC}$), alors $(d_1) \parallel (d_2)$.
\emph{Condition : les points doivent être dans le même ordre sur les deux sécantes (configuration en X ou en \textasciicircum).}
\end{propriete}

\begin{propriete}[Contraposée de Thalès]
Dans la même configuration, si les rapports ne sont \textbf{pas égaux} ($\frac{OA}{OC} \neq \frac{OB}{OD}$), alors les droites $(d_1)$ et $(d_2)$ ne sont \textbf{pas parallèles}.
\emph{C'est l'outil pour \textbf{prouver que deux droites ne sont pas parallèles}.}
\end{propriete}

\begin{methode}[Produit en croix]
Pour résoudre $\frac{a}{b} = \frac{c}{x}$ (avec $x$ l'inconnue) : $a \times x = b \times c \implies x = \frac{b \times c}{a}$.
\end{methode}

\courssub{Démarche et corrigé commenté}
\begin{exoresolu}[Rapport type de l'équipe \og Les Géomètres de Nkolbisson \fg{}]
\textbf{RAPPORT DE MISSION — JOURNÉE PORTES OUVERTES}
\vspace{0,2cm}

\noindent\textbf{1. Mission 1 : Hauteur du mât du drapeau}

\noindent\textbf{Hypothèses et Modélisation :}
\begin{itemize}
    \item Le mât $[MF]$ et le bâton $[BR]$ sont verticaux, donc \textbf{parallèles} ($MF \parallel BR$).
    \item Le sol est horizontal, les points $F$ (pied du mât), $R$ (pied du bâton) et $S$ (bout des ombres) sont alignés.
    \item Les rayons solaires sont parallèles (source lointaine), donc les droites $(MS)$ et $(BS)$ sont confondues (c'est la même droite sécante), et $(MF)$ et $(BR)$ sont les deux droites parallèles coupées par les sécantes $(MB)$ et $(FS)$.
    \item On se place dans le triangle $\triangle MSF$. La droite $(BR)$ est parallèle à $(MF)$ et coupe les deux autres côtés. C'est une configuration en \og \textasciicircum \fg{}.
\end{itemize}

\begin{popfigure}
\begin{tikzpicture}[scale=0.55, >=Stealth]
% Sol
\draw[popline, thick] (-0.5,0) -- (9.5,0) node[right] {Sol horizontal};
% Point S (bout des ombres) à l'origine
\coordinate (S) at (0,0);
\fill[popmuted] (S) circle (2pt) node[below left] {$S$};
% Bâton : base R à 1.2, haut B
\coordinate (R) at (1.2,0);
\coordinate (B) at (1.2,1.5);
\draw[poporange, very thick] (R) -- (B) node[above] {$B$};
\fill[poporange] (R) circle (2pt) node[below] {$R$};
\fill[poporange] (B) circle (2pt);
% Mât : base F à 8.4, haut M
\coordinate (F) at (8.4,0);
\coordinate (M) at (8.4,10.5);
\draw[popblue, very thick] (F) -- (M) node[above] {$M$};
\fill[popblue] (F) circle (2pt) node[below] {$F$};
\fill[popblue] (M) circle (2pt);
% Rayons solaires (pointillés)
\draw[popmuted, dashed, thick] (M) -- (S);
\draw[popmuted, dashed, thick] (B) -- (S);
% Cotes ombres (sous le sol)
\draw[<->, popgreen] (S) -- (R) node[midway, below] {\num{1,20}\,m};
\draw[<->, popgreen] (S) -- (F) node[midway, below, yshift=-0.8cm] {\num{8,40}\,m};
% Cotes hauteurs (à droite des objets)
\draw[<->, poporange] (1.7,0) -- (1.7,1.5) node[midway, right] {\num{1,50}\,m};
\draw[<->, popblue] (8.9,0) -- (8.9,10.5) node[midway, right] {$h$ ?};
% Marque angle droit
\draw[popmuted] (1.15,0) |- (1.2,0.05);
\draw[popmuted] (8.35,0) |- (8.4,0.05);
\end{tikzpicture}
\legende{Configuration en \og \textasciicircum \fg{} : Triangle $MSF$ et droite $BR \parallel MF$. L'ordre au sol est $S$--$R$--$F$.}
\end{popfigure}

\noindent\textbf{Propriété utilisée :} \textbf{Propriété directe de Thalès} dans le triangle $MSF$ car $(BR) \parallel (MF)$.
\[
\frac{BR}{MF} = \frac{RS}{FS}
\]
\vspace{0,1cm}
\noindent\textbf{Calculs :}
On note $h = MF$ (hauteur du mât en mètres).
$BR = \num{1,50}$ ; $RS = \num{1,20}$ (ombre du bâton) ; $FS = \num{8,40}$ (ombre du mât).
\[
\frac{\num{1,50}}{h} = \frac{\num{1,20}}{\num{8,40}}
\]
Produit en croix :
\[
\num{1,50} \times \num{8,40} = h \times \num{1,20}
\]
\[
\num{12,60} = \num{1,20} \times h
\]
\[
h = \frac{\num{12,60}}{\num{1,20}} = \frac{126}{12} = \num{10,50}
\]

\noindent\textbf{Conclusion Mission 1 :} La hauteur du mât est de \textbf{\SI{10,50}{m}}. Le devis de peinture sera établi sur cette base.

\tcblower

\noindent\textbf{2. Mission 2 : Vérification du parallélisme des lignes de handball}

\noindent\textbf{Hypothèses et Modélisation :}
\begin{itemize}
    \item Les lignes $(d_1)$ et $(d_2)$ supportent respectivement les segments $[AB]$ et $[CD]$.
    \item Les droites $(AC)$ et $(BD)$ se coupent en $O$ (poteau de but).
    \item L'ordre des points sur $(AC)$ est $O$--$A$--$B$ (donc $OA=\num{3,6}$, $OB=\num{4,8}$, $AB=\num{1,2}$).
    \item L'ordre des points sur $(BD)$ est $O$--$C$--$D$ (donc $OC=\num{2,7}$, $OD=\num{3,5}$, $CD=\num{0,8}$).
    \item C'est une \textbf{configuration en \og X \fg{}} (deux sécantes coupant deux droites en $O$).
\end{itemize}

\begin{popfigure}
\begin{tikzpicture}[scale=0.7, >=Stealth]
% Point O
\coordinate (O) at (0,0);
\fill[popdark] (O) circle (2pt) node[below left] {$O$};
% Droite d1 (OA) - horizontale
\draw[popblue, thick] (O) -- (5.5,0) node[right] {$(d_1)$};
\coordinate (A) at (3.6,0);
\coordinate (B) at (4.8,0);
\fill[popblue] (A) circle (2pt) node[above] {$A$};
\fill[popblue] (B) circle (2pt) node[above] {$B$};
% Droite d2 (OC) - inclinée (angle ~30 deg)
\draw[poporange, thick] (O) -- (4.5,2.6) node[above right] {$(d_2)$};
\coordinate (C) at (2.7*0.866, 2.7*0.5); % 2.7 * cos30, 2.7 * sin30
\coordinate (D) at (3.5*0.866, 3.5*0.5);
\fill[poporange] (C) circle (2pt) node[below left] {$C$};
\fill[poporange] (D) circle (2pt) node[below left] {$D$};
% Cotes sur d1
\draw[<->, popgreen] (O) -- (A) node[midway, above] {\small \num{3,6}};
\draw[<->, popgreen] (A) -- (B) node[midway, above] {\small \num{1,2}};
% Cotes sur d2
\draw[<->, poppurple] (O) -- (C) node[midway, left, sloped] {\small \num{2,7}};
\draw[<->, poppurple] (C) -- (D) node[midway, left, sloped] {\small \num{0,8}};
% Marque angle
\draw[popmuted] (0.5,0) arc (0:30:0.5) node[midway, right, xshift=0.2cm] {\small $30^\circ$};
\end{tikzpicture}
\legende{Configuration en \og X \fg{} : Sécantes $(OA)$ et $(OC)$.}
\end{popfigure}

\noindent\textbf{Propriété utilisée :} \textbf{Contraposée de Thalès} (car on cherche à prouver que les droites \textbf{ne sont pas} parallèles, ou à vérifier l'égalité pour la réciproque). On compare les rapports des segments sur les deux sécantes issues de $O$.

\noindent\textbf{Calculs des rapports :}
\begin{align*}
\text{Secante } (OA) : \quad & \frac{OA}{OC} = \frac{\num{3,6}}{\num{2,7}} = \frac{36}{27} = \frac{4}{3} \approx 1,333\ldots \\[4pt]
\text{Secante } (OB) : \quad & \frac{OB}{OD} = \frac{\num{4,8}}{\num{3,5}} = \frac{48}{35} \approx 1,371\ldots
\end{align*}

\noindent\textbf{Comparaison :}
\[
\frac{4}{3} \neq \frac{48}{35} \quad \text{car} \quad 4 \times 35 = 140 \neq 144 = 3 \times 48.
\]
Les rapports ne sont \textbf{pas égaux}.

\noindent\textbf{Conclusion Mission 2 :} D'après la \textbf{contraposée du théorème de Thalès}, puisque $\frac{OA}{OC} \neq \frac{OB}{OD}$, les droites $(AB)$ et $(CD)$ — c'est-à-dire les lignes $(d_1)$ et $(d_2)$ — \textbf{ne sont pas parallèles}.
Le tracé au sol est \textbf{non conforme} à la réglementation handball. Il faut refaire le traçage.

\vspace{0,3cm}
\noindent\textbf{3. Analyse critique (Réponse bonus)}
Si le bâton n'est pas vertical ou le sol pas horizontal, les droites (mât) et (bâton) ne sont plus parallèles. L'hypothèse fondamentale de Thalès (deux droites parallèles coupées par deux sécantes) n'est plus satisfaite. Les triangles formés ne sont plus semblables par Thalès. Le calcul donnerait une valeur erronée. \textbf{La verticalité du bâton et l'horizontalité du sol sont des conditions nécessaires.}

\vspace{0,3cm}
\noindent\textit{Fait à Nkolbisson, le \today. Signatures : [Mesureur] [Schématiseur] [Calculateur] [Rédacteur]}
\end{exoresolu}

\courssub{Bilan de l'activité}
Cette activité a permis de valider les compétences cibles du programme de 3ème :
\begin{itemize}
    \item \textbf{Modéliser :} Passer d'une situation physique (ombres, lignes au sol) à un modèle géométrique (triangles, configurations en \og X \fg{} ou \og \textasciicircum \fg{}) en explicitant les hypothèses (parallélisme des verticales, horizontalité du sol, alignement des points, ordre sur les sécantes).
    \item \textbf{Raisonner :} Choisir la bonne propriété (directe pour calculer, réciproque/contreposée pour décider du parallélisme) et l'énoncer correctement avec ses conditions.
    \item \textbf{Calculer :} Manipuler les rapports, simplifier les fractions, effectuer le produit en croix sans erreur, gérer les unités (mètres) et la précision (arrondi au cm).
    \item \textbf{Communiquer :} Produire un écrit mathématique structuré (schéma codé $\to$ hypothèses $\to$ propriété $\to$ calculs $\to$ conclusion $\to$ réponse) et le présenter oralement.
    \item \textbf{Esprit critique :} Identifier la fragilité du modèle si les hypothèses idéales (verticalité, horizontalité, parallélisme des rayons) ne sont pas respectées.
\end{itemize}
L'ancrage dans le contexte camerounais (cour de collège, mât du drapeau, lignes de handball, devis peinture) a donné du sens aux abstractions géométriques, transformant le théorème de Thalès en un outil concret de \og géomètre du quartier \fg{}.

\newpage

\coursec{Méthodes et exercices résolus}

\courssub{Méthode 1 : Calculer une longueur avec la propriété directe de Thalès}

La propriété directe de Thalès est l'outil central de cette leçon : dès que l'on reconnaît deux droites sécantes coupées par deux droites parallèles, elle permet de transformer un problème de géométrie en un simple calcul de proportionnalité. Voici la démarche complète, à reproduire à l'identique dans vos copies.

\begin{methode}[Calculer une longueur grâce à Thalès]
Pour calculer une longueur dans une configuration de Thalès :
\begin{enumerate}
\item \textbf{Repérer la configuration} : deux droites sécantes en un point $A$, deux points $M$ et $B$ sur l'une, deux points $N$ et $C$ sur l'autre, avec $(MN) \parallel (BC)$.
\item \textbf{Vérifier et citer les hypothèses} : les points $A$, $M$, $B$ sont alignés ; les points $A$, $N$, $C$ sont alignés (dans le même ordre) ; $(MN) \parallel (BC)$.
\item \textbf{Écrire les égalités de Thalès} : $\dfrac{AM}{AB} = \dfrac{AN}{AC} = \dfrac{MN}{BC}$.
\item \textbf{Choisir les deux rapports utiles} : celui entièrement connu et celui qui contient l'inconnue.
\item \textbf{Faire un produit en croix} : si $\dfrac{a}{b} = \dfrac{c}{d}$, alors $a \times d = b \times c$, donc $a = \dfrac{b \times c}{d}$.
\item \textbf{Conclure} par une phrase avec l'unité.
\end{enumerate}
\textbf{Réflexe} : avant tout calcul, écrire la phrase magique : « Les points $A$, $M$, $B$ et $A$, $N$, $C$ sont alignés dans le même ordre et $(MN) \parallel (BC)$, donc d'après la propriété de Thalès... ».
\end{methode}

\begin{popfigure}
\begin{center}
\begin{tikzpicture}[scale=0.9]
\coordinate (A) at (0,0);
\coordinate (B) at (5,0);
\coordinate (C) at (6,3.6);
\coordinate (M) at (2,0);
\coordinate (N) at (2.4,1.44);
\draw[popdark,thick] (A) -- (B);
\draw[popdark,thick] (A) -- (C);
\draw[popblue,very thick] (M) -- (N);
\draw[poporange,very thick] (B) -- (C);
\fill (A) circle (1.5pt) node[below left] {$A$};
\fill (M) circle (1.5pt) node[below] {$M$};
\fill (B) circle (1.5pt) node[below] {$B$};
\fill (N) circle (1.5pt) node[above left] {$N$};
\fill (C) circle (1.5pt) node[above right] {$C$};
\node[popblue] at (1.9,0.85) {\small $(MN)$};
\node[poporange] at (5.9,1.6) {\small $(BC)$};
\end{tikzpicture}
\end{center}
\legende{Configuration de Thalès : $(MN) \parallel (BC)$, les rapports partent tous du sommet $A$.}
\end{popfigure}

\begin{exoresolu}[Longueur d'un poteau électrique]
Pour électrifier le quartier, un technicien de la SONATREL plante un piquet vertical $[MN]$ de $1{,}5$ m entre l'œil $A$ de l'observateur (posé au sol) et le pied $B$ d'un poteau vertical $[BC]$. On a : $AM = 2$ m, $AB = 10$ m, et $(MN) \parallel (BC)$ (les deux sont verticaux). Calculer la hauteur $BC$ du poteau.
\tcblower
\textbf{Étape 1 : les hypothèses.} Les points $A$, $M$, $B$ sont alignés ; les points $A$, $N$, $C$ sont alignés (droite de visée) ; et $(MN) \parallel (BC)$ car le piquet et le poteau sont tous deux verticaux.

\textbf{Étape 2 : application de Thalès.} D'après la propriété directe de Thalès :
\[ \frac{AM}{AB} = \frac{AN}{AC} = \frac{MN}{BC} \]

\textbf{Étape 3 : choix des rapports.} On connaît $AM$, $AB$ et $MN$ ; on cherche $BC$. On utilise :
\[ \frac{AM}{AB} = \frac{MN}{BC} \quad \text{soit} \quad \frac{2}{10} = \frac{1{,}5}{BC} \]

\textbf{Étape 4 : produit en croix.}
\[ 2 \times BC = 10 \times 1{,}5 \quad \Rightarrow \quad BC = \frac{10 \times 1{,}5}{2} = \frac{15}{2} = 7{,}5 \]

\textbf{Conclusion :} le poteau mesure $7{,}5$ m de haut.
\end{exoresolu}

\courssub{Méthode 2 : Calculer une longueur dans le cas « papillon » (points de part et d'autre)}

Il arrive que les deux droites sécantes se coupent en $A$ et que les points $M$ et $B$ (ainsi que $N$ et $C$) soient situés \emph{de part et d'autre} de $A$ : on parle alors de configuration « papillon » ou « croisée ». Bonne nouvelle : la propriété de Thalès s'applique exactement de la même façon, à condition de travailler avec des longueurs.

\begin{methode}[Thalès en configuration croisée]
Quand les deux droites se coupent en $A$ et que les points sont \emph{de part et d'autre} de $A$ (configuration « papillon ») :
\begin{enumerate}
\item Vérifier que $(MN) \parallel (BC)$ et que $A$, $M$, $B$ d'une part, $A$, $N$, $C$ d'autre part sont alignés.
\item Écrire quand même $\dfrac{AM}{AB} = \dfrac{AN}{AC} = \dfrac{MN}{BC}$ : la propriété reste vraie.
\item Attention : utiliser les \textbf{longueurs} (toujours positives), jamais de signes.
\item Résoudre par produit en croix et conclure.
\end{enumerate}
\textbf{Réflexe} : repérer le sommet commun $A$ ; les rapports partent toujours de $A$.
\end{methode}

\begin{exoresolu}[Largeur d'une rivière]
Pour mesurer la largeur $AB$ d'une rivière sans la traverser, un élève de 3ème place des jalons : $A$ et $B$ sont sur les deux rives, et il construit les alignements $C$, $M$, $A$ et $C$, $N$, $B$ de sorte que $(MN) \parallel (AB)$, le point $C$ étant sur sa rive. Il mesure : $CM = 4$ m, $CA = 12$ m, $MN = 3$ m. Calculer la largeur $AB$ de la rivière.
\tcblower
\textbf{Hypothèses :} les points $C$, $M$, $A$ sont alignés ; les points $C$, $N$, $B$ sont alignés ; et $(MN) \parallel (AB)$.

D'après la propriété directe de Thalès appliquée au sommet $C$ :
\[ \frac{CM}{CA} = \frac{CN}{CB} = \frac{MN}{AB} \]
On utilise les rapports connus :
\[ \frac{CM}{CA} = \frac{MN}{AB} \quad \text{soit} \quad \frac{4}{12} = \frac{3}{AB} \]
Produit en croix :
\[ 4 \times AB = 12 \times 3 \quad \Rightarrow \quad AB = \frac{36}{4} = 9 \]
\textbf{Conclusion :} la rivière mesure $9$ m de large.
\end{exoresolu}

\courssub{Méthode 3 : Démontrer que deux droites sont parallèles (réciproque de Thalès)}

Jusqu'ici, le parallélisme était une \emph{hypothèse} et la longueur une \emph{inconnue}. La réciproque de Thalès inverse les rôles : on connaît des longueurs et l'on veut \emph{prouver} que deux droites sont parallèles. La rédaction change donc complètement : on ne peut plus écrire l'égalité des rapports d'office, il faut la vérifier par deux calculs indépendants.

\begin{methode}[Prouver un parallélisme avec la réciproque]
\begin{enumerate}
\item Identifier les deux droites sécantes en $A$ et les points alignés $A$, $M$, $B$ et $A$, $N$, $C$ \textbf{dans le même ordre}.
\item Calculer \textbf{séparément} les deux rapports $\dfrac{AM}{AB}$ et $\dfrac{AN}{AC}$ (sous forme de fractions, puis simplifier ou donner la même écriture décimale).
\item Comparer : si $\dfrac{AM}{AB} = \dfrac{AN}{AC}$, alors d'après la \textbf{réciproque de la propriété de Thalès}, $(MN) \parallel (BC)$.
\item Rédiger la conclusion en citant la réciproque et l'alignement des points dans le même ordre.
\end{enumerate}
\textbf{Réflexe} : ne jamais écrire l'égalité des rapports avant de l'avoir vérifiée par deux calculs séparés.
\end{methode}

\begin{exoresolu}[Des étagères bien parallèles ?]
Un menuisier de Yaoundé fixe deux étagères $[MN]$ et $[BC]$ sur un montant incliné. Sur le schéma, $A$, $M$, $B$ sont alignés, $A$, $N$, $C$ sont alignés dans le même ordre, avec $AM = 3$ cm, $AB = 7{,}5$ cm, $AN = 4$ cm, $AC = 10$ cm. Les deux étagères sont-elles parallèles ?
\tcblower
\textbf{Calcul séparé des deux rapports :}
\[ \frac{AM}{AB} = \frac{3}{7{,}5} = \frac{30}{75} = \frac{2}{5} \qquad ; \qquad \frac{AN}{AC} = \frac{4}{10} = \frac{2}{5} \]
\textbf{Comparaison :} on constate que $\dfrac{AM}{AB} = \dfrac{AN}{AC} = \dfrac{2}{5}$.

De plus, les points $A$, $M$, $B$ et $A$, $N$, $C$ sont alignés dans le même ordre.

Donc, d'après la \textbf{réciproque de la propriété de Thalès}, les droites $(MN)$ et $(BC)$ sont parallèles.

\textbf{Conclusion :} oui, les deux étagères sont bien parallèles : le travail du menuisier est soigné.
\end{exoresolu}

\courssub{Méthode 4 : Prouver que deux droites ne sont pas parallèles (contraposée)}

Et si les rapports sont différents ? On ne peut pas utiliser la réciproque (qui exige l'égalité), mais on peut raisonner par l'absurde grâce à la propriété directe : c'est la \cle{contraposée} de la propriété de Thalès.

\begin{methode}[Prouver un non-parallélisme]
\begin{enumerate}
\item Calculer séparément $\dfrac{AM}{AB}$ et $\dfrac{AN}{AC}$.
\item Si $\dfrac{AM}{AB} \neq \dfrac{AN}{AC}$, alors les droites $(MN)$ et $(BC)$ \textbf{ne sont pas parallèles} (sinon, la propriété directe de Thalès donnerait l'égalité des rapports, ce qui est absurde).
\item Rédiger : « Si $(MN)$ était parallèle à $(BC)$, on aurait $\dfrac{AM}{AB} = \dfrac{AN}{AC}$ ; or ce n'est pas le cas, donc $(MN)$ et $(BC)$ ne sont pas parallèles. »
\end{enumerate}
\textbf{Attention} : pour la contraposée, l'hypothèse « points dans le même ordre » n'est pas exigée ; il suffit que les rapports soient différents.
\end{methode}

\begin{exoresolu}[Un mur mal construit]
Sur un chantier à Douala, on vérifie la verticalité de deux montants. On a : $A$, $M$, $B$ alignés et $A$, $N$, $C$ alignés, avec $AM = 2$ m, $AB = 5$ m, $AN = 3$ m, $AC = 8$ m. Les droites $(MN)$ et $(BC)$ sont-elles parallèles ?
\tcblower
\textbf{Calcul des rapports :}
\[ \frac{AM}{AB} = \frac{2}{5} = 0{,}4 \qquad ; \qquad \frac{AN}{AC} = \frac{3}{8} = 0{,}375 \]
\textbf{Comparaison :} $0{,}4 \neq 0{,}375$, donc $\dfrac{AM}{AB} \neq \dfrac{AN}{AC}$.

\textbf{Raisonnement :} si les droites $(MN)$ et $(BC)$ étaient parallèles, la propriété directe de Thalès imposerait $\dfrac{AM}{AB} = \dfrac{AN}{AC}$. Or ces rapports sont différents.

\textbf{Conclusion :} les droites $(MN)$ et $(BC)$ ne sont pas parallèles : le mur doit être repris.
\end{exoresolu}

\courssub{Méthode 5 : Rédiger et justifier une démarche complète (problème de vie courante)}

Au BEPC, les exercices de Thalès sont souvent présentés sous forme de problème concret. Le correcteur attend une démarche complète et justifiée, pas seulement un résultat numérique.

\begin{methode}[Résoudre un problème concret du début à la fin]
\begin{enumerate}
\item \textbf{Traduire} la situation par un schéma codé : nommer les points, reporter les mesures, coder le parallélisme.
\item \textbf{Identifier} la propriété à utiliser : directe (on connaît le parallélisme, on cherche une longueur) ou réciproque (on connaît des longueurs, on cherche le parallélisme).
\item \textbf{Rédiger} : hypothèses, citation de la propriété, calculs, produit en croix.
\item \textbf{Vérifier} la cohérence du résultat (ordre de grandeur : le rapport $k$ doit être inférieur à 1 si on passe du petit triangle au grand).
\item \textbf{Conclure} par une phrase répondant à la question posée, avec l'unité.
\end{enumerate}
\end{methode}

\begin{exoresolu}[L'ombre du mât du drapeau]
Au collège, le mât du drapeau $[BC]$ est vertical. À midi, son ombre au sol mesure $AC = 6$ m. Au même moment, un élève place verticalement un bâton $[MN]$ de $1{,}2$ m dont l'ombre mesure $AN = 1{,}5$ m (le soleil envoie des rayons parallèles, donc $(MN) \parallel (BC)$ et les points $A$, $N$, $C$ sont alignés sur le sol, $A$, $M$, $B$ alignés sur le rayon lumineux). Calculer la hauteur du mât.
\tcblower
\textbf{Schéma et hypothèses :} $A$ est l'extrémité commune des ombres ; $A$, $N$, $C$ sont alignés sur le sol ; $A$, $M$, $B$ sont alignés sur le rayon du soleil ; $(MN) \parallel (BC)$ car le bâton et le mât sont verticaux.

\textbf{Propriété de Thalès :}
\[ \frac{AN}{AC} = \frac{AM}{AB} = \frac{MN}{BC} \]
\textbf{Choix des rapports et calcul :}
\[ \frac{AN}{AC} = \frac{MN}{BC} \quad \text{soit} \quad \frac{1{,}5}{6} = \frac{1{,}2}{BC} \]
\[ 1{,}5 \times BC = 6 \times 1{,}2 \quad \Rightarrow \quad BC = \frac{7{,}2}{1{,}5} = 4{,}8 \]
\textbf{Vérification :} le rapport $\dfrac{AN}{AC} = \dfrac{1{,}5}{6} = 0{,}25 < 1$, et $1{,}2 \div 0{,}25 = 4{,}8$ : cohérent.

\textbf{Conclusion :} le mât du drapeau mesure $4{,}8$ m de haut.
\end{exoresolu}

\courssub{Méthode 6 : Calculer une longueur « morceau restant » (type $MB$ ou $NC$)}

Dernière situation classique : l'inconnue n'est pas une longueur partant du sommet $A$, mais un \emph{morceau} comme $MB$. La propriété de Thalès ne donne jamais directement un tel morceau : il faut passer par la longueur totale, puis soustraire.

\begin{methode}[Trouver une longueur partielle]
Quand l'inconnue est un morceau comme $MB$ (et non $AB$) :
\begin{enumerate}
\item Utiliser Thalès pour calculer d'abord la longueur \textbf{totale} $AB$ (ou le rapport $k$).
\item En déduire le morceau par soustraction : $MB = AB - AM$.
\item Conclure avec l'unité.
\end{enumerate}
\textbf{Réflexe} : Thalès ne donne jamais directement $MB$ : il donne des rapports partant du sommet $A$.
\end{methode}

\begin{exoresolu}[La portion restante d'un chemin de câble]
Un câble téléphonique CAMTEL suit la droite $(AB)$. On sait que $A$, $M$, $B$ sont alignés, $A$, $N$, $C$ alignés dans le même ordre, $(MN) \parallel (BC)$, avec $AM = 4$ m, $AN = 3$ m et $AC = 9$ m. Calculer la distance $MB$.
\tcblower
\textbf{Étape 1 : Thalès pour trouver $AB$.} Les points $A$, $M$, $B$ et $A$, $N$, $C$ sont alignés dans le même ordre et $(MN) \parallel (BC)$, donc d'après la propriété de Thalès :
\[ \frac{AM}{AB} = \frac{AN}{AC} \quad \text{soit} \quad \frac{4}{AB} = \frac{3}{9} = \frac{1}{3} \]
Produit en croix : $AB = 4 \times 3 = 12$ m.

\textbf{Étape 2 : soustraction.}
\[ MB = AB - AM = 12 - 4 = 8 \]
\textbf{Conclusion :} la distance $MB$ est de $8$ m.
\end{exoresolu}

\begin{savaistu}[Thalès et la pyramide de Khéops]
Selon la légende rapportée par les historiens grecs, Thalès de Milet (vers 625 -- 547 av. J.-C.) aurait mesuré la hauteur de la grande pyramide d'Égypte en comparant son ombre à celle d'un simple bâton planté verticalement : exactement la méthode de l'exercice du mât du drapeau ! Vingt-cinq siècles plus tard, les géomètres et les topographes camerounais utilisent toujours ce principe pour mesurer des distances inaccessibles : largeur d'une rivière, hauteur d'un pylône, profondeur d'un puits.
\end{savaistu}

\begin{attention}[Les 5 erreurs qui coûtent des points au BEPC]
\begin{enumerate}
\item \textbf{Oublier de citer les hypothèses.} Écrire « d'après Thalès » sans avoir dit que les points sont alignés et que les droites sont parallèles fait perdre des points : la propriété n'est applicable que sous ces conditions.
\item \textbf{Mélanger les rapports.} On ne peut pas écrire $\dfrac{AM}{AB} = \dfrac{MN}{AC}$ : chaque rapport doit être homogène (deux longueurs du même triangle, ou deux côtés correspondants). Vérifie que les numérateurs sont dans le petit triangle et les dénominateurs dans le grand.
\item \textbf{Erreur de produit en croix.} De $\dfrac{2}{10} = \dfrac{1{,}5}{BC}$, on tire $BC = \dfrac{10 \times 1{,}5}{2}$ et non $\dfrac{2 \times 1{,}5}{10}$. Pose toujours l'égalité des produits en croix avant de diviser.
\item \textbf{Confondre directe et réciproque.} Pour prouver un parallélisme, on calcule les deux rapports \emph{séparément} puis on conclut avec la \textbf{réciproque} ; on n'a pas le droit de supposer l'égalité des rapports avant de l'avoir démontrée.
\item \textbf{Oublier l'ordre des points dans la réciproque.} La réciproque de Thalès exige que $A$, $M$, $B$ et $A$, $N$, $C$ soient alignés \emph{dans le même ordre} : sans cette précision, la conclusion peut être fausse et la rédaction est sanctionnée. Pense aussi à conclure par une phrase avec l'unité.
\end{enumerate}
\end{attention}

\newpage

\coursec{Exercices}

\courssub{Je vérifie mes connaissances}

\begin{exercice}[Configuration de Thalès : Vrai ou Faux ?]
Pour chacune des affirmations suivantes, dire si elle est \textbf{VRAIE} ou \textbf{FAUSSE} en justifiant brièvement par un contre-exemple ou une propriété du cours.
\begin{enumerate}
    \item Dans un triangle $ABC$, si $M \in [AB]$, $N \in [AC]$ et $(MN) \parallel (BC)$, alors $\dfrac{AM}{AB} = \dfrac{AN}{AC} = \dfrac{MN}{BC}$.
    \item La propriété directe de Thalès s'applique uniquement si les points $M$ et $N$ sont les milieux des côtés $[AB]$ et $[AC]$.
    \item Si dans un triangle $ABC$, des points $M \in [AB]$ et $N \in [AC]$ vérifient $\dfrac{AM}{AB} = \dfrac{AN}{AC}$, alors les droites $(MN)$ et $(BC)$ sont parallèles.
    \item Dans la configuration « papillon » (deux droites sécantes coupées par deux droites parallèles), les rapports des segments sur une sécante sont égaux aux rapports des segments sur l'autre sécante.
    \item Si $\dfrac{AM}{AB} = \dfrac{AN}{AC}$ mais que les points $M$ et $N$ ne sont pas placés dans le même ordre sur les demi-droites $[AB)$ et $[AC)$, alors $(MN) \parallel (BC)$.
\end{enumerate}
\end{exercice}

\begin{exercice}[Calculs directs d'application]
Dans le triangle $ABC$, on place $M \in [AB]$ et $N \in [AC]$ tels que $(MN) \parallel (BC)$.
Calculer la longueur demandée en justifiant par l'égalité des rapports.
\begin{enumerate}
    \item $AB = 12\,\text{cm},\; AM = 4\,\text{cm},\; AC = 15\,\text{cm}$. Trouver $AN$.
    \item $AB = 20\,\text{cm},\; AM = 8\,\text{cm},\; BC = 18\,\text{cm}$. Trouver $MN$.
    \item $AC = 18\,\text{cm},\; AN = 6\,\text{cm},\; BC = 24\,\text{cm}$. Trouver $MN$.
    \item $AM = 5\,\text{cm},\; MB = 10\,\text{cm},\; AN = 7\,\text{cm}$. Trouver $NC$.
\end{enumerate}
\end{exercice}

\begin{exercice}[Reconnaissance de la configuration : QCM]
On considère la figure ci-dessous où les droites $(MN)$ et $(BC)$ sont parallèles.
\begin{popfigure}
\begin{tikzpicture}[scale=0.9, >=Stealth]
\coordinate (A) at (0,0);
\coordinate (B) at (5,3.5);
\coordinate (C) at (7,0);
\coordinate (M) at ($(A)!0.4!(B)$);
\coordinate (N) at ($(A)!0.4!(C)$);
\draw[thick, popdark] (A) -- (B) -- (C) -- cycle;
\draw[thick, popblue] (M) -- (N);
\draw[dashed, popmuted] (A) -- ($(A)!1.2!(B)$) (A) -- ($(A)!1.2!(C)$);
\node[above left] at (A) {$A$};
\node[above right] at (B) {$B$};
\node[below right] at (C) {$C$};
\node[above left] at (M) {$M$};
\node[below left] at (N) {$N$};
\end{tikzpicture}
\legende{Configuration triangle (sommets $A$ en haut).}
\end{popfigure}
Parmi les égalités suivantes, laquelle n'est \textbf{PAS} une conséquence de Thalès dans cette figure ?
\begin{enumerate}
    \item $\dfrac{AM}{AB} = \dfrac{AN}{AC}$
    \item $\dfrac{AM}{MB} = \dfrac{AN}{NC}$
    \item $\dfrac{AB}{AC} = \dfrac{AM}{AN}$
    \item $\dfrac{MB}{AB} = \dfrac{NC}{AC}$
    \item $\dfrac{AM}{AB} = \dfrac{MN}{BC} = \dfrac{AN}{NC}$
\end{enumerate}
\end{exercice}

\begin{exercice}[Contraposée : Droites non parallèles]
Dans un triangle $RST$, on a placé $U \in [RS]$ et $V \in [RT]$.
Les mesures sont : $RU = 6\,\text{cm},\; RS = 15\,\text{cm},\; RV = 4\,\text{cm},\; RT = 12\,\text{cm}$.
Montrer que les droites $(UV)$ et $(ST)$ ne sont pas parallèles.
\end{exercice}

\courssub{Je m'entraîne}

\begin{exercice}[Calculs de longueurs dans un triangle (Type BEPC)]
Soit un triangle $ABC$ tel que $AB = 10\,\text{cm}$, $AC = 15\,\text{cm}$ et $BC = 12\,\text{cm}$.
On place $M \in [AB]$ tel que $AM = 4\,\text{cm}$.
La parallèle à $(BC)$ passant par $M$ coupe $[AC]$ en $N$.
\begin{enumerate}
    \item Faire une figure à main levée.
    \item Calculer $AN$. Justifier chaque étape.
    \item Calculer $MN$. Justifier.
    \item Calculer $MB$ et $NC$.
\end{enumerate}
\end{exercice}

\begin{exercice}[Problème concret : Hauteur d'un poteau électrique]
À Douala, au quartier Bonanjo, un poteau électrique de la société ENEO projette une ombre de $8,40\,\text{m}$ sur le sol horizontal au même instant où un élève de $1,40\,\text{m}$ de haut projette une ombre de $1,20\,\text{m}$.
On suppose que les rayons du soleil sont parallèles et que le poteau et l'élève sont perpendiculaires au sol.
\begin{enumerate}
    \item Modéliser la situation par un schéma (deux triangles rectangles semblables).
    \item Calculer la hauteur du poteau électrique. Arrondir au centimètre près.
    \item Si l'élève s'éloigne du poteau, la longueur de son ombre change-t-elle ? Justifier.
\end{enumerate}
\end{exercice}

\begin{exercice}[Réciproque de Thalès : Parallélisme ou pas ?]
On donne un triangle $ABC$. Les points $M$ et $N$ appartiennent respectivement aux segments $[AB]$ et $[AC]$.
Dans chaque cas, déterminer si $(MN) \parallel (BC)$ en justifiant par le calcul des rapports.
\begin{enumerate}
    \item $AM = 3,5\,\text{cm},\; AB = 8,4\,\text{cm},\; AN = 2,5\,\text{cm},\; AC = 6\,\text{cm}$.
    \item $AM = 5\,\text{cm},\; MB = 7\,\text{cm},\; AN = 4\,\text{cm},\; NC = 5,6\,\text{cm}$.
    \item $AB = 18\,\text{cm},\; AM = 8\,\text{cm},\; AC = 13,5\,\text{cm},\; AN = 6\,\text{cm}$.
\end{enumerate}
\end{exercice}

\begin{exercice}[Configuration « Papillon » : Deux sécantes et deux parallèles]
Les droites $(BM)$ et $(CN)$ se coupent en $A$. Les droites $(BC)$ et $(MN)$ sont parallèles.
On donne : $AM = 6\,\text{cm},\; AB = 15\,\text{cm},\; AN = 4\,\text{cm}$.
\begin{popfigure}
\begin{tikzpicture}[scale=0.9, >=Stealth]
\coordinate (A) at (0,0);
\coordinate (B) at (-3,-2);
\coordinate (C) at (3.5,-2.5);
\coordinate (M) at ($(A)!0.4!(B)$);
\coordinate (N) at ($(A)!0.4!(C)$);
\draw[thick, popdark] (B) -- (C);
\draw[thick, popblue] (M) -- (N);
\draw[thick, poporange] (A) -- ($(A)!1.3!(B)$);
\draw[thick, poporange] (A) -- ($(A)!1.3!(C)$);
\node[above] at (A) {$A$};
\node[below left] at (B) {$B$};
\node[below right] at (C) {$C$};
\node[left] at (M) {$M$};
\node[right] at (N) {$N$};
\end{tikzpicture}
\legende{Configuration papillon (deux sécantes).}
\end{popfigure}
\begin{enumerate}
    \item Calculer $AC$.
    \item Calculer $MN$ si $BC = 10\,\text{cm}$.
    \item Vérifier la cohérence des données en calculant le rapport $\dfrac{AM}{AB}$ et $\dfrac{AN}{AC}$.
    \item Si on prolonge $[AM]$ au-delà de $M$ d'une longueur $MP = 3\,\text{cm}$, la droite $(PN)$ est-elle parallèle à $(BC)$ ? Justifier.
\end{enumerate}
\end{exercice}

\courssub{Je me prépare au Probatoire}

\begin{exercice}[Partage d'un terrain triangulaire — Modélisation]
Un terrain a la forme d'un triangle $ABC$ d'aire $2700\,\text{m}^2$.
On veut le diviser en deux parties par une clôture $(MN)$ parallèle à $[BC]$, avec $M \in [AB]$ et $N \in [AC]$, de sorte que la parcelle $AMN$ (destinée à l'habitation) représente exactement $\dfrac{4}{9}$ de l'aire totale du terrain.
\begin{enumerate}
    \item Exprimer le rapport des aires $\dfrac{\mathcal{A}_{AMN}}{\mathcal{A}_{ABC}}$ en fonction du rapport de similitude $k = \dfrac{AM}{AB}$.
    \item En déduire la valeur de $k = \dfrac{AM}{AB}$.
    \item Si $AB = 60\,\text{m}$, calculer la distance $AM$ à partir de $A$ pour placer le piquet $M$.
    \item Calculer la longueur de la clôture $MN$ si $BC = 45\,\text{m}$.
\end{enumerate}
\end{exercice}

\begin{exercice}[Mesure de la largeur du fleuve Wouri — Application géométrique]
Pour mesurer la largeur du fleuve Wouri en un point étroit sans le traverser, un géomètre utilise la méthode suivante :
\begin{itemize}
    \item Il plante un jalon $A$ sur la rive de départ.
    \item Il repère un arbre remarquable $B$ sur la rive d'en face.
    \item Sur la rive de départ, il plante un jalon $C$ aligné avec $A$ et $B$ (donc $A, C, B$ alignés dans cet ordre).
    \item Il plante un jalon $D$ sur la rive de départ tel que $(CD) \perp (AB)$.
    \item Il plante un jalon $E$ sur $[AD]$ tel que $(CE) \parallel (DB)$.
\end{itemize}
Les mesures relevées sont : $AC = 30\,\text{m},\; AD = 50\,\text{m},\; AE = 20\,\text{m}$.
\begin{enumerate}
    \item Faire une figure codée de la situation.
    \item Justifier que les triangles $ACE$ et $ADB$ sont semblables (configuration de Thalès).
    \item Calculer la largeur du fleuve $AB$.
    \item Calculer la distance $CE$ si l'arbre $B$ est à $40\,\text{m}$ du pied de la perpendiculaire $D$ (c'est-à-dire $DB = 40\,\text{m}$).
\end{enumerate}
\end{exercice}

\begin{exercice}[Synthèse : Figure complexe et raisonnement par l'absurde]
Soit un triangle $ABC$. On place $M \in [AB]$ et $N \in [AC]$.
On suppose que $\dfrac{AM}{AB} = \dfrac{2}{5}$ et $\dfrac{AN}{AC} = \dfrac{3}{8}$.
\begin{enumerate}
    \item Les droites $(MN)$ et $(BC)$ sont-elles parallèles ? Justifier par la contraposée de Thalès.
    \item On construit le point $N'$ sur $[AC]$ tel que $(MN') \parallel (BC)$. Calculer $\dfrac{AN'}{AC}$.
    \item Comparer les longueurs $AN$ et $AN'$. En déduire la position relative de $N$ et $N'$ sur le segment $[AC]$.
    \item On trace la parallèle à $(AB)$ passant par $N$. Elle coupe $(BC)$ en $P$.
        \begin{itemize}
            \item Calculer $\dfrac{CN}{CA}$.
            \item En déduire $\dfrac{CP}{CB}$.
            \item Si $BC = 20\,\text{cm}$, calculer $CP$.
        \end{itemize}
\end{enumerate}
\end{exercice}

\newpage

\coursec{Corrigés des exercices}

\begin{corrige}[Exercice 1 — Configuration de Thalès : Vrai ou Faux ?]
\begin{enumerate}
    \item \textbf{VRAI.} C'est exactement la propriété directe de Thalès : dans le triangle $ABC$, avec $M \in [AB]$, $N \in [AC]$ et $(MN) \parallel (BC)$, on a $\dfrac{AM}{AB} = \dfrac{AN}{AC} = \dfrac{MN}{BC}$.
    \item \textbf{FAUX.} La propriété directe de Thalès s'applique dès que $(MN) \parallel (BC)$, quelle que soit la position de $M$ sur $[AB]$. Contre-exemple : $AM = 3$, $AB = 9$, $AC = 12$ donne $AN = 4$, alors que $M$ n'est pas le milieu de $[AB]$ (car $AM = 3 \neq 4,5 = MB$).
    \item \textbf{VRAI.} C'est la propriété réciproque de Thalès : les points $A, M, B$ et $A, N, C$ sont alignés dans le même ordre et $\dfrac{AM}{AB} = \dfrac{AN}{AC}$, donc $(MN) \parallel (BC)$.
    \item \textbf{VRAI.} Dans la configuration papillon (deux sécantes coupées par des parallèles), le théorème de Thalès donne $\dfrac{AM}{AB} = \dfrac{AN}{AC} = \dfrac{MN}{BC}$ : les rapports sur les deux sécantes sont égaux.
    \item \textbf{FAUX.} La réciproque de Thalès exige que les points soient placés \textbf{dans le même ordre} sur les demi-droites $[AB)$ et $[AC)$. Sans cette condition d'ordre, l'égalité des rapports ne suffit pas à conclure au parallélisme.
\end{enumerate}
\end{corrige}

\begin{corrige}[Exercice 2 — Calculs directs d'application]
Dans tous les cas, $(MN) \parallel (BC)$, donc par la propriété directe de Thalès : $\dfrac{AM}{AB} = \dfrac{AN}{AC} = \dfrac{MN}{BC}$.
\begin{enumerate}
    \item $\dfrac{AM}{AB} = \dfrac{AN}{AC}$ donne $\dfrac{4}{12} = \dfrac{AN}{15}$, d'où $AN = \dfrac{4 \times 15}{12} = \dfrac{60}{12} = 5\,\text{cm}$.
    \item $\dfrac{AM}{AB} = \dfrac{MN}{BC}$ donne $\dfrac{8}{20} = \dfrac{MN}{18}$, d'où $MN = \dfrac{8 \times 18}{20} = \dfrac{144}{20} = 7,2\,\text{cm}$.
    \item $\dfrac{AN}{AC} = \dfrac{MN}{BC}$ donne $\dfrac{6}{18} = \dfrac{MN}{24}$, d'où $MN = \dfrac{6 \times 24}{18} = \dfrac{144}{18} = 8\,\text{cm}$.
    \item $AB = AM + MB = 5 + 10 = 15\,\text{cm}$. Alors $\dfrac{AM}{AB} = \dfrac{AN}{AC}$ donne $\dfrac{5}{15} = \dfrac{7}{AC}$, d'où $AC = \dfrac{7 \times 15}{5} = 21\,\text{cm}$. Donc $NC = AC - AN = 21 - 7 = 14\,\text{cm}$.
\end{enumerate}
\end{corrige}

\begin{corrige}[Exercice 3 — Reconnaissance de la configuration : QCM]
La bonne réponse est la \textbf{proposition 5} : $\dfrac{AM}{AB} = \dfrac{MN}{BC} = \dfrac{AN}{NC}$ n'est pas une conséquence de Thalès.

En effet, Thalès donne $\dfrac{AM}{AB} = \dfrac{AN}{AC} = \dfrac{MN}{BC}$. Or dans la proposition 5, le dernier rapport est $\dfrac{AN}{NC}$ et non $\dfrac{AN}{AC}$ : en général $\dfrac{AN}{AC} \neq \dfrac{AN}{NC}$.

Vérifions que les autres sont bien vraies :
\begin{itemize}
    \item (1) est l'énoncé direct de Thalès.
    \item (2) se déduit de (1) : si $\dfrac{AM}{AB} = \dfrac{AN}{AC} = k$, alors $MB = AB - AM = AB(1-k)$ et $NC = AC(1-k)$, donc $\dfrac{AM}{MB} = \dfrac{k\,AB}{(1-k)AB} = \dfrac{k}{1-k} = \dfrac{AN}{NC}$.
    \item (3) se déduit de (1) par inversion des rapports : $\dfrac{AB}{AM} = \dfrac{AC}{AN}$ donne, par produit en croix, $\dfrac{AB}{AC} = \dfrac{AM}{AN}$.
    \item (4) : $\dfrac{MB}{AB} = 1 - k = \dfrac{NC}{AC}$.
\end{itemize}
\end{corrige}

\begin{corrige}[Exercice 4 — Contraposée : Droites non parallèles]
Calculons séparément les deux rapports :
\[ \dfrac{RU}{RS} = \dfrac{6}{15} = \dfrac{2}{5} = 0,4 \qquad \text{et} \qquad \dfrac{RV}{RT} = \dfrac{4}{12} = \dfrac{1}{3} \approx 0,333. \]
On constate que $\dfrac{RU}{RS} \neq \dfrac{RV}{RT}$.

Les points $R, U, S$ et $R, V, T$ sont alignés dans le même ordre. D'après la \textbf{contraposée} de la propriété de Thalès, si les droites $(UV)$ et $(ST)$ étaient parallèles, on aurait $\dfrac{RU}{RS} = \dfrac{RV}{RT}$, ce qui est faux.

\textbf{Conclusion :} les droites $(UV)$ et $(ST)$ ne sont \textbf{pas parallèles}.
\end{corrige}

\begin{corrige}[Exercice 5 — Calculs de longueurs dans un triangle (Type BEPC)]
\begin{enumerate}
    \item \begin{popfigure}
\begin{tikzpicture}[scale=0.7,>=Stealth]
\coordinate (A) at (0,4);
\coordinate (B) at (-3,0);
\coordinate (C) at (5,0);
\coordinate (M) at ($(A)!0.4!(B)$); % AM/AB = 4/10
\coordinate (N) at ($(A)!0.4!(C)$); % AN/AC = 4/10
\draw[thick] (A) -- (B) -- (C) -- cycle;
\draw[thick, popblue] (M) -- (N);
\draw[dashed] (A) -- ($(A)!-0.5!(B)$) (A) -- ($(A)!-0.5!(C)$);
% Labels
\node[above] at (A) {$A$};
\node[left] at (B) {$B$};
\node[right] at (C) {$C$};
\node[left] at (M) {$M$};
\node[right] at (N) {$N$};
% Dimensions
\node[left] at ($(A)!0.5!(M)$) {$4$};
\node[left] at ($(M)!0.5!(B)$) {$6$};
\node[right] at ($(A)!0.5!(N)$) {$6$};
\node[right] at ($(N)!0.5!(C)$) {$9$};
\node[below] at ($(B)!0.5!(C)$) {$12$};
% Parallel marks
\draw[popblue] ($(M)!0.3!(N)$) ++(0.2,0.1) -- ++(0.3,0) ($(M)!0.3!(N)$) ++(0.2,-0.1) -- ++(0.3,0);
\draw[popblue] ($(B)!0.5!(C)$) ++(0.2,0.1) -- ++(0.3,0) ($(B)!0.5!(C)$) ++(0.2,-0.1) -- ++(0.3,0);
\end{tikzpicture}
\legende{Figure du triangle $ABC$ avec $(MN)\parallel(BC)$}
\end{popfigure}
    \item Les droites $(MN)$ et $(BC)$ sont parallèles, $M \in [AB]$, $N \in [AC]$. Par la propriété directe de Thalès :
    \[ \dfrac{AM}{AB} = \dfrac{AN}{AC} \iff \dfrac{4}{10} = \dfrac{AN}{15} \iff AN = \dfrac{4 \times 15}{10} = 6\,\text{cm}. \]
    \item Toujours par Thalès : $\dfrac{AM}{AB} = \dfrac{MN}{BC} \iff \dfrac{4}{10} = \dfrac{MN}{12} \iff MN = \dfrac{4 \times 12}{10} = \dfrac{48}{10} = 4,8\,\text{cm}$.
    \item $MB = AB - AM = 10 - 4 = 6\,\text{cm}$ et $NC = AC - AN = 15 - 6 = 9\,\text{cm}$.
\end{enumerate}
\textbf{Barème indicatif (4 pts) :} figure 0,5 pt ; citation de Thalès avec hypothèses 1 pt ; calcul de $AN$ 1 pt ; calcul de $MN$ 1 pt ; $MB$ et $NC$ 0,5 pt.

\textbf{Points de vigilance :} citer explicitement les hypothèses ($M \in [AB]$, $N \in [AC]$, $(MN) \parallel (BC)$) avant d'écrire les rapports ; ne pas confondre $AC$ et $NC$ dans les rapports.
\end{corrige}

\begin{corrige}[Exercice 6 — Problème concret : Hauteur d'un poteau électrique]
\begin{enumerate}
    \item \begin{popfigure}
\begin{tikzpicture}[scale=0.6,>=Stealth]
% Ground
\draw[thick, popmuted] (-1,0) -- (10,0);
% Pole
\draw[thick, popdark] (0,0) -- (0,5) node[midway, left] {$h$};
% Student
\draw[thick, poporange] (6,0) -- (6,1.4) node[midway, right] {$1,40\,\text{m}$};
% Sun rays (parallel)
\draw[dashed, popgold] (-1,6) -- (7,-1);
\draw[dashed, popgold] (0,5) -- (8,0);
% Shadows
\draw[<->, popblue] (0,0) -- (8,0) node[midway, below] {$8,40\,\text{m}$};
\draw[<->, popblue] (6,0) -- (8,0) node[midway, below] {$1,20\,\text{m}$};
% Right angle marks
\draw (0,0) rectangle (0.3,0.3);
\draw (6,0) rectangle (6.3,0.3);
% Labels
\node[below left] at (0,0) {$P$ (pied poteau)};
\node[below] at (6,0) {$E$ (élève)};
\node[below right] at (8,0) {$O$ (bout ombres)};
\node[above left] at (0,5) {$H$ (haut poteau)};
\node[above right] at (6,1.4) {$T$ (tête élève)};
\end{tikzpicture}
\legende{Modélisation : deux triangles rectangles semblables}
\end{popfigure}
    Le poteau $[PH]$ et l'élève $[ET]$ sont verticaux, donc perpendiculaires au sol horizontal : ils sont \textbf{parallèles}. Les rayons solaires $[HO]$ et $[TO]$ sont parallèles (le soleil est très loin). Les points $P, E, O$ sont alignés sur le sol. On est dans une configuration de Thalès (triangles $PHO$ et $TEO$ semblables).
    \item Notons $h$ la hauteur du poteau. Par la propriété directe de Thalès dans le triangle $PHO$ (ou égalité des rapports dans les triangles semblables) :
    \[ \dfrac{h}{1,40} = \dfrac{8,40}{1,20} = 7 \quad \Longrightarrow \quad h = 7 \times 1,40 = 9,80\,\text{m}. \]
    Le poteau mesure \textbf{9,80 m} (soit $980\,\text{cm}$, déjà exact au centimètre).
    \item À \emph{instant fixe} et sur sol horizontal, les rayons solaires gardent la même direction. Si l'élève se déplace sur le sol, la longueur de son ombre \textbf{ne change pas} (les triangles restent semblables avec le même rapport). En revanche, l'ombre varie au cours de la journée car l'inclinaison des rayons solaires change (le rapport $\frac{\text{hauteur}}{\text{ombre}}$ n'est plus le même).
\end{enumerate}
\textbf{Points de vigilance :} bien convertir et garder les mêmes unités ; justifier le parallélisme des deux verticales (toutes deux perpendiculaires au sol) avant d'appliquer Thalès ; ne pas confondre position de l'élève et heure de la journée.
\end{corrige}

\begin{corrige}[Exercice 7 — Réciproque de Thalès : Parallélisme ou pas ?]
\begin{enumerate}
    \item $\dfrac{AM}{AB} = \dfrac{3,5}{8,4} = \dfrac{35}{84} = \dfrac{5}{12}$ et $\dfrac{AN}{AC} = \dfrac{2,5}{6} = \dfrac{25}{60} = \dfrac{5}{12}$.
    Les rapports sont égaux et les points sont dans le même ordre : par la \textbf{réciproque de Thalès}, $(MN) \parallel (BC)$. \textbf{OUI.}
    \item $AB = AM + MB = 5 + 7 = 12$ et $AC = AN + NC = 4 + 5,6 = 9,6$.
    $\dfrac{AM}{AB} = \dfrac{5}{12}$ et $\dfrac{AN}{AC} = \dfrac{4}{9,6} = \dfrac{40}{96} = \dfrac{5}{12}$.
    Rapports égaux, points dans le même ordre : par la réciproque de Thalès, $(MN) \parallel (BC)$. \textbf{OUI.}
    \item $\dfrac{AM}{AB} = \dfrac{8}{18} = \dfrac{4}{9}$ et $\dfrac{AN}{AC} = \dfrac{6}{13,5} = \dfrac{60}{135} = \dfrac{4}{9}$.
    Rapports égaux : par la réciproque de Thalès, $(MN) \parallel (BC)$. \textbf{OUI.}
\end{enumerate}
\textbf{Point de vigilance :} dans la question 2, il faut d'abord calculer $AB$ et $AC$ (sommes) avant de former les rapports — erreur classique que d'écrire directement $\dfrac{AM}{MB}$.
\end{corrige}

\begin{corrige}[Exercice 8 — Configuration « Papillon »]
\begin{popfigure}
\begin{tikzpicture}[scale=0.7,>=Stealth]
\coordinate (A) at (0,0);
\coordinate (B) at (-4,-3);
\coordinate (C) at (3,-2);
\coordinate (M) at ($(A)!0.4!(B)$); % AM/AB = 6/15 = 0.4
\coordinate (N) at ($(A)!0.4!(C)$); % AN/AC = 4/10 = 0.4
\coordinate (P) at ($(A)!0.6!(B)$); % AP/AB = 9/15 = 0.6
\draw[thick] (B) -- (A) -- (C);
\draw[thick, popblue] (M) -- (N);
\draw[thick, poporange, dashed] (P) -- (N);
\draw[thick] (M) -- (C) (N) -- (B); % Secantes
% Labels
\node[above] at (A) {$A$};
\node[left] at (B) {$B$};
\node[right] at (C) {$C$};
\node[left] at (M) {$M$};
\node[right] at (N) {$N$};
\node[left] at (P) {$P$};
% Parallel marks
\draw[popblue] ($(M)!0.5!(N)$) ++(-0.15,0.1) -- ++(0.3,0) ($(M)!0.5!(N)$) ++(-0.15,-0.1) -- ++(0.3,0);
\draw[popblue] ($(B)!0.5!(C)$) ++(-0.15,0.1) -- ++(0.3,0) ($(B)!0.5!(C)$) ++(-0.15,-0.1) -- ++(0.3,0);
\end{tikzpicture}
\legende{Configuration papillon : $(MN)\parallel(BC)$, $(PN)$ non parallèle}
\end{popfigure}
Les droites $(BM)$ et $(CN)$ se coupent en $A$ et $(MN) \parallel (BC)$ : le théorème de Thalès s'applique dans la configuration papillon :
\[ \dfrac{AM}{AB} = \dfrac{AN}{AC} = \dfrac{MN}{BC}. \]
\begin{enumerate}
    \item $\dfrac{AM}{AB} = \dfrac{AN}{AC} \iff \dfrac{6}{15} = \dfrac{4}{AC} \iff AC = \dfrac{15 \times 4}{6} = 10\,\text{cm}$.
    \item $\dfrac{AM}{AB} = \dfrac{MN}{BC} \iff \dfrac{6}{15} = \dfrac{MN}{10} \iff MN = \dfrac{6 \times 10}{15} = 4\,\text{cm}$.
    \item $\dfrac{AM}{AB} = \dfrac{6}{15} = \dfrac{2}{5} = 0,4$ et $\dfrac{AN}{AC} = \dfrac{4}{10} = \dfrac{2}{5} = 0,4$. Les rapports sont égaux : les données sont cohérentes avec le parallélisme de $(MN)$ et $(BC)$.
    \item On a $AP = AM + MP = 6 + 3 = 9\,\text{cm}$. Alors $\dfrac{AP}{AB} = \dfrac{9}{15} = \dfrac{3}{5} = 0,6$, tandis que $\dfrac{AN}{AC} = \dfrac{4}{10} = 0,4$. Comme $\dfrac{AP}{AB} \neq \dfrac{AN}{AC}$, par la \textbf{contraposée} de Thalès, la droite $(PN)$ n'est \textbf{pas parallèle} à $(BC)$.
\end{enumerate}
\textbf{Point de vigilance :} dans la configuration papillon, $A$ est le point d'intersection des sécantes ; les rapports s'écrivent toujours à partir de $A$.
\end{corrige}

\begin{corrige}[Exercice 9 — Partage d'un terrain triangulaire — Modélisation]
\begin{enumerate}
    \item Comme $(MN) \parallel (BC)$, le triangle $AMN$ est une réduction du triangle $ABC$ de rapport $k = \dfrac{AM}{AB}$. Les aires sont multipliées par $k^2$ :
    \[ \dfrac{\mathcal{A}_{AMN}}{\mathcal{A}_{ABC}} = k^2. \]
    \item On veut $k^2 = \dfrac{4}{9}$, donc $k = \sqrt{\dfrac{4}{9}} = \dfrac{2}{3}$ (on garde la valeur positive, $k$ étant une longueur relative).
    \item $AM = k \times AB = \dfrac{2}{3} \times 60 = 40\,\text{m}$. Le piquet $M$ se place à \textbf{40 m} de $A$ sur $[AB]$.
    \item Par Thalès : $MN = k \times BC = \dfrac{2}{3} \times 45 = 30\,\text{m}$. La clôture mesure \textbf{30 m}.
\end{enumerate}
\textbf{Barème indicatif (5 pts) :} rapport des aires $k^2$ 1,5 pt ; valeur de $k$ 1 pt ; $AM$ 1 pt ; $MN$ 1 pt ; rédaction et unités 0,5 pt.

\textbf{Points de vigilance :} le rapport des aires est le \emph{carré} du rapport des longueurs (ne pas écrire $k = \frac{4}{9}$ directement) ; prendre la racine \emph{positive} ; vérifier la cohérence : $\frac{2}{3}$ des longueurs donnent bien $\frac{4}{9}$ de l'aire.
\end{corrige}

\begin{corrige}[Exercice 10 — Mesure de la largeur du fleuve Wouri]
\begin{enumerate}
    \item \begin{popfigure}
\begin{tikzpicture}[scale=0.5,>=Stealth]
\coordinate (A) at (0,0);
\coordinate (C) at (3,0);
\coordinate (B) at (7.5,0);
\coordinate (D) at (3,5);
\coordinate (E) at ($(A)!0.4!(D)$); % AE/AD = 20/50 = 0.4
\draw[thick] (A) -- (B);
\draw[thick] (C) -- (D);
\draw[thick, popblue] (C) -- (E);
\draw[thick, popblue] (D) -- (B);
% Right angle
\draw (C) rectangle ++(0.3,0.3);
% Parallel marks
\draw[popblue] ($(C)!0.5!(E)$) ++(0.1,0.15) -- ++(0.2,-0.1) ($(C)!0.5!(E)$) ++(0.1,-0.15) -- ++(0.2,0.1);
\draw[popblue] ($(D)!0.5!(B)$) ++(0.1,0.15) -- ++(0.2,-0.1) ($(D)!0.5!(B)$) ++(0.1,-0.15) -- ++(0.2,0.1);
% Labels
\node[below left] at (A) {$A$};
\node[below] at (C) {$C$};
\node[below right] at (B) {$B$};
\node[above] at (D) {$D$};
\node[left] at (E) {$E$};
% Dimensions
\node[below] at ($(A)!0.5!(C)$) {$30\,\text{m}$};
\node[below] at ($(C)!0.5!(B)$) {$45\,\text{m}$};
\node[left] at ($(A)!0.5!(E)$) {$20\,\text{m}$};
\node[left] at ($(E)!0.5!(D)$) {$30\,\text{m}$};
\node[right] at ($(C)!0.5!(D)$) {$40\,\text{m}$};
\end{tikzpicture}
\legende{Dispositif de mesure de la largeur $AB$}
\end{popfigure}
    \item Les points $A, C, B$ sont alignés (sur la rive de départ), les points $A, E, D$ sont alignés (car $E \in [AD]$), et $(CE) \parallel (DB)$. On est donc dans la \textbf{configuration de Thalès} appliquée au triangle $ADB$ : les triangles $ACE$ et $ADB$ ont leurs côtés proportionnels, ils sont semblables.
    \item Par la propriété directe de Thalès dans le triangle $ADB$ :
    \[ \dfrac{AC}{AB} = \dfrac{AE}{AD} \iff \dfrac{30}{AB} = \dfrac{20}{50} \iff AB = \dfrac{30 \times 50}{20} = \dfrac{1500}{20} = 75\,\text{m}. \]
    La largeur du fleuve est $\mathbf{AB = 75\,\text{m}}$.
    \item Toujours par Thalès : $\dfrac{AE}{AD} = \dfrac{CE}{DB} \iff \dfrac{20}{50} = \dfrac{CE}{40} \iff CE = \dfrac{20 \times 40}{50} = \dfrac{800}{50} = 16\,\text{m}$.
\end{enumerate}
\textbf{Barème indicatif (6 pts) :} figure codée 1 pt ; justification de la configuration (alignements + parallélisme) 1,5 pt ; égalité des rapports 1 pt ; calcul de $AB$ 1,5 pt ; calcul de $CE$ 1 pt.

\textbf{Points de vigilance :} citer les deux alignements de points \emph{et} le parallélisme avant d'écrire Thalès ; bien identifier le triangle « contenant » ($ADB$) et le triangle « réduit » ($ACE$).
\end{corrige}

\begin{corrige}[Exercice 11 — Synthèse : Figure complexe et raisonnement par l'absurde]
\begin{popfigure}
\begin{tikzpicture}[scale=0.6,>=Stealth]
\coordinate (A) at (0,4);
\coordinate (B) at (-4,0);
\coordinate (C) at (6,0);
\coordinate (M) at ($(A)!0.4!(B)$); % AM/AB = 2/5
\coordinate (N) at ($(A)!0.375!(C)$); % AN/AC = 3/8
\coordinate (Np) at ($(A)!0.4!(C)$); % AN'/AC = 2/5
\coordinate (P) at ($(C)!0.625!(B)$); % CP/CB = 5/8
\draw[thick] (A) -- (B) -- (C) -- cycle;
\draw[thick, popblue] (M) -- (Np);
\draw[thick, poporange, dashed] (M) -- (N);
\draw[thick, popgreen] (N) -- (P);
% Parallel marks
\draw[popblue] ($(M)!0.5!(Np)$) ++(-0.1,0.15) -- ++(0.2,-0.1) ($(M)!0.5!(Np)$) ++(-0.1,-0.15) -- ++(0.2,0.1);
\draw[popblue] ($(B)!0.5!(C)$) ++(-0.1,0.15) -- ++(0.2,-0.1) ($(B)!0.5!(C)$) ++(-0.1,-0.15) -- ++(0.2,0.1);
\draw[popgreen] ($(N)!0.5!(P)$) ++(0.1,0.1) -- ++(-0.15,0) ($(N)!0.5!(P)$) ++(0.1,-0.1) -- ++(-0.15,0);
\draw[popgreen] ($(A)!0.5!(B)$) ++(0.1,0.1) -- ++(-0.15,0) ($(A)!0.5!(B)$) ++(0.1,-0.1) -- ++(-0.15,0);
% Labels
\node[above] at (A) {$A$};
\node[left] at (B) {$B$};
\node[right] at (C) {$C$};
\node[left] at (M) {$M$};
\node[right] at (N) {$N$};
\node[right] at (Np) {$N'$};
\node[below] at (P) {$P$};
% Dimensions
\node[left] at ($(A)!0.5!(M)$) {$2$};
\node[left] at ($(M)!0.5!(B)$) {$3$};
\node[right] at ($(A)!0.5!(N)$) {$3$};
\node[right] at ($(N)!0.5!(Np)$) {$1$};
\node[right] at ($(Np)!0.5!(C)$) {$4$};
\node[below] at ($(B)!0.5!(C)$) {$20$};
\end{tikzpicture}
\legende{Figure de synthèse : $(MN')\parallel(BC)$, $(NP)\parallel(AB)$}
\end{popfigure}
\begin{enumerate}
    \item Comparons : $\dfrac{AM}{AB} = \dfrac{2}{5} = \dfrac{16}{40}$ et $\dfrac{AN}{AC} = \dfrac{3}{8} = \dfrac{15}{40}$. Donc $\dfrac{AM}{AB} \neq \dfrac{AN}{AC}$.
    Les points $A, M, B$ et $A, N, C$ sont alignés dans le même ordre. Par la \textbf{contraposée} de la propriété de Thalès, si $(MN) \parallel (BC)$ on aurait égalité des rapports, ce qui est faux. Donc $(MN)$ et $(BC)$ ne sont \textbf{pas parallèles}.
    \item Puisque $(MN') \parallel (BC)$ avec $M \in [AB]$ et $N' \in [AC]$, la propriété directe de Thalès donne :
    \[ \dfrac{AN'}{AC} = \dfrac{AM}{AB} = \dfrac{2}{5}. \]
    \item $\dfrac{AN}{AC} = \dfrac{3}{8} = 0,375$ et $\dfrac{AN'}{AC} = \dfrac{2}{5} = 0,4$. Donc $AN < AN'$ : le point $N$ est situé \textbf{entre $A$ et $N'$} sur le segment $[AC]$ ($N$ est plus proche de $A$ que $N'$).
    \item La parallèle à $(AB)$ passant par $N$ coupe $(BC)$ en $P$. On applique Thalès dans le triangle $CAB$ avec $N \in [CA]$, $P \in [CB]$ et $(NP) \parallel (AB)$ :
    \begin{itemize}
        \item $\dfrac{CN}{CA} = 1 - \dfrac{AN}{AC} = 1 - \dfrac{3}{8} = \dfrac{5}{8}$.
        \item Par Thalès : $\dfrac{CP}{CB} = \dfrac{CN}{CA} = \dfrac{5}{8}$.
        \item $CP = \dfrac{5}{8} \times BC = \dfrac{5}{8} \times 20 = \dfrac{100}{8} = 12,5\,\text{cm}$.
    \end{itemize}
\end{enumerate}
\textbf{Barème indicatif (6 pts) :} comparaison des rapports et contraposée 1,5 pt ; $\frac{AN'}{AC}$ 1 pt ; comparaison et position relative 1 pt ; $\frac{CN}{CA}$ 1 pt ; $\frac{CP}{CB}$ et $CP$ 1,5 pt.

\textbf{Points de vigilance :} pour la contraposée, rédiger proprement le raisonnement par l'absurde (« supposons les droites parallèles, alors... contradiction ») ; dans la question 4, le sommet de référence de Thalès est $C$ (et non $A$) car le parallélisme est avec $(AB)$ : les rapports s'écrivent $\frac{CN}{CA} = \frac{CP}{CB}$.
\end{corrige}

\newpage

\coursec{Fiche bilan}

\begin{popfigure}
\begin{center}
\begin{tikzpicture}[mindmap, grow cyclic, every node/.style={concept, font=\small},
root concept/.append style={concept color=popblue, font=\bfseries},
level 1/.append style={level distance=3.4cm, sibling angle=51.4},
level 2/.append style={level distance=2.2cm, font=\scriptsize}]
\node[root concept]{Thalès dans le triangle}
  child[concept color=poporange]{ node {Configuration}
    child { node {Deux droites sécantes coupées par deux parallèles} }
  }
  child[concept color=popgreen]{ node {Propriété directe}
    child { node {$\dfrac{AM}{AB}=\dfrac{AN}{AC}=\dfrac{MN}{BC}$} }
  }
  child[concept color=poppurple]{ node {Réciproque}
    child { node {Égalité des rapports $+$ alignement dans le même ordre} }
  }
  child[concept color=poppink]{ node {Contraposée}
    child { node {Rapports différents $\Rightarrow$ non parallèles} }
  }
  child[concept color=popgold]{ node {Applications}
    child { node {Hauteurs, distances, agrandissements} }
  }
  child[concept color=popblue!70]{ node {Méthodes}
    child { node {Identifier, écrire, calculer, conclure} }
  }
  child[concept color=popgreen!70]{ node {Vie courante}
    child { node {Ombres, plans, topographie} }
  };
\end{tikzpicture}
\end{center}
\legende{Carte mentale de la leçon : Thalès dans le triangle}
\end{popfigure}

\begin{bilan}
\textbf{1. Configuration de Thalès.}
\begin{itemize}
\item $ABC$ est un triangle ; $M$ est un point de la droite $(AB)$ et $N$ un point de la droite $(AC)$, avec $(MN) \parallel (BC)$.
\item On distingue deux cas : la \cle{configuration « papillon »} (droites sécantes, parallèles de part et d'autre du sommet) et la configuration « emboîtée » ($M$ et $N$ sur les côtés du triangle).
\end{itemize}

\textbf{2. Les deux théorèmes à connaître par cœur.}
\begin{center}
\begin{tabularx}{\linewidth}{|l|X|X|}
\hline
\rowcolor{popblueL} \textbf{Théorème} & \textbf{Hypothèses} & \textbf{Conclusion} \\
\hline
Propriété directe de Thalès & $M \in (AB)$, $N \in (AC)$ et $(MN) \parallel (BC)$ & $\dfrac{AM}{AB} = \dfrac{AN}{AC} = \dfrac{MN}{BC}$ \\
\hline
Réciproque de Thalès & $A, M, B$ et $A, N, C$ alignés \textbf{dans le même ordre} et $\dfrac{AM}{AB} = \dfrac{AN}{AC}$ & $(MN) \parallel (BC)$ \\
\hline
Contraposée & $\dfrac{AM}{AB} \neq \dfrac{AN}{AC}$ & $(MN)$ et $(BC)$ \textbf{ne sont pas} parallèles \\
\hline
\end{tabularx}
\end{center}

\textbf{3. Méthode express pour calculer une longueur.}
\begin{enumerate}
\item Repérer la configuration de Thalès (droites parallèles, points alignés).
\item Citer les hypothèses et le théorème dans la rédaction.
\item Écrire l'égalité des trois rapports, puis garder les deux utiles.
\item Remplacer par les valeurs connues et utiliser le \cle{produit en croix}.
\item Conclure avec l'unité et une phrase.
\end{enumerate}

\textbf{4. Méthode express pour prouver un parallélisme.}
\begin{enumerate}
\item Calculer \emph{séparément} les rapports $\dfrac{AM}{AB}$ et $\dfrac{AN}{AC}$.
\item Vérifier que les points sont alignés dans le même ordre.
\item Si les rapports sont égaux : réciproque de Thalès $\Rightarrow$ parallèles. S'ils sont différents : contraposée $\Rightarrow$ non parallèles.
\end{enumerate}

\textbf{5. Piège classique.} Ne jamais appliquer Thalès sans avoir vérifié (et écrit) l'hypothèse de parallélisme pour la propriété directe, ni l'ordre des points pour la réciproque.
\end{bilan}

\begin{aretenir}[Checklist avant le BEPC]
Avant l'épreuve, vérifie que tu sais :
\begin{itemize}
\item[$\square$] Reconnaître une configuration de Thalès (emboîtée ou papillon) dans une figure.
\item[$\square$] Énoncer la propriété directe de Thalès avec toutes ses hypothèses.
\item[$\square$] Écrire correctement l'égalité $\dfrac{AM}{AB}=\dfrac{AN}{AC}=\dfrac{MN}{BC}$.
\item[$\square$] Calculer une longueur inconnue par le produit en croix.
\item[$\square$] Rédiger la solution en citant le théorème utilisé.
\item[$\square$] Énoncer la réciproque de Thalès et vérifier l'ordre des points alignés.
\item[$\square$] Prouver que deux droites sont parallèles en comparant deux rapports.
\item[$\square$] Prouver que deux droites ne sont \emph{pas} parallèles grâce à la contraposée.
\item[$\square$] Distinguer clairement propriété directe, réciproque et contraposée.
\item[$\square$] Résoudre un problème concret (hauteur d'un arbre, largeur d'une rivière, plan à l'échelle) avec Thalès.
\item[$\square$] Donner le résultat avec l'unité et, si besoin, un arrondi justifié.
\item[$\square$] Vérifier la cohérence du résultat (une longueur cherchée plus petite que la longueur totale, etc.).
\end{itemize}
\end{aretenir}

\findecours

\end{document}
